\part{Nora Adeyemi}

\numbered{1}{Two in the Morning}

\textsc{the call froze at} seven minutes past two, as it had the night before, and Warren Brandt, who had been telling his daughter about a bird he had seen from the car, was left looking at a small paused face with its mouth open on a vowel.

Pia was seven. She was in a kitchen eleven time zones away where it was, by the clock on the wall behind her, the middle of a bright afternoon, and the frozen image had her one hand lifted toward the screen as though she meant to take hold of the bird. Warren waited the forty seconds that he had learned the call would require. The picture did not recover. A grey wheel turned in the corner of the window, the word \emph{reconnecting} appeared beneath it, and then the call dropped entirely and the room around him became very quiet, in the way of hotel rooms at that hour, which are quiet by construction and are never silent.

The room smelled of new carpet and of the faint sweetness of a cleaning product that was trying to be a flower. Under those, if he breathed slowly, there was the smell of every previous guest: a warmth, a dampness, the ghost of someone's hair gel, a trace of the oranges that a traveller had peeled on the bed and thought no one would notice. The air conditioning made a sound he had been ignoring for two nights and now, for the first time, could not. It was a low roar with a ticking inside it, like a distant sea that kept a clock. Somewhere beyond the wall an ice machine dropped a load of cubes into its bin with a noise like a handful of teeth. Through the thin curtain, the street below gave up a single car, its tyres hissing on the wet, and then nothing, and the nothing had a pitch.

He did not call back at once. This was the part of his behaviour that, in later years, the people who studied him would find hardest to render in the language they had available, because it was not patience and it was not resignation. He sat on the edge of a bed that had been made with unusual tightness and he looked at the time, and then he took a sheet of hotel stationery from the desk, uncapped the hotel's pen, and wrote \emph{2:07} at the top of the page, with the date, and under it the word \emph{Tues}, because the night before had been a Tuesday and he saw no reason to leave the comparison to memory.

It would be easy to say that Warren was a man who solved problems, and the sentence would be accurate and would mislead. Plenty of people solve problems. What distinguished him was a habit of the attention that had been present since early childhood and had never been trained out of him, perhaps because nobody had ever been able to say what it was. Offered a difficulty, he went at once to the smallest true thing about it. Other people, offered the same difficulty, went to the most important thing, or the most urgent thing, or the thing that would best be said aloud in a meeting, and these were frequently not true, or were true only at a scale that made them useless. Warren went down. If a call failed at the same minute on two successive nights, then the minute was a fact, and the fact had a cause, and the cause was in the building, and the building was something a person could walk through.

He put on his shoes. They were the black leather ones he had bought for a funeral two years before and worn on every trip since, on the principle that a man who owns one good pair should let as many ceilings as possible see him in it. He did not put on a jacket, which meant that the key card for the elevator, being in the jacket, stayed on the chair, and so he left the room for the corridor of the ninth floor with the heavy door sighing shut behind him and no means of opening it again. He noticed this about four steps along and decided that it was a problem for a later and more resourceful version of himself.

The corridor was long and carpeted in a pattern of interlocking leaves that someone had chosen, he supposed, to hide the evidence of travel. The carpet had the faint give of an underlay laid generously and never replaced, and his shoes made on it no sound at all, so that he moved along the corridor like a man in a dream of being polite. Brass lamps in the shape of closed tulips lit the walls at intervals, and between the lamps the doors went off in pairs, each with its small illuminated number and its small hanging card asking, in three languages, not to be disturbed. At the far end a window gave onto an air shaft, and beside the window, set into the ceiling panel at the point where the corridor turned, was a white plastic disc about the size of a saucer with a single amber light on its rim. The light blinked. It blinked at a rate that Warren, standing beneath it with his head tipped back, found he could count, and he counted it, and he found that there was a second, fainter pattern inside the first, a flutter every ninth beat that suggested the device was doing something other than idling.

He had a theory, which was that the router served the whole floor, and that something on the floor, some scheduled task of the building's, woke at two and took the connection for several minutes. He had another theory, which was that the router served the building, in which case the task was larger and the building's habits ought to show elsewhere, in the elevators, in the lights, in the door at the bottom of the stairs. He did not hold either theory firmly. What he held firmly was the conviction that one of them could be tested before breakfast by anyone willing to ask a person who knew how the hotel was wired, and that this person was almost certainly awake.

A housekeeping cart stood against the wall near the lift, abandoned, its top shelf lined with small bottles and folded towels and a single paperback with a cracked spine. Warren looked at the paperback for a while, and then back up at the disc, and then, because there was no one to hear it, said aloud to the corridor, ``One per floor, or one for the whole building,'' and he was pleased with the sentence, as one is pleased with a tool that has fitted a bolt.

It was a good question. This should be said now, because the book that follows will spend several hundred pages on what became of it. It was a clear question, it had a definite answer, and any competent night engineer would have given the answer in a sentence and a half. It had been formed in the right way, out of a fact, in response to a failure, by a person who wished to know. Nothing was wrong with it. The difficulty, which nobody on the ninth floor of that hotel could have foreseen, was entirely a matter of who would later be standing within earshot when he asked it again.

He went back down the corridor to his door and stood in front of it and remembered about the key card. After some thought he knocked, softly, on the door of the room next to his, behind which a deputy from the company's legal team was asleep, and then thought better of it and stopped. He sat on the carpet with his back against his own door and his knees drawn up, and he took the stationery from his pocket, where he had put it without noticing, and under \emph{2:07} he wrote a second line, \emph{router, 9th fl, amber, 9-beat flutter}, and under that, after a moment, \emph{ask}.

Eight floors below, in a conference room that had not yet been booked for the morning, a minister of one of the host country's ministries was rehearsing, to a printout of her own notes, a single sentence about accountability that she had been refining for three weeks and intended to deliver at a quarter past nine. She had it nearly right. Neither of the two questions that would be put in that room at a quarter past nine was a bad one, and each would be heard, by the people who heard it, as something other than what had been said.

Warren Brandt waited in the corridor for a housekeeper to come and let him in, and, because it was the only thing he could usefully do, fell asleep sitting up.

\numbered{2}{The Interview}

\textsc{nora adeyemi had spent} eleven years of her working life being the person in the room who knew what the other side wanted before they had finished saying it, and she arrived at Harbor's headquarters for the final interview with the settled calm of someone who has been told the question in advance. She had not been told the question. She had simply assumed, on the strength of long experience, that any company which sent a car for a candidate and then kept her waiting nineteen minutes in a glass atrium was running a test, and that the test was the waiting. She had decided in the car to pass it by looking entirely unbothered, and had spent the first six minutes of the wait perfecting an expression of calm that she now suspected was visible from the mezzanine.

It was a good atrium. A tree grew in the middle of it, an actual tree, a Japanese maple lit from below, and around the tree a ring of low chairs had been arranged at a distance from one another that suggested the designer had thought carefully about how strangers feel when seated near strangers. Nora sat in the chair that faced the lifts and watched people come and go. She noticed that nobody looked up at the tree. She noticed that the stairwell, off to the left behind a frosted door, was used by perhaps one person in ten, and that the ones who used it carried themselves a little differently from the others, with a faint air of having made a decision.

The atrium had a smell, too, which Nora identified after a minute as the smell of money that had been told to be calm: cold stone, expensive coffee, and the faintly resinous wet of the mulch around the maple's roots, where an invisible system had been dripping water at intervals since six that morning. The lifts announced themselves with a note, one soft chime, a perfect fourth, and she found herself counting the chimes and then, to her own irritation, trying to guess from the interval which of the four lifts had arrived. Heels crossed the floor. Trainers crossed the floor. A man in a fleece carrying a bowl of soup in both hands moved through the space as though crossing a river on stones. Nobody spoke above a murmur. The building seemed to have agreed that raised voices were a technical fault.

She had been a trade negotiator for the better part of her career, first for a ministry and later for an intergovernmental body whose name no one outside the field could ever reproduce, and the work had taught her that most of what is called persuasion is the patient management of what the other party already believes. She was good at it. She was known to be good at it, in the small and unforgiving world where such things are known, and it was because of this reputation that Harbor had approached her, in the person of its chief operating officer, Diane Kessler, who had called her at home on a Sunday and said, with no preamble, ``We are about to enter forty countries and we have no one who has ever had to give anything up.''

The interview took place on the eleventh floor, in a room with a long table and a single window. Diane was there, and beside Diane was a man in a grey sweater who had been introduced as Warren and was holding, as far as Nora could tell, a paper cup of water that he had not drunk from. He was younger than his photographs, or perhaps only less arranged. He shook her hand and held it a moment too long, in the manner of a person who has noticed something about the hand and is deciding whether to mention it.

``How did you get here today?'' he said.

It was not a question she had prepared for, and for a fraction of a second she felt the particular exhilaration of the trained negotiator who realises the real talks have begun. A man with a company of that size did not ask a candidate about her route. He asked it to see what she would do with an opening. She considered the possibilities. A direct answer would be a small thing, honest and forgettable. An elaborate answer would show that she knew it was a test. She decided that the correct response was to treat the question as an instrument for something larger and to answer both layers at once.

``By train, and then on foot,'' she said. ``Forty minutes, door to door. I prefer it. I do some of my best thinking in transit, and I like to arrive having already had the argument with myself.''

``Do you take the stairs?''

``When I can.'' She let a beat go by. ``In my experience, a person who takes the stairs has made a small decision that nobody asked her to make, and the habit carries over. I'd rather work with people who make those decisions. They tend to be the ones who read the clause before they sign.''

Diane smiled, very slightly, at the table. Warren nodded, slowly, as though she had said something that agreed with a private view of his own.

What Warren Brandt wanted to know was how long it took to get from the street to the eleventh floor by lift, because for two weeks he had been timing the lobby, and the numbers did not make sense to him. Visitors reached the atrium in about ninety seconds from the street door and then waited, on average, four minutes for a lift, and the building's operations team had reported satisfaction with this on the grounds that the lifts were within specification. He suspected the specification of being wrong, and he had been asking people how they arrived in order to find out whether the wait was something visitors noticed or merely something they bore. Nora's reply had told him nothing at all about the lifts. It had told him that she walked, that she enjoyed arguing with herself, and that she knew what a clause was. He liked her immediately, and for none of the reasons she supposed.

They talked for an hour and ten minutes, mostly about trade. She spoke at length about a customs arrangement she had helped negotiate between two states that disliked each other, and the account contained a number of small reversals that she delivered well. Warren asked three questions, each of them narrow, each of them about a mechanism: who had signed the original draft, how long the translation took, what happened if the officials on the border were replaced. She found them unusually good, the questions of a man who understood how agreements decayed, and she answered with more candour than she had planned. At the end Diane said, ``We would like you to start in the autumn,'' and Nora said she would need to see the terms, and Warren said, ``Of course,'' and then, as they stood, ``Is it all right if I walk down with you? I'd like to try the stairs.''

They went down eleven floors together, not speaking, in a stairwell painted a hospital green. Somewhere around the sixth landing he took out his phone and noted a number. At the bottom he thanked her and went back up in the lift, and she stood in the atrium beside the maple and thought, with a conviction she did not examine, that she had been chosen for her judgment, and that this was exactly what she had wanted to be true.

\interchapter{The Hall Monitor}

\textsc{when warren was ten,} the school he attended in a mild coastal town had a system of hall monitors that consisted of a laminated list, a corkboard, and eleven students who were trusted to wear a yellow band on the left arm and to say ``walking feet'' to anyone who ran. The system was not good. It was not meant to be good. Its purpose was to give eleven children the experience of being trusted with something that did not matter, and it had worked well enough for two generations.

Warren read the list on the corkboard in September and noticed that the same four children appeared on it every week, that two of them were in the same class, and that on Fridays, when the school emptied by a different door, nobody was assigned to it at all. A mitten had gone missing from that door in October, and the matter had been discussed in assembly, and it was this that first drew his attention to the list. By the middle of October he had replaced the list with a rotating schedule of his own devising, colour coded in four shades, with a second sheet for substitutes and a third that dealt with the contingency of an illness. He worked on it for six evenings. It was beautiful. Every child in the school appeared on it exactly once per term, and the Friday door had a monitor, and the alternate schedule folded into the first like a map into a pocket.

The teachers admired it openly. The headmistress put it in the newsletter. A visiting inspector, who came in November, asked to be shown the system and was shown it and said it was the most complete he had seen in a school of that size, and Warren, standing at the corkboard with his hands at his sides, felt a pride so clean that it hurt a little.

In December the vice principal, a tired man called Mr.~Lowry who had been in charge of the original list for nineteen years, found himself standing in front of the board one morning with a mug of tea, reading it, and he asked, to nobody, what it was for. The word was \emph{for}. He did not mean it unkindly. It simply struck him, with the force of a thing noticed late, that the school had been entirely reorganised around a system whose purpose had never been stated, and that he could not now recall what problem it had been introduced to solve. He asked Warren, who was passing. Warren thought for a long moment, honestly, and said that he did not know, and that he had assumed somebody did. Neither of them was troubled. In January, a first-grader named Dennis Pike walked out of the Friday door in his socks, and was stopped at the gate by an eleven-year-old in a yellow band who had been assigned to that door for the first time in the school's history. The school took this as proof that the system worked, which it had, and nobody went back to the question.

\numbered{4}{The Minister}

\textsc{the delegation travelled in} two cars, and Nora, who had expected to be put in the second, found herself placed beside Warren in the first, and understood at once that this was the form the reward would take. Nobody had said anything. A junior aide had opened the door, and she had got in, and in the course of the eleven-minute drive from the hotel to the ministry she saw, with the clarity of a person who has just crossed a line she did not know was there, that she was now somebody whose seat was decided in advance.

The city outside the car windows had the flat, colourless light of a morning that had not yet decided to be a day. There was a smell of diesel and baked bread where the car crossed the market, and for a moment, with the window down, a gust of something sweeter, grilled corn or burning sugar, from a stall that was setting up on a corner. Boys in school uniforms ran alongside the traffic selling packets of tissues. A shutter went up with a long iron shriek. On the pavement outside the ministry a woman in a blue wrap was sweeping the steps with a broom of bound twigs, in short fierce strokes, as though the dust were an argument she meant to win, and when the car stopped she stepped aside without looking up and went on sweeping behind it.

Inside, the building smelled of floor wax and cold radiators and the warm paper smell of files kept too long in rooms with the windows shut. The stairs were marble, worn into shallow dishes at the centre of every tread by a century of ministers' shoes. Somewhere above them a typewriter was being used, actually used, and the sound of it came down the stairwell with a dry, regular patter, like rain on a tin roof in a country where it seldom rained.

She had been with Harbor nine days. In that time she had read four hundred pages of internal documents on the company's approach to the country they were about to enter, and she had formed a view. The view was that the country was a hard case in the specific sense that its government had been burned before: a foreign platform had arrived, eleven years earlier, with promises, and had left behind a regulatory regime that treated every foreign network as a hostile act in waiting. The minister they were to meet, Ines Lindqvist, was known to be intelligent, exacting and indifferent to charm. The briefing papers called her \emph{the most difficult interlocutor in the region}, and Nora, who had negotiated with more difficult ones, had not been alarmed.

The meeting took place in a room on the fourth floor of a building that had been built to be impressive in an earlier century and had since been fitted with the usual apparatus of the present one: cables taped along skirting boards, a conference phone on a rolling cart, a bank of screens at the far end where two technicians sat with headsets. Lindqvist received them standing. She was a small woman with grey in her hair and a way of looking at each person as though recording a measurement. Her handshake was dry and brief, and withdrew at precisely the moment Nora had begun to enjoy it. She spoke in her own language, and a translator, a young man with a thin moustache, sat at her elbow and relayed her words through a headset system that sent them to a small receiver on each visitor's chair.

The first twenty minutes were what such meetings always are. Lindqvist's staff described the country's concerns, which were three, and which concerned data, jurisdiction and the question of who in the country would answer for what in the platform. One of her aides, a lean man named Faruk, spent some time on the matter of infrastructure. The country, he explained, connected to the wider world through a single international gateway, a legacy of an earlier regime, and the ministry was anxious about any foreign system that would make itself dependent on that one point. A fault at the gateway could cut a population of thirty million from the world. A foreign company that arranged its service around the gateway made itself a participant in that risk. He said all this carefully, with diagrams, and Nora took notes and glanced at Warren, who was looking at the diagrams with an expression she took, correctly, to be one of deep attention.

Then Lindqvist spoke. She did not raise her voice, and she had no notes, and the translator took a half second longer than usual to find the English.

``When your system decides what a country hears, and it will, because that is what a system of this kind does,'' she said, ``I would like to know whose name is on the decision.''

The room was still. It was the kind of sentence that a delegation spends a week preparing for and cannot answer in the time available, because the true answer was that there was no name, and everyone on Harbor's side of the table knew it. Diane Kessler, at the far end, began to open her mouth. Nora had her own answer ready, a good one, sound, built from the language of oversight boards and escalation paths, and she was drawing a breath to deliver it when something happened to the translation system.

It dropped. The headsets clicked, the small receivers on the chairs went silent, and for about forty seconds the room was entirely without translation. The translator, startled, began speaking directly into the air in a slightly too loud English, but his sentences were scattered, and the technicians at the far end bent over their screens. Warren, who had been holding his paper cup, put it down and turned in his chair.

``Excuse me,'' he said. ``Is this one router per floor, or one for the whole building?''

He said it to the room, pleasantly, in the tone of a man who has noticed a thing and would like to resolve it. Nobody answered. The translator looked at the minister, and Faruk looked at the ceiling, and the two technicians at the screens, who had heard the question perfectly, exchanged a glance that Nora noticed and did not read.

It was the silence that mattered. Forty seconds earlier a ministry aide had spent ten minutes explaining that the country's whole dependence on a single gateway was the central fear of the state, and now a foreign executive had asked, in the middle of the minister's most important remark, whether the connection ran through one point or many. The question fitted too well. A room of intelligent people, primed to listen for a particular concern, heard that concern in the first plausible sentence offered, and heard it as a rebuke: \emph{you speak of accountability, and I speak of the structure that makes accountability possible.} Faruk's pen stopped moving. Lindqvist, who had been about to say something else, closed her mouth and regarded Warren with an expression that was neither warm nor cold but had acquired, quite suddenly, an edge of interest.

Nora felt the moment come to her like a door opening in a corridor. She understood, in the space of a heartbeat, what they were all hearing, and she understood too that Warren, who looked pleasantly baffled, was not going to supply it. Somebody had to. It was a negotiator's reflex, the oldest she had: when the other side offers you an ambiguous remark, you do not ask what it means, you tell them what it means, in a form that serves you.

``I think what Warren is pointing to,'' she said, and her voice came out steadier than she felt, ``is that accountability begins with architecture. Who answers for a decision depends on where the decision is made, and where it is made depends on how the system is built. If the country's connection runs through one point, then that point is where responsibility concentrates. We would want to design around it. We would want to make sure no single point, ours or yours, becomes the place where a name cannot be found.''

She had not planned any of it. It came out whole, and when she finished, nobody spoke. Then Lindqvist inclined her head by about a centimetre, and Faruk wrote something down, and across the table one of the young staff members, who had not yet written a word, began to take notes in a fast neat hand.

``That is a serious point,'' said the minister, in English, for the first time. ``I accept it, for the moment. And whose name?''

It was said without heat, and it landed harder than anything before it. The room had begun to relax, and Nora felt it tighten again around her. She said that Harbor would put its answer in writing within the week, and Lindqvist said that she would be glad to read it, in a voice that made it plain she would also note how long it took to arrive. The meeting moved on to the matter of where the data would be kept. The question stayed where Nora had left it, in the middle of the table, and when the delegation rose to go only the minister and the young man with the fast neat hand looked at it again.

In the car afterward Diane Kessler leaned across and said, with real warmth, ``That was extraordinary,'' and Warren, looking out of the window at the grey river, said, ``Was it? I was only wondering about the router.'' Diane laughed, as did the junior aide, and after a moment Nora laughed too, because it was funny, because it was exactly the kind of thing he said, because the entire car understood that he was being modest. Warren glanced at her, puzzled, and smiled in the way of a man who is happy to be included in a joke whose content he cannot quite locate.

She told herself, on the way back to the hotel, that she had read the room correctly. She had. That was the difficulty. She had read it flawlessly, and the reading was wrong, and the room would go on acting on it for years, and nothing she had been taught about her own competence had prepared her for the discovery that the two could be the same thing.

\numbered{5}{The Corridor}

\textsc{the debrief took place} in the ninth-floor corridor because the suite's sitting room had been taken over by the legal team, who were arguing in low voices about the meaning of the word \emph{residency}, and because Diane Kessler, who disliked rooms in which people had already begun to take positions, preferred to talk on her feet. She stood with her back to the window at the end of the passage. Hollis Grady, who ran product, leaned against the wall with his jacket over his shoulder, and Tamsin Vo, who ran communications, held a tablet against her chest like a shield. Nora stood a little apart, and then, at a small gesture from Diane, not apart at all.

Above them, set into the ceiling panel where the corridor turned, a white disc with an amber light on its rim went on blinking at its own steady rate. Nobody looked at it. It had been there when they arrived and it would be there when they left, and in the ordinary course of things it had about as much claim on their attention as the pattern of leaves on the carpet.

``I want to talk about the router,'' Diane said.

Hollis laughed, short and delighted. ``He's doing it again.''

``He did it in the exact second she had said the thing,'' said Tamsin. ``The minister. The name on the decision. It's the best line anyone has said to us in a year, and we had no answer, and before the silence could turn into a story he'd moved the room somewhere else. I've been trying to work out whether that's skill.''

``It isn't skill,'' Hollis said. ``That's what's so good about it. He goes straight to the substrate. Everyone else in that room was talking about who answers for what. He asked what the thing was \emph{made of}.'' He shook his head with a kind of admiring exasperation. ``I've worked with him nine years and I still can't tell you how he does it.''

Diane did not answer at once. She was a tall, spare woman in her fifties, who had spent her career in the operations of large organisations and had developed a watchfulness about them that had never been mistaken for warmth and was, in its way, a form of it. She had been with Warren since the company had forty employees. She believed, with no illusions that she could have defended, that he was the most unusual person she would ever work for, and she had long ago stopped trying to say why.

``Nora,'' she said. ``You were closest. What did you see?''

It was a real question, put with real interest, and Nora felt it as a hand held out. She took her time. She was aware that the three of them were looking at her with the unforced attention of colleagues who had been in the room for years and would very much like to be told something they did not already know, and she was aware too that she had the thing to tell them, because she had already told it once, to the minister, and it had worked.

``I think the part that matters is the word \emph{building},'' she said. ``He didn't ask whether there was one router or many. He asked whether it was one per floor or one for the building. He didn't say the country. He pulled the whole question down to the level of a structure we were all sitting in, so that the dependency stopped being theirs. If the whole building hangs on a single point, then we're hanging from it too. Faruk had just spent ten minutes making it their weakness, and in about nine words Warren made it ours. It's the first time they've heard a foreign company say, in effect, \emph{we are also exposed}. And he did it by asking a question he already knew the answer to.''

``Did he know the answer?'' said Tamsin.

``He'd have to,'' Nora said. ``He doesn't ask otherwise.''

There was a short silence in which the amber light blinked eleven times. Then Diane said, quietly, ``That's it. That's exactly it,'' and turned to Hollis and said that she wanted the point written up for the team before they left the country, and Hollis said he would do it tonight. Tamsin was already typing. Diane put a hand briefly on Nora's arm.

``You read him better than anyone I've brought in,'' she said. ``I'm not sure why that surprises me. You're paid to read people who don't want to be read.''

It would have been a mistake to say that Nora was flattered, or not only that. What she felt was the particular relief of a person who has been working for years at a high level of skill and has seldom had it recognised by anyone in a position to use it, and the relief was real, and it was not a vice. She thanked Diane. She was careful not to sound as though she had been waiting to be thanked.

Behind them, at the lift end of the corridor, a night manager in a dark suit came through the fire door with a clipboard under his arm. He was perhaps sixty and walked with the stiff economy of a man who had been on his feet for most of his life. He glanced at the group, as hotel staff glance at guests, and nodded at no one in particular. Then he stopped under the white disc, tipped his head back, and checked something on it against the page in his hand, and made a short mark with his pen. He moved on towards the stairs. It took less than twenty seconds, and nobody in the corridor turned to see it.

Nora telephoned her mother that evening from the chair by her window, with her shoes off and the lights of the river laid out below like a thing arranged. Folake Adeyemi answered on the fourth ring, in the voice she used on the telephone, which was louder and more formal than the one she used in person, and said that she was well, that she had made rice, that Nora's cousin had been to see her and had talked for three hours about a van.

``And the hip?'' said Nora.

``The hip is the hip.''

``Mama.''

``I have an appointment on the fourteenth. It is only an appointment.''

It was a habit of Nora's, formed in childhood and sharpened by her profession, to hear what was under a remark, and she heard it now: a small lengthening of the second syllable of \emph{appointment}, a pause before the verb. ``You're worried,'' she said.

``I did not say that.''

``I could come.''

``Nora.'' Her mother sighed, and the sigh came down the line from eleven hours away, and was followed by a silence in which Nora could hear, quite distinctly, the sound of a kettle. ``It is an appointment, not a verdict. You always want to explain me to myself. I am sixty-eight years old, I know what I am worried about.'' And then, after another moment, in the other voice, the real one, so low that Nora had to press the phone closer to her ear: ``A little. I am a little worried. You are not wrong. But I would have said it when I was ready.''

Nora sat for a while after the call with the phone in her lap. She had been right, and she had also been told something, and she was not sure which of the two she would carry out of the conversation.

\interchapter{The Party}

\textsc{at nine, warren was} taken to a birthday party for a boy in his class named Calder Penhale, whose father owned the hardware store and whose mother had hired, for the occasion, a man in a waistcoat who made animals out of balloons. The man's name was Desmond. He had a hand pump, a battered case, and a patter that had been refined over twenty years of Saturdays and was no longer, strictly speaking, directed at anyone.

Warren did not watch the animals. He watched the pump. He stood so close to the man's elbow that Desmond, who was used to a ring of children at a respectful distance, stopped twice to reposition himself, and each time Warren moved with him. At the end of the third giraffe Warren asked how the pump worked, and the man said that you pushed the handle, and Warren asked what happened inside when you pulled it back, and the man, after a moment's hesitation, said there was a valve, and Warren asked how the valve knew which way to go. This went on for some time. The giraffe, forgotten in Desmond's hand, began very slowly to sag.

Across the room, a girl of five called Maisie Forde, whose balloon had burst during the second round of animals, was crying in the steady, uninflected way of a child who has been crying for some minutes and has decided to see it through. Nobody was attending to her. Her mother was in the kitchen. The other children were arranged in a half circle around the entertainer, and the adults, who were standing along the walls with paper cups of juice, had turned their attention to the odd sight of a boy taking apart a pump with his questions.

``He's interested in the mechanism,'' said one father, with approval.

``Look at the concentration,'' said another.

``That one will run something one day,'' said Mrs.~Penhale, in the voice of a woman who is never wrong at parties, and the remark went down the wall like a coin passed hand to hand, and by the time it reached the end of the room it had become a settled fact. Warren, in the middle of the carpet, heard none of it. He had just understood the valve, and the understanding was so purely satisfying that he turned to share it with someone and found no one looking at him except the small crying girl, who stopped, briefly, to stare.

It was Mrs.~Penhale's sister who eventually crossed the room, picked up Maisie, and carried her out to the kitchen, where there was an unopened packet of balloons and a drawer of felt-tip pens. Nobody would remember this in later years. The prodigy, on the other hand, had a long life.

\numbered{7}{The Entry}

\textsc{the negotiation for the} country lasted nine hours and fourteen minutes, and for the last four of them it took place over lunch, which had not been planned. The delegations had begun at eight in a conference room on the second floor of the partner's offices, a modern building of pale stone on the east bank of the river, and had made good progress until a quarter to eleven, when the discussion arrived at the question of what Harbor would do when the authorities asked it to remove content, and the room, which had been moving along in a spirit of mutual goodwill, stopped as though it had walked into glass.

Nora had known it would come to this. She had spent the three weeks since the ministry meeting building a structure that she believed, with the specific confidence of someone who has examined her own weak points, was the best thing she had ever designed. It rested on three ideas. The first was that Harbor would not operate in the country as a foreign company providing a service from outside. It would incorporate a local subsidiary, licensed under the country's law, staffed by residents, accountable to its courts, and holding certain categories of data on servers within the borders. The second was that this would be matched by a commitment that answered the minister's fear in the ministry's own terms: Harbor would route its traffic through two independent international gateways, one of them satellite, and would not build any function that depended on the single legacy gateway whose failure the state dreaded. It would cost money. It would also make Harbor's service, in that country, more resilient than the government's own. The third idea was the one she cared about most. Harbor would establish a safety hotline staffed by speakers of the country's languages and dialects, with authority to act on reports of threats, self-harm and coordinated abuse within a defined number of hours, and not merely to forward them.

The partner's team liked all three. They were a sharp, experienced group, led by Samir Qadri, who ran the country's largest media and telecom company and whose family had been in the business under four governments. He was a heavy, genial man with a habit of rubbing his thumb along the edge of the table when thinking, and he had read Nora's proposals the night before with a pencil. What he disliked, or rather what he could not agree to without consulting someone, was the clause in which Harbor would decline to remove content unless the request came in writing, from a designated authority, citing a legal basis, with a right of appeal.

``You understand,'' he said, gently, ``that this will be taken as a lack of confidence.''

``It is a lack of confidence in the process, not in the people,'' Nora said. ``It protects both sides. If your authorities ask for something lawful, they'll be able to point to the paper. If they ask for something that isn't, we'll both be glad there's paper.''

``The Directorate is not accustomed to paper.''

``Then it will have to become accustomed.''

He rubbed his thumb along the edge of the table. There was a pause. The Directorate of Information was the body that handled such matters, and it was headed, as everyone in the room knew without saying, by Dr.~Yusuf Qadri, the partner's cousin, a courteous man who had trained as an engineer and who had, by general agreement, no particular wish to be unpleasant. Nora knew this. It had been in the briefing papers, a single sentence, and she had weighed it and had decided that in a country of that size, with a business class of that density, almost everyone was somebody's cousin, and that to refuse to deal with the partner on those grounds would be to refuse to deal with the country. The fact had gone into her mind and had been classified as a detail.

By a quarter past eleven the room had stopped moving. Qadri's lawyers had asked for a recess, and Harbor's lawyers had asked for a longer one, and Diane, who had flown in the night before, was out in the corridor on the telephone to someone in a different time zone. Warren, who had said almost nothing since eight o'clock, had been sitting at the end of the table with his hands in his lap, watching a young man in a white jacket who came in and out of the room with a tray.

``What is that?'' he said, to nobody, when the young man came in again.

``The soup,'' said Qadri, a little blankly. ``For lunch. My people have prepared it. It's a local thing, a fish soup.''

``What's in the broth?''

There was a short pause. Then Qadri, who had in the previous two hours been guarded, courteous and formidable, smiled with something close to relief, in the way of a man who has been allowed for a moment to talk about something he understands, and said that it was his mother's recipe, that the fish was caught on the delta in the early morning, that the secret was the burnt onion, which gave the broth its colour and a faint bitterness that tasted, he said, like the smell of the river in October.

Warren asked whether the onion was burnt on a flame or in a dry pan. Qadri said, a flame, directly, on the stove ring, and that his mother had always said you could not do it any other way. Warren asked how she knew when it was done. By then Qadri had stood up. He called for the cook. The cook came, a large unsmiling woman in an apron, and the question was put to her, and she answered it at length with her hands. The delegations, who had been seated at opposite sides of the table, found themselves standing in a rough ring in the corridor outside the kitchen, listening to a woman describe the smell of an onion going black, and in the course of twenty minutes the quality of the room changed.

Nobody could afterwards say exactly how. Lunch was taken in the dining room with the lawyers present, and the conversation was not about the clause. It was about the soup, and about the delta, and about a bridge that Qadri's grandfather had helped to build, and Samir Qadri told a story about the first time he had visited Harbor's headquarters and had been unable to find a place to sit. At two o'clock he put down his spoon and said, to Nora and to nobody else, ``The paper may stay. I will speak to my cousin. I think he will agree to it, if we say it is to protect him.''

It was Diane who first used the word. She said it that evening, on the terrace of the hotel, with her feet on a chair, to a small group that included Hollis and the head of legal. ``He did the soup,'' she said. ``I've been trying to remember if I've ever seen anything like it. He took a room that was locked on principle and he got it unlocked by asking about an onion.'' Everyone laughed, and Nora, who had helped to build the structure that the soup had made possible, laughed with them. She thought that the soup had been luck, and she also thought that Warren had a gift for finding the thing in a room that people were willing to talk about, and she saw no contradiction between these.

The licence was signed eleven days later. The subsidiary was incorporated, the gateway commitment was written into the schedule, and the safety hotline, which was to be staffed by thirty-four people at launch, was given an annex of operating procedures that Nora drafted herself over two sleepless nights and that was widely admired. In the same annex, at the request of the Directorate, and after a short and cordial exchange, a provision was added allowing the Directorate to send Harbor's hotline staff \emph{regional safety notices} identifying categories of concern in the local languages, to which the staff would respond within the hours already agreed. It was a reasonable provision. It reflected a real need, since the people who spoke those dialects understood things that Harbor's engineers did not, and Nora read it three times and approved it.

She was good at this. The structure she had made would give the country a more resilient connection than it had ever had, it would give its people a place to report a threat and have the report acted upon, it would put the company under the discipline of local law, and it would do so in a form the minister could examine and the partner could defend. Anyone who read the licence in good faith and without foreknowledge would have said that it was an excellent piece of work. Nora, reading it a final time at her desk on the night before signing, thought so too, and the thought was true.

\numbered{8}{The Clause}

\textsc{the answer to the} minister's question went out on the sixth day, and it took Nora four drafts.

She had begun, in the first, by writing a name. She had thought about it with some care. The question had been \emph{whose name is on the decision}, and the plain response, the one that honoured the question, was to put a person's name at the head of the letter and say that this person would answer for what the platform did in the country. She had considered offering her own. In the first draft it read: \emph{The person accountable to you for these decisions is Nora Adeyemi, Head of Global Policy.} It was the sentence she was proudest of, and it lasted a day.

Diane Kessler came to read it in Nora's office on the seventh floor, a glass room with a view over the car park, and sat on the arm of the sofa with the printout on her knee. She read to the bottom and then went back to the top.

``I like the instinct,'' she said. ``I don't think you should do it.''

``Why not?''

``Because you'll leave. Not soon. I don't mean that. But in three years you'll be doing something else, and in five I'll be somewhere else, and the minister will still have your name in a file and a decision from year four that you had nothing to do with. A name is a promise about a person. We can only make promises about functions.'' She turned the page. ``Put the function. Make it a good one. Say it has the authority to escalate to the top of the company, say it reports in writing, say it will publish. But put the title.''

It was good advice, and Nora knew it was good advice, because it was exactly the advice she would have given herself on any other day. Persons leave, titles remain, and an institution that makes promises in the first person is an institution that will be accused, a few years later, of having made promises nobody can keep. She rewrote the paragraph. It read now, \emph{Accountability for these decisions rests with the Global Policy function, which has direct authority to escalate to the Chief Executive and which will report to your ministry in writing, each quarter.} It was a clear, committed, well-made sentence. It contained no name.

The letter went out on the sixth day, at the hour Nora had promised, and a reply came back from the ministry eleven minutes later. It consisted of a single word in English, \emph{Received}, followed by the initials \emph{I.L.} Nora looked at it for some time. It was not a rebuke. It was a record.

The sign-off memo was a different matter, and in that one the difficulty lay in a single paragraph.

Every market entry at Harbor required one. It was the document by which a recommendation became a decision, and it was addressed, by custom, to the two people whose approval it needed, and it ran to between four and nine pages. Nora's ran to seven. She wrote it over a weekend in the empty office, in the particular state of lucid, slightly feverish concentration that comes at the end of a long piece of work, and when she finished the first six pages she went for a walk round the floor, came back, and read them as a stranger might, and they were good. They described the structure. They set out the benefits, which were real, and the costs, which were manageable, and the commitments Harbor was making, which she had personally negotiated and which she believed in.

The seventh page began with the paragraph she had been dreading.

\begin{quote}
The Annex to the hotline procedures provides that the Directorate may submit regional safety notices in the local languages identifying categories of concern. The intended purpose is to allow staff with local knowledge to respond to threats and abuse that Harbor's systems would not detect. A less favourable reading is available. A notice framed as a category of concern could be directed at particular individuals or groups, and the hotline staff, who are residents of the country and subject to its law, might not feel able to decline. This risk cannot be eliminated by drafting. It can be reduced by the written-request requirement, by the quarterly report, and by our own review of notice volume. Our assessment is that, with these controls, the risk is manageable.
\end{quote}

She had written this in one sitting and had not changed a word. It was an accurate paragraph, and it was almost exactly as strong as she believed the evidence to be. What troubled her was the last sentence. \emph{Manageable} was a word that belonged to a certain kind of document, and she knew precisely what it did there: it took a finding that was open and gave it the shape of a conclusion. It did not misstate her view. It simply stopped one step short of what she actually felt, which was that she did not know.

She thought about adding a sentence. She thought of several. Each of them, written out, looked slightly hysterical beside the careful prose around it, the sentence of a person who had noticed a draught in a warm room and could not stop mentioning it. In the end she wrote something else. She added a closing clause to the paragraph, one of the ordinary, defensible clauses she had written a hundred times at the foot of other people's decisions:

\begin{quote}
Residual risk of this kind is ultimately a matter of judgment, on which reasonable people may weigh the commercial and human considerations differently, and the final assessment rests with leadership.
\end{quote}

It made her feel better at once. This was the strange thing, and she noticed it, and did not examine it. The clause did not alter the facts. It did not change the paragraph's content by a syllable. But it placed the paragraph, with an almost audible click, inside the frame of a decision that someone else was empowered to take, and Nora, reading it back, felt the small settled relief of a person who has put down something heavy on a table where it belongs.

At the foot of the final page she also set out, in a short section headed \emph{Data}, the retention terms for the hotline's operational logs: ninety days, in line with the company's standard policy, which she had approved, with a dozen other routine items, the week before. It was a sensible figure. It protected the privacy of people who had reported something frightening, and it limited the quantity of material that could ever be demanded. She did not think about it again.

Diane read the memo on Monday morning, standing, in the doorway of Nora's office with a coffee in her other hand, and she took four minutes over it. At the end of the seventh page she looked up.

``This is excellent.''

``Thank you.''

``The risk paragraph is exactly right. Honest, calibrated, and it leaves the door open for us.'' She tapped the last line with a finger. ``I particularly like this. Some people don't understand what the clause is for.''

``What is it for?'' Nora said. It was meant as a joke, or half a joke, and Diane took it as one.

``It's for us,'' Diane said, warmly. ``It's so that when we decide, we know that we decided.''

She signed the memo that afternoon. Warren signed it on Tuesday, in the margin, with a small mark beside the paragraph on the gateways, where he had written, in a hand like a draughtsman's, \emph{good, two not one}. That was all. He did not comment on the seventh page. The clause went through the building without a murmur, unremarked, and this was what persuaded Nora in the end that it had been harmless: nobody had noticed it, so it could not have done anything.

\scenebreak
\noindent In the third month, Nora began to understand that there was a second company inside the first, and that it had no org chart.

The first sign had been the car. On the morning of the delegation, three vehicles had waited at the kerb, and she had been shown to the first without anyone appearing to decide it; the seat beside Diane, the seat behind Warren, a place from which one could be heard without raising one's voice. Nobody had mentioned it. She had noticed it at once, and noticed herself noticing, and had been pleased in a way she would have been embarrassed to describe. It was a very small thing to be pleased by, and she was a little surprised to find how much of her day it had come to organise.

There were others, and they accumulated like small coins. There was the unplanned dinner, the Thursday after the launch, in the glass-walled kitchen on the sixth floor where Warren sometimes ate a bowl of noodles standing up, and to which, without announcement, a handful of people drifted over the evening: Diane, Hollis Grady, Tamsin Vo, and, on this occasion, Nora. Nobody had been invited. All of them had known to come. There was the number, a mobile that rang in his pocket and not on his desk, which Hollis gave her one afternoon as one hands over a spare key, saying only, ``Don't abuse it.'' There was the privilege of disagreeing with him in the Monday meeting, which three people had and eleven did not, and which she discovered she held on the day she said ``I don't think that's right'' and nothing happened, except that a junior director at the end of the table looked at her with a new kind of attention. And there was the knowledge itself: to hear about a decision on Tuesday that would be announced on Friday, and to find that she had already begun to arrange her face for the people who would hear it on Friday.

None of it was given as a reward. That was what she found hard to hold on to later. No one had weighed her merit and distributed a seat. It had simply been understood, as weather is understood, that she was now a person to whom certain things were said first, and the understanding sustained itself, because the people who were not told learned not to expect it and the people who were told learned to stop noticing.

On a Wednesday in that month she sat in a conference room for twenty-two minutes while six senior people discussed which of three visiting ministers should receive a personal introduction to Warren. It was a serious question. A minister who was introduced would feel favoured; one who was not would feel overlooked, and might hold a grudge; and each of the three headed a ministry with a bearing on Harbor's licence. The discussion was intelligent, and generous, and funny. It considered who had been snubbed last time, and what the third minister's chief of staff was like, and whether a handshake in a corridor would count. Nora contributed a point about timing and was thanked for it.

On her way out she passed the open door of a smaller room, where a printed recommendation from the hotline team lay on an unoccupied table. She stopped to glance at it, as she would have stopped for a leaf. It was four pages, politely worded, and it said that the staff who handled the most distressing reports in the delta region worked without any support for the effects of the material, and recommended a modest programme. It had been sitting there, she saw from the date, for nineteen days. It was not a matter of anyone's neglect. It simply had no one in the corridor who could be told first.

She picked up her coffee, and resolved to raise it, and did, in a meeting the following week, in a sentence that was well received. It was minuted. It was the kind of thing she was good at.

\numbered{9}{The Table}

\textsc{the dinner to celebrate} the approval was held on the top floor of the hotel, in a private room with a long table, a view of the river and a waiter for every four guests. There were thirty-one people present, among them the full delegation, the partner's family, two members of the ministry's staff who had come in a private capacity, and Samir Qadri's mother, a very small, upright woman in black, who had been brought at her son's insistence and who sat beside Warren, to his evident pleasure, for the whole of the first course, explaining the burning of onions.

Nora was seated at the second table from the head, between Hollis Grady and the head of legal. She had expected this. After the second course, when the plates were taken, Diane stood, tapped a glass, and said that there were people who should be thanked and that she would not name them because they would not like it, which was a form of naming; and when Nora looked across the room she found that two or three people were smiling at her. Then a screen at the end of the room lit up.

It was Hollis's doing. He had been making the slides in the car, between the ministry and the hotel, and had been unable to resist. The first slide, white on a deep blue, was headed \emph{Infrastructure-First Engagement}. Below it was a photograph of a hotel corridor with a leaf-patterned carpet and brass tulip lamps, at the far end of which a small white disc had been circled in red. Beneath the photograph, in a smaller type, were the words \emph{The Brandt approach}.

There was a ripple of appreciative laughter. Hollis explained, with the mock solemnity of a man who is enjoying himself, that the team had been asked to describe the method by which Harbor had unlocked the country, and that after careful analysis they had concluded that the method was this: you identify the structure that everyone is dependent on, you ask a plain question about it, and you allow the room to discover that the question was the real subject. He clicked to a second slide, which listed three short lines, and a third, which showed a diagram of the hotel's wiring that he had had a junior sketch from memory.

Warren sat at the head of the table with his chin in his hand, and the look on his face was one of friendly, uncomprehending interest. After a moment he leaned towards Samir Qadri's mother and murmured something to her, and she nodded.

``That's the nine-beat flutter,'' he said, to the table in general, pointing with his fork at the circled disc. ``That's the one I was looking at. It does something every ninth blink.''

Nobody heard him. The remark, which was quite audible, passed through a gap in the laughter as a sound passes through a place where many people are listening for something else. The head of legal said, \emph{Of course it does}, in a tone of reverence. Samir Qadri raised his glass and said that he proposed to put the photograph in his own office. Diane, from the head of the table, said that it was now official doctrine, and a great cheer went up.

Nora laughed. It was an easy laugh, the laugh of a person who is happy and who has no reason to withhold it. She knew, of course, where the phrase \emph{infrastructure-first} had come from. She had said it herself, five weeks earlier, in a ministry meeting, in the half second after the forty-second silence, and Hollis had written it on a legal pad in the corridor outside the room. She knew too that Warren had not meant by his question anything remotely like what the slide said, and that he was at that moment looking at a photograph of the corridor in which he had sat on the carpet in his socks, trying to understand a flutter. She noticed all of this, and it did not seem to her an important thing to say. It would have been absurd, in the middle of a toast, to rise and explain that the doctrine was a misreading. Besides, the misreading had worked. The country had said yes.

Warren, catching her eye down the length of the table, lifted his glass a little in her direction with a look that she took as thanks. It was in fact the expression of a man trying to ask, across thirty people, whether she happened to know what happened to a router at two in the morning. She raised her own glass and smiled at him warmly, and the evening went on.

\scenebreak
\noindent Nora's flat was on the fourth floor of a building with a broken lift, and on Sunday evenings it filled, without anyone having arranged it, with the people she thought of as her household.

There were three of them that night. Odile Fenn, a paediatric nurse with a dry laugh and a habit of eating standing up, had arrived first and was already at the stove, taking over the rice. Ruth Halvorsen, who lived on the floor above and was seventy-one, had come down with a bottle of something Norwegian and the expression of a woman who had been invited to a trial. Kobi Asante, a joiner, large and gentle and slow to speak, had brought a small walnut box he had made for no reason and set it on the sideboard without comment. Nora had known Odile since university, Kobi since a bad winter five years before when her boiler failed and he had arrived with tools and stayed to dinner, and Ruth for as long as the building had been hers. None of them worked in her world. This was, she sometimes thought, the reason she could breathe in it.

The conversation that night was about the launch. She had agreed to tell them about it, in outline, and she did so with some care, and a good deal of animation, because the week had gone better than she had dared hope. She described the minister, and the question, and the answer she had written, and the sentence she had been proud of, and Odile listened with her mouth full and said that it sounded like talking a parent down at three in the morning, which was the highest praise she had.

Kobi said nothing. He ate steadily. He looked at his plate and, now and then, at the wall behind Nora's head, and when she asked him what he thought he said ``Mm'' and reached for the bread.

She was aware of it at once, and she did what she always did. She read it. Kobi did not usually say very much, but he usually said something, and the silence had a texture. Perhaps he disapproved of the work: of a company, of its reach, of her part in it. Perhaps he thought she had been smug. He had grown up in a town where the government had once taken the land under his grandfather's feet, and he had never liked to see people comfortable in front of an official. By the time Odile had finished the rice, Nora had assembled a complete and coherent account of his feelings, with a history, and had decided to apologise for the smugness before pudding.

``I think I may have gone on,'' she said, in a gentle voice, carrying plates. ``About the job. I know it isn't something you'd necessarily love.''

Kobi looked up, startled. Then he put a hand to his jaw. ``I'm sorry,'' he said. ``It's a tooth. I've had it since Thursday. I can't chew on the left.'' He seemed to feel that this required more. ``The thing with the minister sounded good. I'd have liked to hear more. I just didn't want to talk, with my face like this.''

There was a short pause. Odile, who had not looked up from the sink, began to shake very slightly with laughter. Ruth said, ``Well, there it is,'' and poured herself another glass.

``I had a whole theory,'' said Nora.

``I'm sure it was a good one.''

``It was very good. It had your grandfather in it.''

Kobi stared at her. Then he laughed, painfully, holding his jaw, and said that he would have to remember that, and that it was a shame to waste it. They ate the pudding. Ruth, over the dregs of the wine, said the thing that Nora would think about afterward.

``I translated for the courts for thirty years,'' she said. ``I was good at it. I was famous for it, in a small way. What I was never asked to do was say what the witness meant. Only what they said. The interpreter who tells you what the witness meant has stopped being an interpreter and become a second witness.'' She turned the stem of her glass. ``You do it beautifully, my dear. I've watched. You read people like a document and you're usually right. I would only say that being usually right is the dangerous kind of good.''

The telephone rang at ten. It was Nora's father, who always rang on Sunday evenings at the end of whatever he was doing, and who spoke in the warm, rumbling voice that had read her bedtime stories and corrected her grammar. He asked how the important work was going. She told him, briefly, that it had gone well. He said that he was proud of her, and then, after the pause she had been braced for, that she should not take it all so hard, that she had always felt things more than other people, and that it was a gift but also a weight. He said it lovingly. She said that she knew. She said goodnight.

Afterward, when the others had gone and the plates were stacked, she stood at the window for a while with the walnut box in her hands, running her thumb along the join, which was invisible. She noticed that she was composing, in her head, a gentle note to her father about what he had really meant. She put the box down and did not write it.

\numbered{10}{Launch}

\textsc{harbor went live in} the country at six o'clock in the evening, local time, on a Thursday. This was the hour that the engineers had chosen, with some care, to coincide with the end of the working day, when people reached for their telephones.

The hotline had been set up on the third floor of the subsidiary's offices, in a room that Nora had specified in detail and that smelled, on the opening night, of new carpet and nervous coffee. There were thirty-four people in it, at desks arranged in pairs, with headsets and two screens apiece, and on the wall at the far end hung a large display that showed, in a column of unadorned numbers, the number of messages sent, the number reported, and the number awaiting review. For the first hour the third column read zero, and the staff, who had been trained for crisis, sat in an atmosphere of faintly embarrassed calm.

Nora stood at the back. Diane was beside her, and somewhere near the door Warren was talking to a young woman at one of the desks about her headset, which he suspected of picking up the fan. At twenty to eight the first report arrived: a man in the south had been sent a threatening message by an ex-employer, and the hotline staff, working from the script Nora had drafted, took it apart in eleven minutes and had the message removed and the account locked. It was an unremarkable event. Nora felt it enter her like a coin dropping into a slot. It was working.

At eleven minutes past nine, a message came through that was not a report at all. It had been posted by a woman called Leyla Haddad, who lived in the north of the country, and it was a photograph and four lines of text. The photograph showed a young man, perhaps twenty, standing in front of a blue door, and the lines said that this was her brother, that he had been lost in the floods of nine years before when she was fifteen and he was eleven, that she had not seen him since, and that if anyone knew him she would be grateful, and that his name was Tarek.

A hotline operator named Basma, who had been scanning the queue for flags, stopped on it. It had not been flagged; it had simply reached the top of the screen because she had been trained to read the first line of everything. She read it twice. Then she turned in her chair and called across the room to the supervisor, a heavy, quiet man called Idris, who came over and read it standing behind her, and who said, after a moment, in the voice of someone speaking at a funeral, ``Leave it up. Watch it.''

They watched it for forty minutes. The post was shared, then shared again, then picked up by a group in the capital, then by a group of women from the delta region who had been in the same relief camps. At nine fifty-two a man in the second city replied in the thread. He said that he knew the boy in the photograph, or rather the man, that the door was the door of a bakery on a street he could name, and that the man's name was not Tarek but Tareq, that he had been taken in by an aunt, and that he had spent nine years in the belief that his sister had died.

The room went quiet in a way that the engineers' graphs could not have recorded. Basma was crying openly. Idris had taken his glasses off. Nora, standing at the back with her arms crossed, found that she could not speak, and that her face was wet, and that she had no memory of the moment at which it had become so. She had not wept at work for eleven years. She did not try to hide it. Beside her Diane had pressed the back of her hand to her mouth, and across the room Warren, who had stopped talking about the headset, was looking at the screen with an expression of simple, open attention, the expression of a man who has just been shown what the machinery is for.

At ten past ten the woman and her brother spoke for the first time, by video, from two cities. The operators did not watch. They had turned their chairs, one by one, to face the other way. There was a long applause when it was over, and then a longer one, and Idris, the supervisor, who did not clap, lifted his hand and held it up for a moment to the room as though to say that this was enough. Hollis took a photograph of the display on his phone. Somebody opened a bottle, and was told off, and was forgiven.

Nora went back to the hotel at midnight and could not sleep. At two in the morning she gave up, put on a jersey, and went out into the corridor of the ninth floor for no reason she could have explained, only to be somewhere other than the room. The corridor was empty. The lamps in the shape of closed tulips burned along the walls at their same low intensity, and the doors went off in pairs with their cards hung out. She stood at the turn, under the white disc in the ceiling, and thought about the young man in front of the blue door.

The amber light on its rim blinked, and blinked, and every ninth beat it did something else.

She did not look up.

\partendnote{End of Part One}

\notesheading{Notes to Part One}

{\small

\emph{These notes are keyed to the text by chapter and phrase. They are not required reading. Readers who find that the novel already has enough opinions in it are advised to proceed directly to Part Two.}

\textbf{Chapter 1, ``the people who studied him.''} There were never any such people in the sense the sentence suggests. The phrase belongs to the narrator's habit of speaking from a position slightly later than the events, a habit the reader may resent and is asked to forgive. The people who studied Warren later were, in the main, a small number of graduate students in management, one documentary producer, and a retired judge with a hobby. None of them rendered his behaviour well.

\textbf{Chapter 1, ``2:07.''} The date is not recorded. Warren recorded the time and the weekday and, by the habit of a man who trusts his own memory for everything except dates, forgot to add the date itself. The stationery survives. It is kept flat, between two sheets of acid-free paper, in a drawer in an office on a floor where the drawers are not often opened. The ink has faded to the colour of weak tea. The \emph{ask} at the bottom has faded least.

\textbf{Chapter 1, ``a bird he had seen from the car.''} It was a heron. The narrator knows this because Warren knows it, and has said so more than once, and because Pia has never forgotten it. Warren has no memory of describing the heron. Pia remembers the description, which she insists included a number of details, among them the colour of the legs, that Warren says he cannot have known.

\textbf{Chapter 2, ``a Japanese maple lit from below.''} The tree was replaced twice in the first year. Maples dislike glass atriums, and this one disliked it with an evident and expensive intensity, dropping its leaves in March and holding the bare branches out over the visitors like a request. A facilities manager named Dunstan eventually saved it by lowering the temperature by four degrees and raising the humidity by twelve percent, and was thanked in an all-hands email that no one remembers and Dunstan keeps.

\textbf{Chapter 4, ``a translator, a young man with a thin moustache.''} His name was Emil. He had translated for the ministry for six years without once being introduced to the people he translated for. Nora, who noticed everything, noticed that on this occasion he translated the minister's last sentence about a half second late, not because he was uncertain but because he wanted, for a moment, to hear it in his own head before he passed it on. In the written record of the meeting this half second does not appear.

\textbf{Chapter 5, ``the sound of a kettle.''} Folake's kettle was a stovetop kettle with a whistle that she had asked her husband to repair for eleven years. He had repaired it twice. It whistled in two notes, one flat, one sharp, and she said that she had never been able to decide which was the original. Nora can hear it still, in the moments before sleep, when she is not trying to interpret anything.

\textbf{Chapter 7, ``the soup.''} The recipe is not given in the text, at the request of the family. It has been described by those who ate it as an onion soup that had been told a secret. The onions were sweated for two hours in the fat from a fish that is caught nowhere else on the river. If it has been reproduced, the reproduction was not authorised.

\textbf{Chapter 8, ``ninety days.''} It is hard to say in how many documents a number appears before it becomes a law. The ninety was chosen because the company's general counsel at the time, a careful man named Thiessen, believed that ninety days was the longest period that a reasonable person could imagine remembering something without writing it down. He was, as it turned out, entirely correct. He was also wrong about what that would mean.

\textbf{Chapter 9, ``the Infrastructure-First Engagement.''} The slide was made in forty minutes by a junior designer called Wendeline, who then went home. It has since been shown, to the narrator's knowledge, in eleven countries, twice in a university lecture and once, without attribution, at a funeral. Wendeline has never been asked what she thought of it. She thought the arrow on the left was too thick.

\textbf{Chapter 10, ``the Leyla and Tareq reunion.''} On the evening of the launch, in a café in the city of the first country, two people met who had been separated for nine years. They sat at a table in the corner. The waiter, who knew nothing, brought them water and left them alone for forty minutes, which is the best thing a waiter has ever done. Neither of them has ever said what they said.

}

\part{Diane Kessler}

\numbered{11}{The Apology}

\textsc{the outage began, as} the post-incident review would later establish, with a certificate.

It was a small one. It belonged to an internal service that signed the requests passing between two of Harbor's regional systems, and it had been issued four years earlier by an engineer who had since moved to the infrastructure group, then to a different company, and finally, according to his last known profile, to a sailing school on the coast of Portugal. The service had been built by a team that no longer existed under that name. In the second of the company's three reorganisations it had become the responsibility of a group whose charter did not mention it, and the calendar reminder for the certificate's renewal had been set up on a shared account whose owner's mailbox had been closed. At four minutes past three on a Sunday morning, Pacific time, the certificate expired, and the two systems, which had spent four years trusting each other, stopped.

Diane Kessler was woken at 3:19 by a telephone call from a person on the operations desk whose voice she did not recognise, and she was in the building by half past four. She was fifty-three and had been called out at night often enough to have a routine for it: she dressed in the clothes she had laid out, she ate a banana in the car, and she did not allow herself to form a view of the cause until she had seen the shape of the damage. The banana was a small superstition. She had eaten one before every bad night of her career and had come, without ever quite deciding to, to regard it as the cost of admission. The shape of the damage was a landlocked country in the north, population twenty-six million, in which Harbor had become, over six years and without anyone having decided that it should, the place where most people went to find out what was happening. The country's two national broadcasters were state-run and had long since lost the habit of being believed. Harbor's news feed was what remained. It was down, and it stayed down for seventy-one hours.

The building at that hour had a different body. The lobby lights had been dimmed to a lantern glow, and the great atrium tree stood in its pool of uplight like a thing that had been caught doing something private. The air smelled of floor polish and of the electric warmth that comes off machines left on all night. The cleaners had gone. Coffee rings stood on the reception desk in the shape of a small solar system. In the lift, a recorded woman's voice announced the floors in a tone of such perfect composure that Diane, who had been called out at night often enough to resent it, wanted for the first time to say something rude to it. On the fourth floor, the operations room glowed through its glass wall like an aquarium, blue and busy, with figures moving slowly in it and the low, continuous murmur of forty people trying not to sound frightened.

Nobody died of it. That was said repeatedly in the first week, in the internal review and then in the external statements, and it was true and was also beside the point. What happened was a great many small things. A farmer in the northern uplands did not learn that a bridge had been closed by flooding and drove four hours to a river crossing that was not there. A woman waiting for the result of a medical test, which her clinic sent by message, waited three days. In a city of two million, a rumour that the water supply had been contaminated ran through the neighbourhoods by word of mouth, without the correction that would normally have followed within the hour, and by Monday afternoon there were queues at the standpipes and two arrests. Diane read all of it as it arrived, on a long table in the operations room, on printouts, which she preferred to screens, and she marked them in pencil.

The communications team produced a first draft of the apology on Sunday afternoon. By Monday midday there were nineteen drafts and by Tuesday evening forty-two. They were the work of good, tired professionals who had written a great many such statements and understood the grammar of them with a precision that was close to musical. They contained the phrases \emph{we take this extremely seriously} and \emph{we are committed to learning from this incident}. They contained a sentence in the passive voice, \emph{service was disrupted}, around which Tamsin Vo, in the third version, had put a small border in red ink and written, in the margin, \emph{who disrupted it?} They all, in one way or another, spent their first paragraph on the company and their second on the customer.

Warren read them on Tuesday night, all forty-two, standing at the end of the long table with a cup of water he did not drink. He read them slowly, which was how he read everything, with the end of a pencil tracing the lines. When he had finished the last he put the pile down square and looked at it for about a minute.

``I don't think any of these is wrong,'' he said. ``I think they're all written for somebody who is not there.''

``They're written for the public,'' said Tamsin.

``I know. But the public is the farmer.'' He said it without any particular emphasis, and then, after another pause, ``Can someone find out where he was going?''

It took the operations team three hours. The farmer, whose name was Iosif Marek, was found through a regional newspaper that had written a short item about the closed bridge, and he had indeed driven to a crossing on the river Vara that no longer existed, and had waited there for most of a day, and had come home. Warren asked for the item and read it twice. Then he sat at an empty desk and began, with a pencil, to write.

The statement he produced ran to nine sentences and was recorded on a camera in the operations room at six the following morning, with the staff still at their desks behind him. It began without preamble. \emph{For seventy-one hours, you could not find out what was happening in your own country. That was our failure. A certificate that nobody was responsible for expired, and because nobody was responsible for it, nobody renewed it, and our systems stopped speaking to each other. A man named Iosif Marek drove four hours to a bridge that was gone. A woman did not get a result she was waiting for. Neighbourhoods were afraid of their own water. We caused that. Here is what we are doing: from today, every system we run will have a named owner, and the owner's name will be on a list that you can read. I am sorry.} He did not say that the company took the matter seriously. He did not say that he was committed to learning. He looked at the camera for about three seconds after the last sentence, as if checking that it had ended.

Nobody in the room spoke. Tamsin Vo was crying, quietly and professionally, behind her tablet. A young engineer at one of the desks, who had been awake for two days, had put his face in his hands. Diane, standing by the door with her arms folded, found that her throat had closed, and felt, along with the emotion, a distinct and rather cold interest in her own reaction, as she always did. She watched the faces of the staff. She had seen men and women apologise on behalf of companies for thirty years, and she knew the difference between an apology that the room endured and one that it received, and this was the second kind. It was good. It was, she thought, quite possibly the best thing the company had ever said in public.

It was Nora Adeyemi, at her elbow, who put it into words. She had been at Harbor eleven months and had been brought in from two time zones away for the incident, and she spoke softly, as though in a church. ``He's better than the process,'' she said. ``Every one of those forty-two drafts is what the process produces. That's what a person sounds like.''

Diane nodded. She did not say what she was thinking, which was that it was perhaps not safe to say aloud that the founder was better than the process, and that the process might be listening.

The statement was published at seven. By noon it had been shared more than any message in the company's history in that country. The reaction was extraordinary, and, to everyone's surprise, forgiving; a columnist in the capital called it the first honest sentence she had heard from a technology firm in a decade. The farmer, Iosif Marek, was reported to have said, in a radio interview, that he would still use the application because the man had sounded as though he meant it.

What happened next took eleven days, and it happened without anyone doing anything wrong. Warren's promise that every system would have a named owner was received by the infrastructure group, which was obliged to do something about it, and which created a working group. The working group produced a taxonomy of ownership. The taxonomy identified some forty thousand services, of which eleven thousand had no discernible owner, and it was agreed by the group that the first step was to define what an owner was. A Systems Accountability Register was established. By the end of the quarter it contained eight thousand names and was reported to the board as a major reform. Of the eleven thousand orphaned services, a little over three thousand had been assigned, and a considerable number of those to the same twelve people.

Diane followed it all with her pencil. She did not oppose the Register. It was a sensible thing, and she had more or less authorised it. But she noticed, in the way she noticed the movement of money in a budget, that a sentence which had cost the founder a night to write and a morning to say had been converted, in the space of a fortnight, into a procedure with an acronym, and that the procedure was now more durable than the sentence. The apology would be forgotten within months. The Register would be there in ten years, asked about by auditors, and whatever it measured would be what the company had decided to mean.

\scenebreak
\noindent What Diane remembered, when she thought about the seventy-one hours in later years, was not the apology. It was hour forty-four.

By then the cause had been found and the certificate replaced, and the company was in the position of a man who has found his lost keys and cannot get them into the lock. The certificate authority, the small, quiet service that vouched for the identity of every other service, was back. The moment it came up, eleven thousand services that had been waiting in the dark asked it, all at once, to confirm who they were. It collapsed under the load. They retried. It came up again and collapsed again. Each failure made every service more insistent, and the pattern was repeating at intervals of about ninety seconds, and in the operations room, forty engineers stood at a wall of screens watching a thing that looked exactly like a heartbeat, and was a patient dying.

They tried the obvious remedies. They doubled the capacity, and the herd was larger than that. They put the services in an order and brought them up in turns, but the order had to be decided by someone, and eleven thousand services have eleven thousand people who believe that theirs is the important one. A committee formed in the corner and began, with courage and good humour, to rank them.

Warren had been awake for forty hours. He stood at the back, in the same grey jumper, with his hands in his pockets, watching the graph, and he did not join the committee. He was looking at the shape of the failure. It was a spike, and then a gap, and then the same spike, and the spikes were regular, and a regular thing has a cause that is also regular. It was not that the herd was too large. It was that the herd moved together. Every service, having failed, had been told by its own code to wait ninety seconds and try again, and so all of them waited ninety seconds and tried again, in the same second, like a stadium full of people rising at the same whistle.

He went to the whiteboard. He wrote ten lines. They said that each service, on failure, should wait not ninety seconds but ninety seconds plus a delay computed from its own name: a number, between zero and ninety, obtained by taking a short fingerprint of the name and reducing it. No service would be told its turn. No one would hold a list. The names themselves would scatter the herd across the whole interval, evenly and without coordination, because a fingerprint of eleven thousand different names falls out almost uniformly. He added a line so that the delay would grow with each failed attempt. That was all.

It is worth saying what was beautiful about it, because the people in the room mostly did not say so, and the engineers who read the ten lines afterward understood it before they could have explained it. The committee had been trying to solve a problem of order. Warren had observed that the problem was one of simultaneity, and that simultaneity could be dissolved by giving each member of the crowd something that only it possessed. The fix needed no leader. It needed no list, no priority and no argument. It carried the solution inside the services themselves, so that the crowd arranged itself, and the order it chose was nobody's, and was fair.

It took eleven minutes to push to the first hundred services. The graph, for the first time in nine hours, did not spike. It rose, and rounded, and flattened, as a hill does. Over the next six hours the company came back, service by service, in an order no one had decided.

Diane watched the committee dissolve without a word, and watched the faces of the engineers, and saw something that she had not seen on a face in the building for years, which was delight. She wondered whether she would see it again. Warren stood with his hands in his pockets looking at the curve and said, mildly, that it was a little like a sleepover he had once attended, and that he would explain later, which he did not.

It was, she thought, the most complete thing she had ever seen a person do. The system had been broken, and he had mended it with ten lines, alone, and the mending was better than the thing it mended. For some hours afterward she allowed herself to believe that this was what the rest of the company needed: that the right mind, applied to the right problem, would dissolve it.

\numbered{12}{The Doctrine}

\textsc{the idea that the} hotel question was a method rather than an event originated, as far as anyone could later reconstruct, in a Wednesday afternoon meeting of the growth group, at which a junior product manager named Ottilie Barr described how she had managed to get an unresponsive team to fix a bug by asking a very specific question about a coffee machine.

It had happened in the usual way. Ottilie had been trying for six weeks to persuade a platform team to prioritise a defect in the way notifications were delivered, and had made her case in the usual language of impact and severity, and had been politely ignored. In the sixth week, after reading the story of the ministry meeting in an internal newsletter, she had gone to the platform team's floor and, instead of arguing, had asked one of the engineers whether the coffee machine in the corner ran on its own circuit or shared one with the kitchen. The engineer, startled, said it shared. She asked what happened when the microwave was running. The engineer explained, at some length, and in the course of the explanation reached the point of realising that the notifications bug was caused by an identical sort of contention, and that he had been looking past it for a month. He fixed it that afternoon.

It was a good story, and it spread, and the way it spread is worth following. Hollis Grady told it to his own team in the form: \emph{ask the oddly specific question, and let them discover why it matters}. Tamsin Vo put the phrase in a leadership email. A deck was made, called \emph{Hotel Thinking: A Practical Introduction}, in which the corridor photograph appeared on the second slide, circled in red, and the founder's name appeared on the fourth. By the end of the quarter the phrase \emph{do a hotel} had entered the language of the company, as in \emph{I'm going to do a hotel on the vendor} or \emph{can you do a hotel on that estimate?}, and a quiet engineer in the payments group was heard to remark that he had spent the morning being hotelled.

Diane watched all this with her pencil and with a kind of half-attention. She was not against it. A company of eleven thousand people needs a vocabulary that lets a junior person ask a question of a senior one without feeling that she is being impertinent, and \emph{do a hotel} was, as vocabularies went, cheap, friendly, and rather good. It gave people a script. It permitted a twenty-four-year-old to walk into a director's office and ask why a number was what it was. She had seen worse ideas become programmes at greater expense. What interested her was something else: the moment at which an incident became an institution had come, in this case, with no one noticing, between two Wednesdays.

She noticed that Nora Adeyemi was at the centre of it. Nora had not set out to be. She had been asked by Hollis, in a corridor, to say a few words at a leadership offsite about how the ministry meeting had really gone, and she had given a ten-minute talk that was clear, funny and carefully balanced, in which she said that the lesson of the router question was not that it was clever but that it was \emph{concrete}, and that concrete questions disarm people because they cannot be answered with a position. Two hundred people heard it. It was received as a revelation. Within a month Nora had been asked to design a one-day workshop, and she did, and it was a good workshop.

Diane attended the first. It was held in a converted warehouse the company used for training, a long bright room with a rough wooden floor, and twenty-four people sat in a horseshoe of chairs. Nora, standing at the front in a plain dark jacket, explained the method, using a handful of the best examples. She gave a short account of the ministry. She invited the participants, in pairs, to pick an object in the room and prepare one oddly specific question about it, and then to ask it of the person beside them.

What followed was one of the oddest sights Diane had seen in her working life. The twenty-four participants, among them a senior director of finance and a lawyer with twenty years of experience, turned to one another with grave, attentive faces, and the questions began. \emph{Does that radiator run on a loop or is it pumped from the basement?} A man asked another what the carpet was glued with. A woman asked, \emph{How many of these chairs were bought at the same time?} Each question was received in absolute silence, and then discussed, at length, with every sign of profundity. A hush of nodding came down on the room. One participant, after a long pause, said slowly that the radiator question had opened for her a whole new way of thinking about dependency, and the others murmured, \emph{Yes.}

Diane glanced at Nora, who was standing at the back with a clipboard, and saw on her face a quick, unguarded flicker of something that could have been alarm. It passed in half a second. Nora smiled and moved to the front and thanked everyone, and began the debrief, and it was a good debrief.

The promotion came that afternoon. It had been in preparation for weeks, and the timing was an accident. Nora was made a senior vice president, with the responsibility for the global policy and trust functions, and the announcement went out at four o'clock. Diane signed the paper and called her in to say it in person, and found Nora white-faced behind her desk, with the workshop materials stacked in front of her.

``I wrote the first version of that,'' Nora said. ``The method. I wrote it in a car.''

``I know.''

``It isn't what he meant.''

``I know that too.'' Diane sat down opposite. ``Nora. Does it matter? It works. Ottilie Barr fixed a bug. A man in payments asked a better question about a vendor contract than he's asked in six years. Whatever it was when he asked it, it's something else now, and the something else is useful.'' She paused. ``You built a curriculum out of a misreading. That is what most good curricula are.''

It was a persuasive thing to say, and Diane believed it. She watched Nora take it in, and watched her decide to accept it, and was aware of a small cool movement somewhere inside her own mind, like a hand adjusting a lamp, that told her it had also been said because she needed Nora to stay.

That night, because she could not sleep and did not want to think about the afternoon, Nora went looking for the phrase. She had never asked where ``the hotel question'' came from. It had arrived in her life with its meaning attached, and she had been too busy using it to wonder whether anyone had used it before. She liked to know where things came from, and the discovery that she had gone three years without asking struck her as a small domestic failure, like finding a room in one's own house one had never entered.

Someone had. In an archive of Victorian periodicals, a satirical weekly called \emph{The Tomahawk} carried a piece under that exact title, dated the fifth of September, 1868. It was a collection of letters from travellers in Jersey, Switzerland, Germany and France, and the writers were a mixed company: a retired lieutenant-colonel, a government clerk, the father of a large family, and an Irish peer writing from Dinan, who had brought his wife, six children, a governess and two maids, and who thought the bill for the party excessive. One correspondent observed that the journey to the Rhine had become wonderfully cheap, and that the hotels along it had noticed. Another, in Jersey, allowed that the food and wine were reasonable and objected only to paying nearly five shillings a day for a bedroom. None of them accused any single charge of being a swindle. Each charge, taken alone, could have been defended by the man who wrote it.

She read the whole thing twice, with the particular pleasure of finding a stranger's handwriting in a house she had thought was new. The editors closed with a confident sentence: if English holiday-makers would compare their experiences and profit by them, the solution of the great hotel question would be a matter of little difficulty.

Nora laughed, once, alone in her kitchen. It was the confidence she liked, the faith that a problem was merely a shortage of comparison. Then she stopped laughing, because the sentence was not foolish. It was what she herself believed on most days. Harbor published its prices, its policies and its complaint statistics, and made comparison easier than any institution had before. In 1868 the editors had believed that the hotel question could be settled by comparing the bills. A century and a half later the bills were perfectly comparable. It was the settlement that had become impossible.

She closed the archive. She did not put the article in the workshop materials, and she could not have said why, except that it would have made the method sound as though someone had thought of it before, and it was still, for the moment, hers.

\interchapter{The Playground Game}

\textsc{when warren was eleven} he invented a variant of tag for the yard of his school, and the game was so well made that within a fortnight the entire yard was playing it.

The rules were simple at the outset. One player was \emph{It}. Anyone who was tagged became a \emph{Lantern}, which meant that they stood where they were with one arm raised until a free player ran under the arm, at which point they were released. This was the whole of it. The innovation, which Warren had arrived at after three days of watching the yard, was a single short sentence: \emph{a Lantern may call out the name of one free player per minute.} It gave the standing players something to do. It turned the game from a chase into a conversation, with shouted names and sudden alliances, and by the end of the second week children from the senior forms had begun to turn up at the edge of the yard to watch.

The rules multiplied. A boy called Rafferty pointed out that nothing was said about walls, and Warren added a clause about walls. A girl named Jocelyn asked whether the drinking fountain counted as a place, and the matter went to a vote. By the fourth week there were nine rules, a committee of five that met on Tuesdays in the shade by the bicycle shed, and an informal tribunal, presided over by a solemn twelve-year-old with a notebook, which settled disputes. The most famous was the case of the shoe. A boy named Pim, being tagged at the moment of losing a shoe, had left it on the ground and hopped to safety; the question, which took forty minutes of deliberation and was recorded in the notebook in two pages, was whether a shoe could count as base.

Warren listened to the whole of it, attentive and patient, and at the end he proposed an amendment, which was elegant and was carried unanimously. He was happy. It was one of the happiest afternoons of his childhood. He did not notice that a small boy named Colm, who was seven and had been tagged in the first minute and had stood for most of the afternoon with one arm raised, waiting for someone to run under it, had at some point lowered the arm, put his hands in his pockets, and walked quietly inside.

\numbered{14}{The Election Pause}

\textsc{the election was in} nineteen days, in a country of fourteen million that Harbor's briefing documents described as \emph{mid-sized, democratic and fragile}, and the decision before the company was whether to switch off a single feature.

It was a small one. The feature, called Quickshare, allowed a user to forward a post to up to two hundred groups with a single tap, and it had been introduced eight months earlier for the benefit of community organisers, for whom it had proved invaluable: it let a neighbourhood association tell every other neighbourhood association about a meeting in a minute. It had also, as the data showed, begun to behave oddly during the final weeks of the campaign. A forged notice purporting to come from the electoral commission, announcing that polling stations in one province would open a day late, had been forwarded eleven thousand times in six hours before the commission's own communications officer, a woman working alone from a borrowed laptop, had managed to get a correction out.

Nora wrote the memo in a night. She did it at the kitchen table of the flat she was renting, under a lamp, with a mug of cold tea, and it was four paragraphs long. She recommended that Quickshare be suspended in the country until the day after the vote. She set out the evidence in two sentences. She set out the cost, which was real: the community organisers would be inconvenienced, a small fraction of the country's usage would fall, and the growth group would object, as would the head of product. She set out the alternative, which was to leave it on and rely on the commission's ability to respond, and she said, in plain terms, why she considered that insufficient. She read it through once. Then she got to the foot of the page, where she had written a hundred times a sentence beginning \emph{Residual risk of this kind is ultimately a matter of judgment}, and she stopped, with her pen above the paper, and did not write it.

The memo came to Diane at eight the next morning, and she read it twice. She was less interested in its argument, which she found sound, than in what it left out. She noticed the absence of the clause the way one notices a missing step on a stair.

``You've left something off,'' she said, when Nora came to see her.

``I know.''

``It isn't like you.''

``No.'' Nora was quite calm. ``I thought, this once, I'd like it to be a recommendation.''

Diane considered her. It occurred to her that she was in the presence of something she had seen only a few times in her career: a competent person doing the plain thing, deliberately, to see what happened. She did not try to dissuade her. She took the memo upstairs, and at ten o'clock she and Nora stood in Warren's glass office on the fifteenth floor, where he was eating a pear.

He read it slowly, as he read everything. He put the pear down. He asked two questions. The first was how many of the community organisers' messages would still be deliverable by hand, and Nora said, nearly all, at the cost of a few taps. The second, after a long pause, was whether the forged notice had spread because of the two hundred groups or because of the one tap. Nora said she thought the tap. He nodded, slowly, and said, ``Then what if we kept the two hundred and added the pause?'' and Nora said that was a different feature, and that it would take eleven weeks, and he said, ``Right. Then switch it off,'' and handed the memo back.

It was the work of nine minutes. Diane, who had been present at four hundred decisions of comparable size, found herself noticing that it was \emph{clean}. There was nothing in it to wade through. A person had proposed a thing in plain words and another person had asked what a careful person would ask, and the thing had been done. She had a brief and disquieting thought, which she put away, about how rare this was.

The pause held. Hollis Grady objected, loudly and well, on behalf of the growth group, and the head of product objected, in writing, and the organisers sent a letter, and none of it moved anything. For nineteen days Quickshare was dark in the country. The forged notices continued to appear, as they were always going to, but they travelled at the speed of ordinary people, and each of them was met by a correction before it had reached a thousand. There was a second forgery, more skilful than the first, which was shared four hundred times and then died. On the day after the vote the commission issued a short statement thanking the platforms for their cooperation. The election was held, close and contested, and its result was accepted by both sides.

Nora was, for a few weeks, the hero of a certain kind of internal presentation. Tamsin Vo's team put out a short film. The head of product, in the grudging style of a good loser, sent her a note that said \emph{I was wrong and so were you about the cost}. Warren said nothing publicly, but he sent her, through the post, a postcard of a lighthouse on which he had written, in his draughtsman's hand, \emph{thank you for writing the paragraph without the other paragraph}. She kept it on her desk.

It was Diane who gave the episode its meaning, at a leadership dinner some weeks later, in the course of a short toast. She said that the company had spent two years building a process for deciding things, that the process had been laughed at by every new hire, and that the election was the proof that the process worked. A person had a concern, the concern went up, and the right people made the right decision. She meant it. It was the sincere conclusion of a woman who had spent her life building the thing she was praising, and there were nods around the table, and Nora, at the far end, raised her glass with the others and smiled.

She did not say what she was thinking, which was that the memo had worked, and that she could not have said why she had not written the others the same way.

\scenebreak
\noindent Nora won the Meridian agreement in four days, in the garden of a man who kept a stuffed heron.

The man was Count Ludovic Varnay, financier, chairman of the holding company that owned Meridian Telecom, and the possessor of a reputation for eccentricity so thorough that his lawyers had stopped trying to correct it. He conducted his negotiations on a croquet lawn. The lawn was a rectangle of imported turf in the middle of a dry plain, watered each night from a tanker, and it was surrounded by a low hedge of clipped juniper in which, at intervals, stood small statues of birds. At the head of the lawn, on a plinth, stood the heron, wearing a very faded ribbon. Each concession, the count explained, had to be announced to the heron before it was written down, so that nothing be agreed that could not be said aloud to a witness who could not be bribed.

It was absurd. It was also, Nora saw within the hour, completely serious. The count had made his fortune in telecommunications in three countries and had lost two wives and a great deal of faith in the people who came to negotiate with him, and the croquet was a way of keeping them in the sun, slightly off-balance, long enough to see who they were. His own team consisted of a nephew, an elderly lawyer with an ear trumpet, and a woman called Perpetua who kept the score and, Nora suspected, ran the whole enterprise.

What Harbor wanted was simple. Meridian carried fourteen million subscribers in the delta country, most of whom paid for data by the megabyte, and who therefore used Harbor sparingly, like salt. Harbor wished the core services to be free on Meridian's network: no charge for the data they used. The count wished for something he would not say. For three days Nora played croquet, badly, and listened, and announced a series of concessions to the heron, each of which the count received with grave courtesy and none of which moved him.

On the third evening Perpetua, adding up the scores in a notebook by lamplight, said without looking up that the count's first wife had been a surveyor, and that she had mapped the plain on foot, and that he had once proposed to name a building for her and been refused by the ministry on a technicality. Nora went to bed and did not sleep. In the morning she proposed, in a sentence addressed to the heron, that the new data centre which Harbor was to build in the region be named for the surveyor.

The count stood perfectly still. He had lowered his mallet. He said, after a time, in a voice with a different register, that it would be a great honour to her memory, and he put out his hand.

She signed that afternoon. It was an excellent agreement, and she knew it as she signed. Fourteen million people would use Harbor without counting, and Harbor would be, for them, the first thing on the phone and the last. She flew home in a state of lucid, shining pleasure. She had found the lock, in a garden, among statues, by listening to a woman with a notebook, and the door had opened. She told Warren on the telephone, and he said that it was extraordinary, and she heard in his voice that it was, and she felt it somewhere in her chest like a held note.

The letter arrived six weeks later. It came from a young man in the delta country called Pavle Dimitri, who had written, with a few friends, a small messaging and forum application called Kolo, popular in the region's towns, which charged nothing for its basic use but depended on people being willing to spend their data on it. He wrote politely. He said that since the Meridian agreement, his users had left. They had not left because Harbor was better. They had left because it was free, and Kolo, however slight, cost them something to open. He did not say that this was unfair. He asked, in the last paragraph, whether Harbor might consider extending the same arrangement to local services, since it would be a great kindness to a community.

Nora read it twice at her desk. It was a reasonable request. The agreement had a clause, a short one, which allowed Harbor to propose such an extension, and Meridian to refuse it, and the refusal, she knew at once, would be given. The count had no interest in subsidising a rival to the thing for which he had been paid. She had not written the clause to exclude Pavle Dimitri. She had written it to be fair. She sat for a while with the letter, with the unfamiliar sensation of having done something excellent and being unable to find the fault, and having been in the wrong all the same. She drafted a reply. It was a good reply. It said that she would raise the matter.

\numbered{15}{The Listening Session}

\textsc{the three visitors arrived} on a Tuesday in March and were shown, by an assistant who had been instructed to be warm and had interpreted the instruction generously, into a small room on the fourteenth floor that had been furnished with a round table, three chairs facing a fourth, and a plate of pastries that nobody touched. They had flown eleven hours. They had come at the company's expense and at the company's request, and they had agreed to come, as they separately admitted afterwards, mostly out of curiosity, and partly because it is difficult to refuse an invitation from a person who has the power to ruin your work.

The first was Ilse Varga, a fact-checker who ran a four-person verification desk in the capital of the delta country and who had spent six years being polite to the platform in public and furious at it in private. She was in her late thirties, small and exact, and she had brought a folder. The second was Dalia Rusu, a journalist of the older kind, heavy-shouldered and unhurried, who had covered three elections and two floods and who had opinions about Harbor that she had been advised not to print. The third was a former judge, Anselm Okoye, seventy, silver-haired, soft-spoken, who had retired from the bench of the country's second court and had since become, to his own surprise, a person to whom ministries and newspapers turned when they needed someone who could say that a thing was lawful or was not.

Warren came in at ten past ten, alone, without an assistant, carrying a notebook of the kind sold in stationery shops, with a green cloth cover and a pencil in the spine. He apologised for being late by ten minutes. He sat in the fourth chair, put the notebook on his knee, and said that he did not want to say anything for the first part of the meeting. He wanted to hear.

Diane sat in the corner, where she had been told by Warren in advance that she might, on condition that she say nothing. She had the sense, which she had at his meetings with some regularity, of watching something she had no category for. She had sat in perhaps six hundred sessions of this kind, and they had a standard form. A senior person listened for twenty minutes with the expression of one who is attending, then spoke for forty minutes about the company's commitment, then invited questions that were answered by a colleague. This was not that. Warren listened, and wrote. He wrote in a small upright hand, and when one of the visitors said something that he wanted more of he looked up and said, ``Can you say that again?''

Ilse Varga went first. She spoke for twelve minutes without once looking at her folder. She described how a false claim moved through her country, from the first post to the corrective, and how the gap between the two had widened by a third since Quickshare was introduced, and how her desk had stopped trying to answer every claim and had begun to triage by a method that she was not proud of. She said that Harbor's own tools for flagging misinformation sorted claims by language, but that her country's most dangerous rumours were not in one language, they were in a hybrid that the villages used and the classifiers did not recognise. She said that she had written to the company six times. She had received two replies.

When she finished, Warren said, ``Which hybrid?''

She told him. He wrote it down. He asked her what the villages called it, and she said that they called it nothing, it was just how they talked, and he nodded and asked whether she could send him ten examples, not of misinformation, just of ordinary messages in it, and she said she could. He asked how she had known the classifiers did not recognise it. She said she had tested them herself, with a list. He asked how long the test had taken, and she told him, four days, and he said, ``I'm sorry. That's four days of your life we should have spent,'' and wrote something else down, and Diane saw Ilse Varga look at him with an expression of puzzled suspicion, as at a door that has opened too easily.

Dalia Rusu spoke next, and then the judge. Each spoke for the better part of a quarter of an hour. The questions Warren asked them were, as his questions generally were, narrow and structural and a little unnerving in their accuracy: to the journalist, how a story got from a regional paper into the platform's feed, and who in the middle could change its position; to the judge, what happened, in practice, when an order from one of the country's courts was served on a foreign company, and who the court believed it was serving. The judge answered this one slowly. He said that the court believed it was serving the company, and that it would be useful if the company believed so too. Warren said, ``I don't think we do. I think we think we're being served by a department,'' and the judge, after a moment, laughed softly and said that he had never heard a chief executive say so.

They finished at five to twelve. The visitors, who had expected a ceremony of regret, had been asked for ninety minutes of expertise, and had given it, and they went out into the corridor with a lightness that Diane had seen before only in people who had been properly examined. Dalia Rusu, waiting for the lift, said to Diane, in an undertone, ``I came to be angry with him. It's very inconvenient.''

By the end of the week Warren's notebook had been typed up. The typing was done by an assistant in the strategy group who, finding forty pages of small upright handwriting punctuated by questions in the margin, understood that she was holding something important and treated it as such. A summary was circulated. A senior vice president in the trust function, reading it, observed that it would be a pity if the insights were lost, and proposed that they be captured in a programme; by the following Thursday there was a Stakeholder Insight Series. By the second month it had nine working groups, each chaired by a director, covering language, correction speed, court orders, regional press, community moderation, data access for researchers, appeals, trust metrics, and a ninth whose title nobody could subsequently remember. A quarterly report was commissioned. It was to be called \emph{What We Heard}.

The first issue ran to sixty-two pages. It was illustrated, attractively, and it contained a section on each of the three visitors, with a photograph and a short quotation chosen from the transcript. Diane read it in the car on the way to the airport and was surprised to find it accurate. It did not distort what had been said. It did not soften the judge's remark about the department. But every page of it had been written to be read by a person who had not been in the room, and none of the sixty-two read like what had happened there, which was that three people had been listened to by a man with a pencil.

She thought of Ilse Varga's ten examples, which had been sent within the week and had been logged by the language working group as item 41 of a backlog, and which, as it happened, were still there in the autumn.

\numbered{16}{The Translation}

\textsc{the expansion into the} hill districts had been planned for a year and was, in the language of the briefing documents, a \emph{tier two linguistic environment}, which was a polite way of saying that three languages were spoken in it, none of them well served by the existing machinery.

Nora had been preparing for it for eight months. She had argued, against the head of product, the head of growth and, at one stage, in a long and unpleasant meeting, the head of finance, that the company should not enter the region with its standard product translated at the last moment by a vendor. She wanted native-language review of every interface string. She wanted local moderators hired from the districts, not flown in. She wanted a pilot in one town before any general release. The finance argument, which she had expected, was about cost, and was answered with a number. The product argument, which she had also expected, was about speed, and was answered by Diane, who sided with her, in the end, with the remark that the company had learned in the delta country what it cost to arrive without understanding.

The translation exercise was the part of it that she remembered afterwards with the most affection, and the most discomfort. Three translators had been engaged, one for each language, from three different agencies, and each had been given the same glossary of two hundred and fourteen key terms and asked to produce an equivalent. For the large majority of terms they had agreed within a day. For eleven they had not, and of those the most troublesome was the first on the list, a word that appeared in the company's mission statement, in the sign-up screen, in the community guidelines, and in the name of the feature that let users form groups: \emph{community}.

The Tavish translator offered a word that meant, literally, \emph{the people one sits with}. The Orrun translator offered a word that meant \emph{those who share a river}, and explained that it carried an implication of obligation, because the river was shared and so were its floods. The Selk translator offered a word that meant, more or less, \emph{the household and its debts}. Each was a good word. Each was a real translation, in the sense that a speaker of that language, hearing it, would recognise the thing. And each was a different thing. Nora spent a week on the problem. She convened the three translators on a video call and asked them to argue, and they argued, with great dignity, for two hours, and at the end of it she knew a good deal more about the three cultures than she had at the start and was no nearer a decision.

It was decided in the end by a slide. Hollis's team had prepared, for the launch event, a presentation in which the Tavish phrase had been used, because it had been the first to arrive; the slide had been shown to Warren, who had liked it, and it had been shown to the head of product, who had said it was beautifully simple. By the time the matter reached the table it was in the deck. Harbor shipped the word for \emph{the people one sits with} in all three languages, and the Orrun and the Selk, who had read it and found it strange, were told that a unified term was important to the brand.

Nora raised no objection. She had lost a battle, she told herself, in a war whose aims she shared, and the point was to preserve her capital for the issues that mattered. She was never to know how often she told herself this, because it was the kind of reasoning that does not record itself.

The pilot in the town of Vessen began in June. It was a town of thirty thousand, at the bottom of a steep valley, with a market square, a half-built cathedral and two rival football clubs, and Harbor's moderation team there consisted of nine people, all from the district. In winter the valley held the smoke of its stoves like a held breath, and the light came late over the ridge and left early, so that the market opened in blue and closed in blue and the colour of the day was a rumour. The lead moderator, Mirela Costa, was thirty-four, a former schoolteacher, who spoke all three languages and a fourth that was not on anyone's list, and who had been hired by Nora, over some objection, after an interview in which she had been asked how she would handle a rumour and had replied that she would first find out who in the village the rumour had been made \emph{for}.

On a Thursday in the pilot's seventh week, at ten past six in the evening, a message appeared in a Vessen group saying that a boy from one of the football clubs had been killed by men from the other. It was false. The boy had been taken to the hospital with a broken wrist after falling off a wall. But the message was in the hybrid register of the valley, with a particular verb ending that implied the writer had seen it, and it had been forwarded forty times in six minutes. Mirela saw it at 6:14. A classifier would not have flagged it, because there was nothing in it that a classifier could recognise as hostile. She read it as a person reads a message from a neighbour.

She did not take it down. She telephoned the boy's mother, whom she knew from the school, and got the whole account in four minutes; then she posted a short message in the same register, in the same verb ending, saying that she had spoken to Danil's mother, that Danil was in the hospital with a broken wrist, and that he had asked for his football boots. She posted it on the group where the rumour had started. She posted it at 6:31. By 6:45 it had been forwarded more than the rumour. The next morning there were no crowds in the square.

The field report was forwarded to Diane on the Friday, with a note from Nora that said only, \emph{this is what it's for}. Diane read it on the train to the airport, twice, and thought of the three translators and the word for \emph{community}. She thought that Mirela Costa had known, as no classifier did, what the people of Vessen were, and that it was precisely this knowledge that the company had set out to buy, and had paid for, and had very nearly decided it could do without in a matter of a single slide. The pilot had been the kind of thing that made a person believe in the work. She believed in it. She also observed, in the cool inner voice that accompanied her on trains, that there were a great many towns like Vessen, and a good many more than nine moderators, and that no number of Mirelas could close the distance between a company that had understood a valley and a company that had decided what a valley was.

\interchapter{The Science Fair}

\textsc{at twelve, warren entered} the school science fair with an apparatus that took him five weeks to build and an hour to explain, and the subject, which nobody had assigned him, was the question of how long a candle burns under a jar.

He had noticed the matter in a kitchen, at a power cut, when his mother had placed a drinking glass over a tealight to keep off a draught and the flame had gone out after a few seconds. He wished to know how few. He measured. He tried twelve sizes of vessel, from a thimble to a preserving jar, and for each he timed the burn forty times with a stopwatch and a mirror, recording the results on graph paper in a hand so small and exact that the judges later had to use a magnifying glass. The apparatus itself consisted of a plank, a rack for the vessels, and a system of strings and a brass bell by which the candle, on going out, released a weight that rang. It worked perfectly. On the Saturday of the fair, he stood behind it in a clean shirt, in a hall that smelled of floor polish, and waited.

The first judge, a retired chemist, spent twenty minutes with him and was thorough and kind. The second, a woman from the education office, asked about the bell. The third, who had been standing back with his hands behind him, approached the display last, read the title card, looked at the rack of twelve vessels, and asked the question he had been working up to.

``Why did you choose jars?''

Warren considered it. He had the sense that this was an important question, because the man had a serious face, and he wanted to give a true answer. He thought back to the day in the hardware shop. ``They were on sale,'' he said.

There was a short silence. The third judge raised his eyebrows. The second judge, who had been writing, stopped writing. In the quiet that followed, Warren thought he ought to add something, and he said, ``There was a box of twelve for four pounds. I couldn't have bought twelve different kinds of glass for that.''

``Ah,'' said the third judge, softly. He turned to his colleagues. ``A constraint of resources, chosen as a design principle.'' The retired chemist nodded slowly. The woman from the education office wrote a line and underlined it. ``You limited yourself,'' the third judge said, turning back, ``to what could be replicated by any child in any kitchen. That is precisely the point of an experiment.''

Warren said that he supposed it was. He took first prize.

\numbered{18}{The Hearing}

\textsc{the hearing was scheduled} for a Wednesday in October, in a high-ceilinged committee room in the capital, before a panel of fourteen senators chosen, in the usual way, to represent the major positions in the chamber, from the sympathetic to the openly hostile. The room smelled of lemon oil and wet wool, and the microphones, which were old, gave a small electric cough each time one was switched on, like a throat being cleared in a church. It would run for four hours. It would be broadcast, and the clips of it would be on the news by evening. The company had been preparing for it for three weeks.

The preparation was done in the large conference room on the fifteenth floor, which had been cleared of everything except a table, a lectern, and eleven interns. They were bright, young, ambitious, and they had been told to be vicious, and they took to it with an enthusiasm that surprised even their supervisors. One of them, a lanky twenty-two-year-old, played the part of the senior senator from the opposition with a gravity so complete that he had adopted a limp. Nora ran the sessions. She sat at the end of the table with a stopwatch and a notepad, and interrupted, and made him do it again.

She had expected the preparation to be hard, since Warren had always been bad at being managed. She found that it was easy. He had changed. He had founded the company fourteen years earlier; for the last two he had been, she thought, quietly working on himself, in the way that some people do after a shock, and the results were visible. He answered directly. He did not wander into the mechanism of a router when asked about a fine. When the intern-senator asked him whether he took responsibility for the outcome of the election in a particular province, he paused, and said, ``I take responsibility for the decisions that I made. I don't think I can take responsibility for the result, because I don't know enough about what we did and did not do. I'd like to know more than I do,'' and the interns, who had been trained to pounce, found that there was nothing to pounce on.

``That's the best answer you've given,'' Nora said, writing. ``Don't change it.''

He looked pleased, with the open pleasure of a man given a grade.

The difficulty, as always, was legal. The company's general counsel, a courteous, silver-haired man named Brendan Pryce, sat in the corner of every session and spoke seldom and decisively. On the Monday before the hearing he drew Nora aside and said that there was a question the senators would almost certainly ask, concerning the hill districts, where tensions between two communities had begun to escalate in September and where there had been, in the previous week, two killings. The question would be whether the platform had contributed. The honest answer, as Nora and Brendan both knew, was that nobody yet knew. There was a team looking into it. Warren's instinct would be to say so.

``I don't want him to say that,'' said Brendan. ``Not on the record. \emph{We don't know whether we contributed} is a sentence a plaintiff would frame. I want him to say what we are doing.''

``That's true,'' Nora said, slowly. ``It's also what he'd say if he were a person.''

``He is a person. He is also a defendant.''

She wrote the statement that evening. She wrote it with care. It said that the company took reports of violence seriously, that it had convened a review of its operations in the affected districts, that it was working with local civil society partners, and that it would share findings with the committee as soon as they were available. Every sentence of it was true. It was, she thought, a good statement of its kind. It said what the company was doing, and did not say what it did not know.

The hearing went well. This was the plain truth of the matter, and the clips on the evening news bore it out. Warren was calm. He was courteous. He answered questions about fines and about Quickshare and about a child-safety measure with a clarity that made two of the hostile senators visibly rearrange their notes. At a quarter to four, the senior senator from the opposition, a gaunt, formidable woman named Halloran who had made her career on asking one question and waiting, leaned into her microphone.

``Mr.~Brandt. In the hill districts. Two men are dead. Did your platform play a role?''

Warren looked at the paper in front of him. It was the statement. He read it aloud, steadily, in its entirety. Then he put it down, looked at the senator, and said, in his own voice, ``I would like to know more than I do. I'm sorry that I don't.''

It was an extraordinary thing to say, and it was said sincerely, and it landed. The room was quiet. The senator held his gaze for a moment, and then nodded, once, and said that she appreciated his candour, and turned to the next item on her list.

The coverage the next morning was, by the standards of such things, warm. A newspaper that had been hostile for two years ran a headline that said \emph{Brandt, Unusually, Admits Limits}. Brendan Pryce sent a congratulatory note. The interns were given a lunch. Diane, reading the clips in the car, noticed with a slight chill that nobody had followed up. Not the senator, and not the journalists, and not the committee, whose staff had received the statement and filed it. The question about the hill districts had been asked, answered, and received with appreciation, and the appreciation had closed it. It was gone from the cycle by evening.

Nora, for her part, felt the satisfaction of a hearing well prepared, and under it a faint, persistent discomfort that she did not trace. She had written a true statement and a true statement had been read and a true sentence had been added to it, and the sum of these truths had been that nothing was asked. She went home, ate, and slept well.

\scenebreak
\noindent Diane's husband was a retired finish carpenter named Lorne Kessler, and in forty years of marriage he had never once asked her how work was. He asked, instead, what she had had for lunch, and she had learned that it was the same question.

She came home that evening at eight and found him in the garage with the door open, planing the edge of a piece of ash held in the vice. The shavings lay about his feet in long pale curls. He did not look up. He ran his thumb along the edge, felt something she could not feel, and made one more pass with the plane. Then he put it down and asked what she had eaten.

``A sandwich,'' said Diane. ``At my desk.''

``Good one?''

``I don't remember.''

He wiped his hands on his apron. He was seventy-three and had hands like old tools, the knuckles enlarged and the pads thick, and they had hung four thousand doors in the city's better houses, each of which opened without a sound. She had loved those hands for decades. She sat on an upturned crate by the bench and watched him sweep the shavings into a pile with the side of his foot.

``You've got a face on,'' he said.

``There was a hearing, last month. I haven't stopped thinking about it. A senator asked Warren something that I'd have thought was a simple question, and Warren said he'd like to know more than he did. And it was true. I sat there and I thought, Lorne, nobody in that building knew the answer. Forty thousand people, and no one.''

``And?''

``And I'm the one who built the way we know things. All of it. The reports, the reviews, the registers. I bought the chairs, at the start, and then I bought the process to go with the chairs.'' She laughed shortly. ``I spent twenty years making it so nobody could be blamed for not knowing. It worked.''

Lorne considered this, picking a curl of ash off his sleeve. ``I never had a client,'' he said, ``who asked me whose house it was.''

``What?''

``It's a funny thing. In the trade you always know. There's the owner, there's the architect, and there's you, and when the door sticks it's yours, because you hung it. Nobody has to write that down.'' He turned the piece of ash and sighted along it. ``The ones that went wrong, in my experience, were the jobs where the house belonged to a company. Developers. A committee. The door'd stick, and I'd ring the number, and a very nice young woman would tell me she'd raise it. And six months later it still stuck.''

``That's not the same thing.''

``No,'' he said agreeably. ``But I'll tell you what I noticed. In the old jobs, I got called back. A man'd phone me at home and say, Lorne, the door. In the new ones, nobody phoned. And I used to think, well, that's fine, they're not having trouble. And then I'd drive past, and there'd be a nice new door on the same frame, put in by somebody else.'' He put the plane back on its hook. ``They'd never told me. There was nobody to tell.''

She said nothing for a time. The garage smelled of cut wood and cold metal. It was what she had come home to every night for forty years, and she realised, with a slight shock, that it was the one room in her life that did not have a register.

``You're the only person I know,'' Lorne said, not unkindly, ``who's never been wrong in a way that cost anyone anything. That's a hell of a thing to be proud of, Di. I'd be careful.''

``I'm not proud of it.''

``No,'' he said, and switched off the light over the bench. ``That's what worries me.''

\numbered{19}{The Metrics Room}

\textsc{the civic integrity dashboard} had begun life as a spreadsheet kept by a single analyst, and by the time it was shown to Diane Kessler on a Thursday in November it occupied a dedicated room on the twelfth floor, with a wall of screens, a standing desk for the presenter, and a long table at which sat the people who were supposed to be able to read it.

It tracked four hundred and six metrics. They covered the speed at which false claims were corrected and the share of reports resolved within a day. They covered the fraction of political advertising that had been labelled, the number of coordinated accounts removed, the median age of an unresolved appeal. For each of Harbor's hundred and forty markets there was a row, and for each metric a column, and in the cell where the two crossed, a number, a colour, and a small arrow. It was an extraordinary piece of work. The wall of screens gave off a faint high whine that its occupants had stopped hearing in the first week, and which visitors noticed only when the displays were switched off and the silence arrived like a change in pressure. Diane, who had spent her career among instruments, admired it without reservation, and would have said that it was the most complete picture of a platform's civic health that any company had ever assembled.

It was also unreadable, and everyone knew it. A dashboard of four hundred and six metrics cannot be read by a person with an hour. The review that Thursday was the culmination of a six-week effort, directed by a vice president named Gareth Linde, to reduce the dashboard to a form that executives could use. The brief had been a single sentence at the top of a page: \emph{Executive attention is the scarce resource; surface what matters.} The result was eleven screens, each with a theme, a headline, and nine or ten figures chosen by a team of three, and the presentation of these eleven screens, by a young analyst named Priya Raman, was the item on the agenda.

Priya was twenty-six. She had been at Harbor for eighteen months, having come from a public health institute where she had spent three years building models of how infections moved through crowded cities, and she carried herself with the careful stillness of a person who has learned that being right is not enough. She stood at the standing desk with a clicker in one hand. She was very good. Diane saw this within a minute: she had a gift for stating a number and its limits in a single breath, and she did not pad.

The screens went by. The first showed the median correction time across all markets, which was down eleven percent. The second showed a map of the world, shaded in a gentle gradient from green to amber, in which nothing was red. The third was a trend line. The fourth, headed \emph{Regional Health}, showed eight regions, each with a single composite score between zero and one, and the scores ranged from 0.74 to 0.91, and the room, which had been alert, relaxed a little.

``Region Four,'' said Gareth Linde, reading from the screen. ``Eighty-two. That's the north and the hill districts?''

``The north, the hill districts, and the delta,'' said Priya. ``Fourteen markets.''

``And that's stable.''

``The composite is stable, yes.''

Diane had been watching Priya's hand, which was on the clicker, and which had tightened very slightly on the word \emph{composite}. She did not know why she noticed it. It was the smallest thing, the movement of a single knuckle, but she had been reading faces and hands for a very long time, and she leaned forward.

``Priya,'' she said. ``What's under the composite?''

There was a pause. It lasted about as long as it takes to decide whether to be brave. Then Priya said, in the level voice of a person reading from a prepared statement, ``Within Region Four, the composite is a weighted average of fourteen markets, each of which has between nine and thirty-one component metrics. For the hill districts, which are market seven, twenty-two of the thirty component metrics are not currently being collected. The data pipeline there is partial. When a component is not collected, the cell is grey on the detailed view, and in the roll-up it is excluded from the average rather than counted as zero.''

``Meaning what?''

``Meaning the composite for the hill districts is calculated from the eight metrics we do have. Those eight are the ones that are collected automatically: uptime, latency, ad labelling, account removal. They're good. The ones we don't have are the human ones. Report resolution time, local-language correction speed, the appeals backlog. Those are the ones that would show a problem.''

``So the score is high because the bad numbers are missing.''

``The score is high because the missing numbers aren't in it.'' She said this carefully, as someone says a thing for the second time. ``I raised it in the working group in September. The decision was that the regional composite should only include collected metrics, so as not to penalise a region for a pipeline problem.''

Nobody spoke. The wall of screens glowed. On the large display at the front, in the detailed view for market seven, the grey cells stood in a tidy column, and at the foot of each a small caption read \emph{No data}, in a font chosen for its calm. The room had looked at that column for an hour and had read it, as Diane now saw, in the way a room reads a blank in a ledger: as nothing to report.

``I think,'' said Gareth Linde, slowly, ``that the confusion is in the labelling. We'll make it clearer.''

``Of course,'' said Diane. She did not look at him. She was watching Priya, who had lowered the clicker and was looking at the floor, and in her face there was an expression Diane would remember: the quiet dismay of a person who has done exactly what she was told, and done it well, and has seen what it made.

It was Nora who stayed behind, after the others had left, and who stood at the back of the room looking at the grey column on the display. Diane stayed too. She asked, in a low voice, who had decided that a missing number should be excluded, and Nora said that she thought it was a sensible rule, in general; that if you counted a missing number as zero you would penalise exactly the places with the weakest infrastructure, which were also the places with the greatest need. She said that it was a good reason. She said that Priya had been right to follow it, and right to say what she said.

``It's a good reason,'' Diane agreed. She was looking at the eleven screens, each with its quiet headline. ``That is exactly what's wrong with it.''

She did not explain. She was not sure that she could. But she had begun to think, in the weeks since the hearing, about the way a reduction works: that to compress four hundred metrics into eleven screens is to decide, forty times, what does not matter, and that the deciding is never done by anyone as an act. It is done by a brief, a template, a rule about missing data. No individual in the room had pruned the hill districts out of sight.

\numbered{20}{The Memos}

\textsc{nora found them on} a Saturday, and she found them by accident.

She had gone into the office to work on the annual policy review, a document she disliked and did every year, and which required her to gather, in one place, every memo her function had produced in the preceding twenty-four months. The archive had a search interface. She typed her own name and the date range into it, and the system returned thirty-one results, and she did what the review required, which was to open them one after another in a row of tabs on the long screen. She then did something the review did not require. She printed them. The printer in the corner of the floor was an old workhorse that warmed up with a groan, and the first page came out warm, with the smell of toner and scorched dust that Nora had always associated with other people's urgency.

There was no particular reason. She was a person who read better on paper, and the annual review needed a margin of her own annotations, and the printer in the corner of the floor was a good one. It took eleven minutes. She laid the pages out on the conference table in the order in which they had been written, thirty-one memos in thirty-one small stacks, and stood at one end with a mug of coffee, and looked along the table.

They were good. She noticed it first. She had forgotten how good they were. They were clear, balanced and exact, and they dealt, one by one, with the questions that had defined her two years: a partnership in the delta country, a change to the appeals process, the hill districts, a feature for sharing images of identity documents that she had thought, at the time, was a mistake. She read the first paragraphs, as she always did, and then she turned over each stack to its last page, as a person turns over a card to look at the back. And she saw that the last paragraphs of the memos were the same.

They were not identical. She was too careful a writer for that. But they had the same bones, the same cadence, and the same operative phrase, which recurred with the regularity of a refrain. \emph{Residual risk of this kind is ultimately a matter of judgment.} \emph{Reasonable people may weigh these considerations differently.} \emph{The final assessment rests with leadership.} She walked down the table and read them one after another. Twenty-two of the clauses were in some form of this wording. Eight were softer, and began with \emph{We recommend that leadership consider}. Each of the thirty ended, in the plainest sense, with a sentence that handed the decision to someone else.

She counted again, standing very still. Thirty. And then, at the far end of the table, on its own, the thirty-first, the election pause, four paragraphs, no clause.

For a long moment she did not feel anything. It was an odd, clinical interval, like the blank after a blow, in which her mind went about its business and assembled the facts in order. Each of the thirty memos had been thanked. She remembered the thanks, which had been warm and specific: \emph{this is a very helpful framing}; \emph{thank you for laying out the options so clearly}. Each memo had been read. Each had been taken seriously, in the sense that a thing is taken seriously when it is placed in a file with a date on it. And in each case the decision that followed had been the decision that would, in any case, have followed. She went through them again, slowly, and in twenty-four of the thirty she could point to the sentence in which the memo had given the decision away.

She understood then what she had done. She had not been deceived. No one had hidden anything from her. There was no villain in the room, no suppressed report, no document found in a drawer. The evidence was her own, in her own words, in her own careful hand, and what it showed was something far more intimate than betrayal. She had learned, in the first month, what kind of sentence a company could hear. She had learned that a warning delivered with a clause was received as a contribution, whereas a warning delivered without one was received as a demand, and that a demand might be refused. She had learned to write the first kind. She had done it so well, and so consistently, and with such an unaffected belief in her own conscientiousness, that she had never once noticed that she was producing a warning designed to be declined.

And it was worse than that, because of the thirty-first. The election pause had worked, and for a while she had taken it as proof that she could be effective when she chose. But on the long table, in the grey Saturday light, she saw what the pause really showed, and it was a subtler thing. It had worked because Warren was in the room that morning, and because the question was narrow, and because the cost was small. It had not shown that she possessed a reliable way of compelling action. It had shown that on one occasion she had been lucky. What she had in the other thirty was not a technique for being ignored. It was a technique for being \emph{disposable}, and she had been brilliantly, reliably good at it.

She sat down on one of the chairs. She was aware of the hum of the building around her, the ventilation, the soft electrical tick of a screen in standby, and aware that she was holding, in her left hand, the nearly full mug.

It was Diane who found her. She had come in to collect a coat. She stood in the doorway for a few seconds, taking it in: the table, the stacks, Nora in the chair with her hands in her lap. She came in without speaking and read the end of the nearest memo, and then of the next, and turned her head.

``You've counted,'' she said.

``Thirty,'' said Nora. ``And one.''

Diane sat down across the table. She did not say that it was a coincidence. She did not say that it was a pattern of drafting that any lawyer would have used. She looked at the pages for some time, and then she said, with no preamble, ``Do you know how it began? The company.''

``The story or the truth?''

``There's a difference, but not as much as people think.'' She leaned back. ``It began in a university. A hundred and twenty students in a dormitory, an online directory with photographs, that Warren, who was a graduate resident tutor there, wrote because he wanted to know who lived on which floor. That is the story, and it is true. What isn't in the story is that when he sent the first invitations, in the second week, he used a mailing list that had a flaw in it. He meant to send them to four hundred people. The list resolved to the whole alumni database. It went to four hundred and eleven thousand. By the end of the first night, nine thousand of them had signed up, most of them out of nothing but surprise, and the servers fell over, and he spent a week sleeping on the floor of the machine room, and the product that he had meant for a floor became a place that people lived in.

``I joined at forty employees. I was the person who knew how to buy chairs. And I have heard the story told, I suppose, three hundred times, as a story of a founder who saw further than anyone else. It's a story about a clerical accident. A mistake in a list. And I have watched what that accident did to the way we think of ourselves, because a company that was founded by a mistake and then succeeded for ten years starts to believe, quite sincerely, that its instincts are better than its procedures. Which is why I built the procedures. It is why, in a manner of speaking, I built you.''

Nora looked up.

``I don't mean that unkindly,'' Diane said. ``Every large organisation needs a person whose job is to turn the founder's instinct into something that can be defended in a room. That's what you are. You are the best I have ever seen at framing options. You take a hard thing and you give it a shape that leaders can hold. I couldn't do it. I've tried. I don't have the gift.''

``I'm not sure that I wanted the gift.''

``No.'' Diane was silent for a moment. ``But you have it, and it's useful, and I would be sorry if you spent the weekend on this table deciding that it isn't.''

It was kindly said, and Nora understood it to be kindly said. That was precisely what she found unbearable. Diane was not telling her to stop. She was not telling her anything but the truth, and the truth had been offered in the form of a compliment, in which the thing that Nora had discovered about herself was named, accurately, as her value. It was gentle and it was complete, and Nora saw, with a flicker of something like vertigo, that there was no sentence she could say in reply that would not make her one of the thirty.

Diane gathered her coat. At the door she turned, and said, as one says something that has been on one's mind, ``I'm not asking you to be less good at it. Only not to be afraid of it,'' and went out.

Nora stayed. The light was going, over the car park. She sat looking along the table at the thirty-one small stacks, in the order in which they had been written, and she did not move them. She left them as they were for the cleaners to find, and went home, and on Monday they had been tidied away, without comment, into a recycling bin, and the review was due on Friday.

\partendnote{End of Part Two}

\notesheading{Notes to Part Two}

{\small

\emph{Further notes, keyed to Part Two. The narrator regrets that some of these are longer than the passages they annotate.}

\textbf{Chapter 11, ``a certificate that nobody was responsible for.''} There are two kinds of certificate in the story of the outage. One is a cryptographic object, a few hundred bytes of text that proves a service is what it says it is. The other is a piece of paper in a file, issued by an engineer to himself, that should have said when the first would die. The second was never found. It is believed to have been in a folder in an office that was cleaned in the spring of the previous year.

\textbf{Chapter 11, ``Iosif Marek.''} He drove four hours and found the bridge gone. He turned the car around, drove home, and told his wife it was a good road. She has said that it was the kindest lie of his life. He is a farmer of leeks. He does not speak of Harbor and he has asked that this note not be sent to him.

\textbf{Chapter 12, ``Ottilie Barr and the coffee machine.''} The coffee machine in question is still in the building, on the ninth floor, though it has been mended four times. Its filter basket has never been replaced because nobody is certain whose job it is. A laminated card taped to its side reads \emph{Please report faults to the Roster}. The Roster lists the machine as owned. The owner has never met it.

\textbf{Interchapter, the Playground Game.} Colm's arm is mentioned only once, but it should be recorded that he lowered it at about a quarter past four and put it in his pocket and walked to the fence. He did not ask to be let in. He did not complain. In later life he became a surveyor, and is known for a habit of standing at the edge of rooms.

\textbf{Chapter 14, ``Quickshare.''} The feature was invented in an afternoon by an engineer named Bram who wanted to forward a recipe to his mother's cooking group. It was never intended for politics. It is a mild and useful function and remains one, in the way that a river is mild and useful until a certain week in a certain March.

\textbf{Chapter 15, ``item 41.''} Item 40 was a request to change the colour of a button. Item 42 was a request to add a hyphen to a word in a help page. They were both completed within a week. The narrator has looked at the whole list. Item 41 is the only one that was neither refused nor done, and the only one that mentioned the word \emph{consequence}.

\textbf{Chapter 16, ``the people one sits with.''} In the Tavish language the phrase is three syllables long and is also used to mean \emph{a table}. In the Orrun language there is a proverb about the river that says, in translation, that one cannot choose the people one sits with but can choose which bank. The narrator was told this by a taxi driver in a country whose name the narrator promised not to disclose.

\textbf{Chapter 18, ``I would like to know more than I do.''} The sentence was Warren's own, delivered without a script. Senator Halloran, who had been prepared to dislike him, wrote it on the back of her hand with a ballpoint pen. She did not wash it off for two days. Her aide, who noticed, assumed she was ill.

\textbf{Chapter 19, ``Priya Raman.''} She was twenty-six. Her mother says she was the kind of child who corrected the radio. She can still be heard to do so, though more gently, in the car, when a news presenter says \emph{a majority of people} and means \emph{some}.

\textbf{Chapter 20, ``Do you know how it began?''} The mailing list that went to four hundred and eleven thousand people had a name. It was called \emph{floor-3-west}. The name still appears in an old configuration file, deep in a repository nobody opens, with a comment that reads \emph{TODO: restrict}. It has been there for fourteen years.

}

\part{Warren Brandt}

\numbered{21}{The Roster}

\textsc{it took warren four} days to understand that he could not simply fire the procedure.

He had arrived at the idea in the shower, the morning after the hearing, with the clarity that he got only in the shower and at two in the morning. The trouble with Harbor, he thought, was that things happened in it which nobody had decided. The expired certificate had expired because no one owned it. The grey column in the metrics room was grey because no one owned the pipeline. The memos that Nora wrote, he did not yet know about; but he knew that nine signatures had been needed to change a line of text in the appeals policy, and that nobody could tell him who had designed the nine. So he would find every such thing and give it an owner. A name, a person, a telephone number. It was the simplest fix he had ever conceived, and he was embarrassed not to have thought of it fourteen years ago.

He asked Diane for a list of every service Harbor ran. It was a Wednesday, and the building was at its midweek hum, the particular hum, a half-tone below the Monday one, of people who have stopped planning and started doing. She sent it in an hour: eleven thousand four hundred and six entries. He asked how many had a single named owner, and a woman in operations, after a day, wrote back that the answer depended on what he meant by owner. There was the person who was paged when the service failed. There was the person whose budget paid for it. There was the person whose name was on the vendor contract, where there was one. There was the team that had written it, which had usually been reorganised. For about four thousand of the services these were four different people. For about nine hundred they were nobody.

That was the first thing he learned: that the thing he had meant to fix was not an absence but a surplus. Every function had been assigned several times over, in several different senses, and each assignment had been made by a reasonable person for a reasonable purpose. The pager rota existed because the engineers wanted to sleep. The budget owner existed because Finance would not release money to a ghost. The contract signatory existed because the lawyers insisted that an agreement have a human being at the end of it. He could remove none of these without breaking the purpose it served, and he found, to his considerable irritation, that he did not want to.

So he did what he had always done when a system proved more complicated than its picture. He built a smaller one beside it. He called it the Roster. It was a single list, with one row for every service and one column for a person who would answer, in plain words and without referring the question elsewhere, for what that service did to anyone. He wrote the definition of that column himself, on a Sunday, in six lines, and read it to Pia over the telephone to see whether she could tell what it meant. She said it sounded like the person who would get shouted at. He said yes, exactly, and she said that nobody would take the job, and he said they would have to be asked.

He went through the building asking. Where the conversation was brief, it went well. A pleasant engineer in infrastructure named Tobias Wren put his own name beside forty-one services before lunch and went back to his desk with the air of a man who has been relieved of a secret. A director of payments, Marguerite Osei, took twenty-three and asked for a deputy. Over six weeks the number of services without a person fell from nine hundred to a hundred and ninety, and in that time a fault in the image-upload path that would have cost three days in the old arrangement was repaired in under four hours, because the man whose name was on it knew that his name was on it and answered the telephone.

It was also the moment he met the wall.

The hundred and ninety services that remained unowned were not an oversight. They were the services that could not be given a person without breaking something outside the building. A third of them ran on infrastructure leased from a partner in a country whose law required that the lessee's accountable officer be a resident citizen; Harbor had none there and could not appoint one without opening a legal entity, which required the approval of the board of that entity, which was the very officer Harbor lacked. Another third sat under contracts whose terms assigned operational responsibility to the counterparty, so that a Harbor employee who put her name beside the service would, in the lawyers' phrase, assume a duty that the contract had deliberately declined to create. The rest were governed by an agreement with a regulator that listed approved roles by title, and a title could be added only by an amendment, which required nine signatures, of which two were held by people in another jurisdiction who answered only through counsel.

Warren read the three explanations one after another on a Friday evening, alone in his office, and what struck him was how good they were. There was not one in which a person had acted unreasonably. Each rested on a contract, a statute or an incentive that he could have defended himself if he had been on the other side of the table. He could, in theory, have overridden any of them; the chief executive can do a great deal. But to override the lease he would have had to breach the contract and invite a suit, and to override the regulator he would have had to withdraw from a market of forty million people, and to override the lawyers he would have had to explain to the board why he was spending an enormous sum to protect an abstraction about who should answer a telephone.

He saw, for the first time and without any satisfaction, what Diane had meant on the night after the hearing when she had said that power in a company this size was mostly the power to ask. He had been proud, in the shower, of his plan to give things owners. It had worked, as far as the building went. The building was simply not the whole of what he had built.

He wrote the number on the whiteboard: a hundred and ninety. Under it, after a minute, he wrote \emph{why}, and under that he wrote nothing, because he was a person who disliked a diagram he could not finish. He left it there. It was still on the board in the spring, when the cleaners had been told not to touch it.

\scenebreak
\noindent In April, Tobias Wren was paged eleven times in a single night.

It was not an extraordinary night. A storage cluster in one of the services he had signed for began to fill, and then an alerting rule he had inherited with it fired, correctly, each time the cluster crossed a threshold, and since his name was now in the column marked \emph{answers}, each alert went to his telephone. He had never been paged before. In the old arrangement, the alerts would have gone to a rota, and the rota to a team, and the team to a shared channel, where they would have been seen at eight by whoever came in first. Now they went to him, on the bedside table, with a sound like a small insect, and each time he reached for it his wife Imke woke a little further.

By four he was at the kitchen table in the dark with a laptop, and the cluster was fine, and he was not. He made tea. Imke came down in a cardigan and sat opposite him, and asked, not crossly, whether this was what he had volunteered for.

``It's forty-one services,'' said Tobias. ``I put my name on them because it was the right thing to do. Someone should answer. If it isn't someone, it's no one.''

``And is it a good feeling? Being the one.''

He thought about it honestly. ``It is. In the afternoon.'' He rubbed his eyes. ``At four in the morning it's a bit like being the only person in the village who owns a telephone.''

He went to the Roster's weekly review that Friday and raised it, mildly, with Marguerite Osei, who ran payments and had taken twenty-three services of her own. She had a deputy now, finally, or half of one: a young woman called Linnea who had been lent for six months from another team, with the understanding that she would go back. Marguerite said that her own count was worse. She was being paged at weekends. She had asked for a second deputy, and Finance had said that the Roster was a voluntary initiative and could not support new headcount, since a headcount request must be attached to a business case, and the business case for accountability was, they said, difficult to quantify.

``They're not wrong,'' said Marguerite. ``That's what makes it so hard to be angry. How do you put a number on a telephone that gets answered?''

``We could count the incidents,'' said Tobias.

``We could count the ones that were fixed fast. We can't count the ones that would have been slow.''

They stood in the corridor with their cups. It occurred to Tobias that he had never before understood the economics of being responsible: that it was a cost, paid daily and in small unrecorded coins, and that the institution had no line in any budget for it. The Roster had been celebrated in a newsletter. Nobody had been given a day off.

He wrote to Warren that evening. It was a short note, polite, which said that he was glad to hold his services, and that he wondered whether the company might consider a modest allowance of time in lieu, for people on the Roster who were paged out of hours. Warren replied within the hour, with real warmth, saying that it was an excellent point and that he would raise it with People Operations. Tobias read it with pleasure. It was months before he learned that it had been entered, as item 2,119, in a backlog.

\interchapter{The Sleepover}

\textsc{at ten, warren was} invited to Anders Pell's birthday sleepover, which was the first he had been invited to, and which his mother treated as a diplomatic event. She packed a sleeping bag, a toothbrush, a second toothbrush in case, and a small speech about not touching things.

There were six boys. After pizza and a film that was too old for them, they lay in a row on the Pells' basement carpet in the dark, and at about eleven, one after another, they stopped talking, because there was a noise.

It was a low, ragged hum with a click in it every forty seconds or so, and it came from behind the wall. Anders said it was nothing. Dov Marsh, who was the oldest, said it was a ghost. A boy named Lucan said it was probably the furnace. By midnight the hum had become the only fact in the room. Nobody slept. Nobody wanted to be the person who said that he was scared, and so they discussed, with great solemnity, the several things it might be.

Warren lay on his back and listened, and at some point after one o'clock he began to count the interval between the clicks. It was not forty seconds. It was thirty-eight, then forty-one, then thirty-eight, then forty-two, and it was louder when the pipes upstairs ran. It was not a ghost, which did not have a duty cycle. At twenty to two he got up, found the wall, followed the noise along it with his hand, and came to a door under the stairs. Behind it, in a cupboard full of coats, was a second refrigerator. It was very old. Its compressor was labouring.

He considered what to do. It seemed to him that the noise was a defect, that the defect had a cause, and that the cause was behind a door he could reach. He pulled the plug.

The silence was tremendous. Five boys sat up in the dark and applauded. Anders said, ``Whoa,'' and Warren, who had not been applauded for anything before, went back to his sleeping bag in a state of such happiness that he fell asleep at once.

Mrs.~Pell found the refrigerator in the morning. It contained, among other things, a birthday cake for the following weekend, which had belonged to a different Pell, and eleven pounds of salmon from her husband's fishing trip. She came upstairs holding the salmon and looked at the boys, and then at Warren, and said in a very gentle voice, ``Sweetheart, you could have just come and asked.''

He said that nobody had been awake, and that he had not wanted to bother her. She said that it was no bother. She said that waking a mother was one of the things mothers were for. He saw that she meant it, and that this was a fact about the world which he had not been told: that there existed, in every house, someone to whom the correct answer to a noise was to say that there was a noise. He filed it. It did not seem to him useful. Mrs.~Pell was kind about the salmon, and nobody ever mentioned the cake.

\numbered{23}{The Mitigations}

\textsc{the request came to} Warren on a Tuesday, from a person he had never met, in the form of a four-line email.

Rukhsana Hadid was the editor of a small investigative outlet in a country in which Harbor had a very large number of users and a very small number of staff. Two of her reporters had been identified, the email said, by a coordinated campaign of mass reporting against their accounts, which was now close to suspending them. The accounts held their sources. If they went dark for the length of the ordinary appeal, which in that country's language queue was nineteen days, the reporters' contacts would be exposed to the people they were writing about. She did not ask for anything unreasonable. She asked whether there was somebody she could telephone.

It was exactly the sort of message that reached Harbor a hundred times a week and was answered, correctly and within a day, by a template. This one reached Warren because Rukhsana Hadid had read about the Roster in a trade newsletter, had looked up the chief executive's address in the way that editors can, and had decided that if the company claimed to have given its problems owners she would find out whether it was true.

He forwarded it to Tamsin Vo, who ran Trust Operations, and Tamsin did something unusual. She did not send it up the nine signatures. She phoned Nora.

Nora was in a meeting about the annual policy review, which she had by then finished and filed, and which sat on her desk like a closed book. She went into the corridor. She heard Tamsin out and asked three questions. How many reports had been filed against the accounts, and in what hours? How many of the reporting accounts had existed for under a month? And were the two reporters in the same language queue, or in different ones?

``The same. About four thousand reports, in nine hours. Most of the reporters are new.''

``Then it is a campaign,'' Nora said, ``and the standard appeal is the wrong instrument. We don't need to find that the reporters are right. We need an interim hold on enforcement, for accounts that show a certain pattern, pending a human review in the first twenty-four hours. It already has a precedent. It's in the election pause.''

``We wrote the election pause for one election.''

``We wrote it as a hold on enforcement for a protected class of content during a defined window of risk. It doesn't say elections. It says, `where there is credible indication of coordinated reporting.'\,'' She paused. ``I wrote that. I put it in because I thought we might need it, and I left the word `election' out of the operative sentence on purpose, and I didn't tell anyone.''

There was a silence on the line, in which Nora understood that she had said something which she did not entirely understand herself.

The mitigation was assembled by four in the afternoon. It was a hold, scoped to accounts with the given pattern, valid for seventy-two hours, requiring a named reviewer in the relevant language to confirm or lift it. It was not a heroic act. It was a form with six fields. It required one signature, Tamsin's, since Nora had written the authority so that only one was needed. The accounts were not suspended. The sources were not exposed. When the review was done, three days later, a reviewer named Arvid Soto read the four thousand reports and found that every one of them had been filed from accounts created within the week, and the hold was converted to a permanent flag on the campaign, and the reporters went on working. Rukhsana Hadid sent Warren a second email, shorter than the first. It said: \emph{Thank you. I did not expect a person.}

He read it several times. It was a very good thing, and he had done almost nothing, and he knew exactly whom to thank.

``I would like to tell you it was luck,'' he said to Nora, when he found her that evening at the end of the sixth floor. ``But it wasn't. You built a door into a wall before anyone knew they would need one.''

``It was already in the draft.''

``You put it there. Why?''

She did not answer at once. She was holding a coffee she had stopped drinking. ``I have a habit,'' she said, ``of leaving myself a way out of the rule. I write the rule, and I write the exception, and I never tell anyone about the exception unless I need it. It's a thing I do when I don't trust the room. I think I got it from my mother, who used to say that one should always leave the house with a way to get back in.''

``That seems to me like good practice.''

``It is. That's what I mean. It is good practice and it is also the same habit that\ldots{}'' She stopped, and Warren, who was not good at hearing what people stopped before saying, asked her what. ``Nothing,'' said Nora. ``I'm glad it worked.''

He watched her go down the corridor. It occurred to him that he had received, in the last ten minutes, the second half of a sentence that he had no way of reading, and that the first half was somewhere he could not get to. He wrote it down on the card in his pocket anyway: \emph{N.: exception, mother, ``way back in'', stopped.}

\numbered{24}{Diane at Dinner}

\textsc{diane kessler had a} rule about dinners with the chief executive, which was that they should be at her house and that she should cook, because a person who has been fed is more likely to say what he thinks.

It was early December. Her husband had gone to bed, and the two of them sat at the kitchen table over what remained of a fish stew, and Warren, who had been quiet for most of the meal, put down his spoon and said that he wanted to understand something, and she said that she had been waiting for him to.

``The Roster worked,'' he said.

``It did.''

``Then why does it feel like it didn't.''

She poured the last of the wine, evenly, into the two glasses. This was a question she had hoped he would ask for about three years, and now that he had she found that she was reluctant to answer it, in the way one is reluctant to open a door onto something one has long known was there.

``Because it fixed the part of the company that you can see from the chair,'' she said. ``There are forty thousand people in this building and in the buildings like it. Every one of them is doing something sensible. You went through the whole structure with a flashlight, and wherever the flashlight went, there was a person and the person was fine. I'd say that is a remarkable achievement. But it was never the structure that was the problem.''

``What was?''

``The things you couldn't point the light at. A contract with a counterparty. A rule in a country where you've no standing. A promise you made in 2019 to a man who has since died. Warren, you wrote a column on a whiteboard, a hundred and ninety, and then you wrote \emph{why} and stopped. I saw it. It's the best thing on any wall in this building. The reason you stopped is that the answer is a list of things that are not yours.''

He turned the glass slowly by its stem. ``I'm the chief executive.''

``You are the chief executive of the part that is yours. The rest is held by other people under agreements that you or your predecessors signed with real consideration and real intent, and those agreements are the reason you can run a company at all. A company this size is a set of promises made to people who are not in the room. You can break one. Nobody can stop you. But every promise you break is a reason the next counterparty will not make you one.''

``So the answer is that I should accept it.''

``The answer is that you should know where it is, and that is all I'm saying.'' She looked at him across the table. He had grown, in these months. She could see the places where he had grown, as one sees on a tree where the new wood goes: he asked better questions, he listened for longer, he no longer opened with the solution. He had learned to say \emph{I don't know} without flinching. She had once thought that this was the thing that would save him. She was beginning to suspect that it was only the thing that made him more honest about the problem. ``I'm not telling you it's good. I'm telling you it's a limit. A lot of the pain, in my experience, comes from trying to treat a limit as if it were a fault.''

``And Nora?''

``What about Nora?''

``I think there is something she is not telling me.''

Diane did not say anything for a moment. She was aware of a decision, and of the fact that she was making it quickly, because she did not want to be seen to hesitate. ``Ask her,'' she said. ``She's better at being asked than she is at telling.''

It was true, and it was also a refusal, and Warren, who noticed most things about a mechanism and almost nothing about a motive, thanked her for the dinner and went home.

\scenebreak
\noindent The Qadri family ate together on Fridays at the Ferry House, and had done so since before there were any ferries to speak of.

It stood at the end of a lane in the old port quarter, a long yellow building with a courtyard and a dry fountain, and it was called the Ferry House because Samir's grandfather Idris had begun there, in the year of the great flood, with two flat boats and a rope. He had carried people across the river for a coin until the coins became a fleet, and the fleet became a concession, and the concession a fortune, which the regime of the day had taken from him in a single morning in the way it took everything. The ferries were nationalised. They rusted at their moorings for thirty years. When the regime fell and the family got back what it could, the boats were beyond rescue, but the Ferry House stood, and so did a good deal of land along the river, and Samir's father, who was a practical man, had put fibre into the ground instead. That was the other Qadri fortune. It was newer, and quieter, and nobody in the family quite believed in it.

Samir arrived at one with a tray of pastries and found the long table already crowded. His uncle Hakim sat at its head, as he had since Samir's father died, an immense soft man in a grey robe who held his glass in both hands. Beside him sat his wife, Zubaida, upright and silent. The cousins lined the sides, with their children moving among the chairs like fish. At the foot sat nobody, because the foot was where Yusuf would sit when he came.

The matter before the table was the land. It had been before the table for a year. A consortium that ran the container terminal at the river mouth wished to buy the Ferry House and the three hectares behind it, to build a depot for the trucks, and had offered a sum that Hakim called, with some ceremony, \emph{not insulting}. The port authority was behind the consortium; the port authority was friendly with the ministries; and Hakim, who was seventy and tired, thought that a family which had been robbed once ought to take its money while it was being offered.

``The terminal will grow with or without us,'' he said, to the table. ``I would rather be paid than flattened.''

Samir said what he had said at every Friday for a year, which was that the land was not Hakim's to sell. It was held in endowment. Idris had placed it in trust in the year of the flood, for the benefit of the quarter, and the deed, which Samir had read in a lawyer's office in the spring with the mounting excitement of a boy reading a map, said that the Ferry House and its grounds could be alienated only with the consent of all adult descendants of the founder. All of them. There were thirty-one.

``A deed from the flood,'' said Hakim, pleasantly. ``Do you suppose a court will honour it?''

``I suppose that if you have thirty-one signatures, it won't be asked.''

``And what do you want with it? A boatyard.'' Hakim shook his head. ``You want to teach boys to build boats in a country that buys its boats.''

It was true that Samir wanted the boatyard. He did not know how to explain it, in a room of people who saw land as money or as history and not as a place where something could be done with the hands. His grandfather's boats had been built by men who learned the work at the age of fourteen from men who had learned it at fourteen. The last of those men was eighty-one and lived two streets away. If Samir could get the shed back, and the slip, and a dozen boys, he could have the trade taught again before it died with the old man. It was the one thing he wanted that was entirely his own. It would earn nothing. It would not be mentioned in any newspaper.

The door opened, and the table went quiet in a different way. Yusuf had come in with two aides, who remained at the wall, and he took the chair at the foot, and kissed his father's hand, and accepted tea. He was thirty-nine and slim and looked at nobody for longer than a second. He said that he had heard the discussion was lively. He said that the port authority was a serious body with a long memory, and that families which met its goodwill with stubbornness had sometimes found their licences, their inspections and their customs clearances slower than they might have been. He did not say that this was a threat. He said it as one remarks on the weather coming from the river.

Nobody spoke for some time. Then Zubaida, who had not spoken at a Friday in Samir's memory, put down her cup.

``My husband learned his trade in that shed,'' she said. She was looking at the fountain. ``Before he was a merchant. When I met him he had black under his nails and smelled of pitch. I would like to see a boy come out of that yard with black under his nails.'' She turned her head. ``I will not sign.''

It was not an argument. Hakim opened his mouth, closed it, and looked at his wife with an expression Samir had never seen on him, the helpless respect of a man who has been married for forty-four years and has just been told something new. Yusuf smiled faintly, and finished his tea, and said that of course the family must do as it saw fit. He left before the pastries were passed. Samir sat at the table with his hands flat on the cloth and did not look at anyone, because he had learned long ago that in this house a victory is accepted by not noticing it.

\numbered{25}{The Partner}

\textsc{the annex was forty-one} pages long, and Warren had not read it because nobody had ever told him that there was a reason to.

It was one of fourteen schedules attached to the licence under which Harbor operated in the country that Nora, a year and a half earlier, had been so proud of. The licence itself had been an excellent piece of work. It gave Harbor the right to serve some thirty million people, on terms that protected the users from arbitrary shutdown and Harbor from arbitrary taxation, and it had been negotiated, in large part, by a partner named Samir Qadri, whose family firm held the local infrastructure rights. The licence was why Harbor was there. The annex was why, as it turned out, Harbor was there on the Directorate's terms.

He found it because of Priya Raman. She had come to see him, uninvited, after her presentation in the metrics room, and had said, without ceremony, that she had finished the reconstruction she had begun for the grey column and had found something that was not a data pipeline. In the delta country, the missing metrics were missing because the annex required that certain categories of user information be held locally, under a custodian chosen by the Directorate. The custodian was a company called Veridian Local, and the sole shareholder was the cousin of the man who headed the Directorate, Dr.~Yusuf Qadri, and the cousin of the man who had negotiated the licence for Harbor. The metrics could not be collected because the data on which they depended was not Harbor's to read.

``It is also,'' Priya said, ``what the annex says Veridian does with the data.''

He read it that night, at the kitchen table in his flat, with a pencil. The flat was on the fifth floor of a building that had been, in a former life, a seed warehouse, and when the radiators came on they ticked like a slow metronome and the room smelled faintly of warm iron and old dust. The clause was on page twenty-six and it was written in the plainest language in the document. The custodian would hold, for ninety days, the content and metadata of accounts ``of interest to the public order'', as designated by the Directorate. The designations were made monthly. The list of designated accounts was an attachment to the annex, and the attachment was not included in Harbor's copy.

He telephoned Nora at eleven.

``Did you know?'' he said. He had not intended it to sound like an accusation, and it did.

``I knew there was an annex,'' said Nora. ``I knew it covered local custody. I didn't know about the list.'' She was quiet. ``The licence required local custody for safety reasons. For crime, for child safety. I thought I'd limited it. I wrote a clause that said designations had to relate to the categories in schedule four, which are the crimes.''

``The designation is made by the Directorate.''

``Yes.''

``Who decides whether it relates to schedule four?''

``The Directorate,'' said Nora, after a very long moment. ``I wrote that they must relate to schedule four. I did not write who decides whether they do.'' He heard her breathe. ``That is exactly the kind of thing I'm supposed to catch.''

It was not an accusation, and he did not make it one. He said that he was sorry to have called so late, and that it was nobody's fault, and she said, in a quiet and level voice that he had not heard from her before, that it was somebody's fault, and that the difficulty was that it was everybody's, and that she would be in at seven.

He sat for a long time with the annex. Warren, who built things, understood what had been built. It was a good licence. It had been, on every measurable dimension, in the interests of the thirty million, and it had given the Directorate a lever that nobody had named, and the lever was now attached to a mechanism at which no one in his company was permitted to look. He could withdraw from the market. He had not forgotten the arithmetic: thirty million people, who would be served by something else, a worse thing, run by people who would not have written a clause about schedule four at all. He could stay, and know. He could stay, and ask.

The next morning, at seven, Nora was in his office with a draft. It was eleven pages and it proposed three options and a recommendation. He read it closely, and it was clear and fair and honest. At the bottom of the last page, he found the sentence he was looking for, and he read it twice. \emph{Ultimately, how to weigh the risk to users against the cost of withdrawal is a matter of judgment, and the final assessment rests with leadership.}

``Nora,'' he said.

``Yes.''

``You've done it again.''

She looked at the page, and then at him. She did not ask what he meant. She picked up a pen from his desk, with a steadiness that must have cost her something, and crossed out the sentence, and wrote underneath it in small upright letters: \emph{I recommend that we withdraw the custodian clause from the next renewal, give the Directorate eighteen months' notice, and tell the users in that country what we know. My name is on this.} She signed it and handed it back.

He read what she had written, and it was the most useful thing anyone had put in front of him in a year, and he understood that it had cost her more than the eleven pages. ``Thank you,'' he said, which seemed inadequate, and she nodded and went out. He held it in his hand.

\scenebreak
\noindent The Vessen town council met on the first Tuesday of each month in the old grain exchange, under a ceiling painted with wheat, and on the Tuesday in question there were more people in the public gallery than the room had held since the cathedral vote.

The business was the pitch. The town had received, from a provincial fund, a grant to turn a sloping field by the river into a playing ground, and had spent nine months deciding who should have it. There were two clubs. Vessen Athletic had been founded in 1923 by the men of the quarry, and its supporters lived on the high side of the valley, in the Tavish streets, and spoke of the club as one speaks of a parish. FC Dolm had been founded in 1951 by the river families, who were Orrun, and its supporters sat in the low town and did not forgive the quarry men for the fact that their forefathers had cut the weir. Selk families, who were few and mostly farmed the far slope, supported whoever was winning, and were not trusted by either.

The grant could build one ground. The mayor, a decent harried man called Radomir Vass, had proposed that the clubs alternate Saturdays. Athletic had accepted this in principle and objected to every Saturday on which Dolm would play. Dolm had accepted it in principle and objected to the first Saturday. By eight o'clock they were on the question of which club would be allowed to hang its banner from the pavilion, and a retired quarryman in the gallery had said something about the weir, and the room had risen.

Mirela Costa sat at the back, on a bench by the radiator, with a bag of marking on her knee. She was not on the council and did not intend to speak. She had been invited by Danil's mother, Ilinka, who sat on it and who had taught the same class as her at the age of twelve, and she had come because it was cold outside and the room was warm, and because she had taught, at one time or another, nineteen of the people now shouting.

She watched, and listened as she did with a classroom, not to what was said but to who it was said at. The quarryman was not angry about the weir. He had said it to the man across the aisle, a haulier called Petrit Oanc, and Petrit Oanc's father had owned the haulage firm that had gone under in the strike of 'eighty-four, when the quarry men had not been paid and the river families had not been paid by the quarry men for the carrying. It was a debt. It had been a debt for forty years, and it was never going to be mentioned, and a man cannot be asked to share a field with a family that owes him.

She said nothing. She waited until the mayor had called for order for the third time, and then she got up and walked to the front, which she had not planned to do, and stood beside Ilinka, and asked, in the register of the valley, whether she might say one thing. The hall was so used to her voice at the front of a room that it quietened out of habit.

She did not make a speech. She said that the Selk had a word for the thing under all this, which was not an argument about a pitch. She said it, and a few people nodded, and the quarryman looked away. Then she said, to the mayor, that the schedule could stay as it was, with the Saturdays alternating, and that in the first year the first match should be played by the juniors of both clubs, on a mixed team, with half of each side's families providing the tea. She said that Petrit Oanc's wife made the best seed cake in the low town and that she hoped she would be asked to make it. She did not look at Petrit Oanc when she said this.

There was a silence. A woman in the gallery laughed. Then it spread, and the mayor, who knew when a thing was over, called the vote, and it passed on a show of hands.

She got her price afterward, in the corridor, by the coat hooks, from Ilinka, who was flushed and smiling and had her hand on her arm.

``What do you want?'' said Ilinka. ``I owe you. Whatever you like.''

Mirela had been waiting a year to be asked. She said that the school on Mill Street had a room that had not been used since the heating failed, and that she wanted it on Wednesday and Thursday evenings, for a class. A class for the women of the river streets, who had never been taught to read in their own language because the school had always taught in Tavish. Twelve of them had said they would come. She had the primers. She needed a key, and someone to mend the boiler.

``The boiler will cost,'' said Ilinka.

``I know.''

``I'll put it to the committee.''

``I'd rather you put it to Petrit Oanc,'' said Mirela. ``He has a brother in heating.''

Ilinka stared at her for a moment and then laughed, aloud, with her head thrown back, in the corridor of the grain exchange, and the two of them went out into the cold together with their collars up. The first snow was beginning. Down in the low town, the lights of the river streets were coming on one by one.

\numbered{26}{The Fact-Checker}

\textsc{ilse varga had waited} eleven months for an answer to item 41, and when it came it was a man in a grey jumper standing in the doorway of her office in a city Harbor did not have an office in.

She ran a small fact-checking organisation, one of two hundred and six that Harbor paid, through a consortium, to examine claims flagged by its users. She had met Warren once, at the listening session, where she had asked a question that was entered in the minutes as ``item 41: whether partner checks could be weighted by consequence''. The minutes had gone to a backlog. The backlog had been triaged by a committee, which had judged the request meritorious and unscheduled it, and in this way, for eleven months, it had been neither refused nor done.

Warren had found it in the Roster. When he asked who answered for the partner programme, the name was Hollis Grady's, and Hollis had said, with real embarrassment, that he answered for it in the way a man answers for weather. The programme's terms had been set by the consortium and approved by the lawyers. Hollis could not change them. He could, however, show Warren the backlog, and Warren had read item 41, then item 40, then the other forty-one hundred, and had taken the night train.

``I don't have a decision,'' he told Ilse, when she had made tea. ``I came because I read your question and I did not understand why it was hard.''

She considered him for some time. She was a small, severe woman of fifty with a bad knee, and she had learned over a long career that senior people who arrived without an appointment generally wanted something. Her office smelled of tea and printer toner and, under both, of the pencils she sharpened by hand and kept in a jar, which she said were the only things in the building that still told the truth about how much of them had been used. ``It isn't hard,'' she said. ``It's expensive. Do you know how Harbor pays us?''

``Per check.''

``Per check, at a flat rate, regardless of the claim. A claim that a celebrity has died, which I can settle in six minutes with a telephone call, pays what a claim about a dam pays. The dam takes me nine days.'' She poured. ``I am not a saint. I have nine people to keep employed. I do the dead celebrities.''

He took out the card he kept in his pocket and wrote on it. ``The money rewards the easy ones.''

``The money rewards the number. I have no complaint about it. It is a very sensible system. Anyone can audit it: you count the checks and divide into the budget. If you paid by consequence, someone would have to judge consequence, and then you'd have to explain the judgments, and the consortium would lose a quarter of its members in a week.'' She drank her tea. ``The people who designed this were not stupid, Mr.~Brandt. They wanted something that couldn't be argued with.''

He asked what it would take to change it. Ilse said that it would take an amendment to the consortium agreement, which required two thirds of the members, who would be asked to vote against the rate that kept them solvent. She said it would take about eighteen months. She said this without bitterness, as one reports a tide table.

What he did next was small, and he knew it was small, and he did it anyway. In the weeks after he came home he could not touch the contract, but he could touch his own budget. He established a fund, outside the consortium, that paid partners a supplement for any check that took more than three days, on the single condition that the claim be one that a Harbor reviewer in the relevant country had marked as likely to affect a decision someone would make. He wrote the criteria on one page. It cost a little under two million dollars a year. In the first quarter, Ilse's organisation completed the dam, and two bridges, and a claim about a vaccine lot that had been circulating for a month.

It was a thing that would be gone the day he left, or the day a new finance chief decided it was a perk, and Ilse knew that. She wrote to him in March, a short note. \emph{Thank you for the supplement. I have hired a tenth person. I want you to understand that I will not rely on it.} He read this twice and understood that she was telling him something about the company that he had not been able to see from inside: that the reasonable response to a benefit one cannot count on is to enjoy it and plan as though it will end.

\scenebreak
\noindent The verification desk of the Halden Bureau had four members and one table, and on Monday mornings the table was where the week's arguments were conducted.

It stood in a room above a printer's shop, in a city of the northern plateau that was nobody's idea of a capital, and it was the oldest desk of its kind in the country. There was Emeric Doss, who was sixty-one and had been checking claims since before the word was fashionable, and who ate an apple at every meeting, cutting it into eighths with a pocketknife. There was Bahar Talwar, who was twenty-seven, quick, and funny, and had joined from a regional radio station. There was a part-time statistician called Gunnvor, who joined by telephone from a kitchen. And there was Ilse, at the head of the table, with her bad knee on a second chair, holding a list.

The list had eleven items. Each Monday it was longer than the desk could handle, and the work of the meeting was to cross out all but four.

``The pension fund,'' said Emeric, before she had finished reading. He cut a slice. ``The Minister's statement on Friday. He said the fund was solvent for thirty years. I have the actuarial report. It says twelve.''

``It's a big one,'' said Ilse.

``It's the biggest one. Four million people. If he is lying, or wrong, then everyone over fifty has been told something false by a man who is about to ask for their vote.''

``He'll say it was a different scenario.''

``Then we'll find out which scenario, and we'll print it.''

Bahar had her feet on the rung of her chair. ``I think we should do the poet,'' she said.

Emeric closed his eyes.

``I'm serious. The quote. It's been shared two hundred thousand times. It's attributed to the old national poet, and it's in the speech of the Interior Minister, and the poet never wrote it. It took me twenty minutes. The line is from a pamphlet by an unknown schoolmaster in the 'thirties. It's \emph{beautiful} to check. It's the sort of thing people read to the end.'' She leaned forward. ``I'm not saying it's more important than the pensions. I'm saying that people who read us for the poet will read us for the pensions. If I put out nothing but actuarial tables, in a year we'll have eleven readers, and they'll all be Emeric.''

``There is a respectable argument in that,'' said Emeric, with dignity, ``and I reject it.''

Ilse listened. She had been listening to this argument, in one form or another, for twenty years, and she had never persuaded herself that either side was wrong. Bahar was right that trust was built in small coins. Emeric was right that the small coins were not what the trust was for. It was the sort of disagreement that could only be settled by a judgment of what the desk was, and the desk had never decided.

She turned the page of her list. Item seven was a line she had underlined twice.

``I want the ward,'' she said.

They looked at her.

``The maternity ward at Kessel. The regional board says it closed because the number of births fell below a threshold. I have the figures from the district register. I don't think the number fell.'' She hesitated. ``I think they counted the births in the next town, where the mothers were sent after it closed, and then described the fall as the reason.''

``That would take weeks,'' said Emeric.

``Probably two months.''

``For a hospital in a coal town of nine thousand people.''

``Yes.''

Gunnvor's voice came from the speaker in the middle of the table, slightly delayed. ``It's an interesting number,'' she said. ``I'd have to see how they define a birth.''

``Everybody wants to do something that nobody else would be interested in,'' said Bahar, but she was smiling. ``It's very Ilse.''

``It's not about interest. I don't know if it's important. I know that no one else is going to look at it.'' Ilse put down the list. ``Someone should have to say how the number was made.''

They voted, as they always did, with a show of hands. Emeric took the pensions. Bahar took the poet. Gunnvor, from her kitchen, took the pensions. Ilse took the ward. The ward lost, two to one to one. She said that it was a reasonable decision, and wrote \emph{pensions} and \emph{poet} at the top of the week's page, and Emeric cut himself the last piece of apple and handed her half of it across the table without a word, which was how they had always said the thing they did not say.

She took the ward home that evening in a folder. She spread the district register on her kitchen table beside the cat, and put on her glasses, and began at the first page.

\interchapter{The Lost Dog}

\textsc{when warren was eleven} the dog went missing for three days, and he organised a search.

It was a cheerful, undistinguished spaniel named Biscuit. On the first evening, after the neighbours had been through the usual calls and the walk to the creek, Warren got out his father's street map of the district and a pencil, and divided it into squares of equal area, numbering them from one to sixty. He assigned each square to a household by knocking on doors. Nobody refused. It was the sort of thing that people, faced with a serious boy with a map, found it simpler to accept.

By the second morning the grid was in force. Each household checked its square twice a day, and left a card in the mailbox of the square beyond it, to say that it had done so. Warren walked the perimeter on his bicycle with a clipboard. He was precise. There was a column for sightings, a column for hours, and a column, at his mother's suggestion, headed \emph{remarks}, in which the neighbours entered things like \emph{saw a cat} and \emph{Mrs.~Aldous's gate open}.

The dog came home on the third afternoon, under its own power, wet to the chest, from a direction no square had covered. It had been, it later emerged, at the house of a widower on the far side of the creek who had given it a sausage, and then another.

Everyone was delighted. Warren was praised for his thoroughness. He accepted this, and said that he supposed he should take down the map. But his father, who had a kind of mild attentiveness, said that he might leave it for a while, and he did.

The grid stayed on the kitchen wall for four years. The households went on filling in the cards, for a time, as a courtesy, and then because someone had started to use the remarks column to announce a garage sale, and then because a woman in square thirty-one had a fall, and her neighbours found her in the morning thanks to a card that had not been collected. The system had never found the dog. It found other things. At fifteen Warren took the map down at last, and found that the clipboard had accumulated, in its remarks, four years of what happened in a street: a list of who had been ill, whose roof had leaked, who had moved. He read through it standing in the kitchen, and thought, with the clearest feeling of pride he would remember, that he had made a thing which was better than the thing he had made it for. Then he put it in a drawer, and did not think of it for twenty-five years.

\scenebreak
\noindent The Joint Committee on Communications sat late on the last Thursday of January, and by eleven o'clock the minister had not eaten since breakfast.

The bill under discussion was the Gateway Renewal Bill, which had occupied her for fourteen months, and which concerned, in the dull language of statute, the future of the cable that brought the country to the world. There was only one such route. It ran under the sea from a landing station at the fishing town of Tessaly to a junction in a neighbouring country, and from there to everywhere else. It had been laid by the old regime, who had treated it as a military asset, and it had since been maintained by a state operator, Nordgate, whose engineers had kept it running with great skill and very little money. In four years the cable would reach the end of its life. Something would have to be done, and the committee had spent the evening disagreeing about what.

There were two proposals. Augustin Tavoli, the Finance Minister, wanted the second cable laid through the neighbour's territory, from the capital's side. It was cheaper by a third and could be finished in nine months. It would also pass through a state with which the country had an unresolved border and a long history of quarrels over water. Ines Lindqvist wanted a second domestic route, to a new landing station on the western coast, and a renovated station at Tessaly. It was dearer, it would take three years, and it would employ eleven hundred people, who would otherwise lose their work.

She did not pretend that she was disinterested. Tessaly was her constituency. She had been born in it, in a house at the edge of the harbour wall, and her father had been a cable splicer, who had worked the station for thirty-one years and had come home each evening with his hands smelling of resin and a small brass tool in his pocket that she still carried in her handbag. She had not told the committee that. They knew. Everyone in the building knew where she came from, and a sizeable number thought that it was why she was fighting.

``You are proposing,'' said Tavoli, who was a precise, courteous man with a talent for making his opponents sound sentimental, ``to spend four hundred million more than is necessary, in order to protect the employment of a district.''

``I am proposing to spend it on a route that does not depend on the goodwill of the neighbouring government in a dispute about the river.''

``The neighbouring government has no interest in cutting a cable.''

``The neighbouring government has changed three times in a decade. I would rather not depend on the temperament of the fourth.''

There was a murmur. It was not unfriendly. At the end of the table sat Dorrit Vance, the leader of the Agrarian Union, whose eleven votes would decide it, and who had so far said nothing. She was a heavy, shrewd woman in a green cardigan, who had come into politics from a cooperative dairy and regarded the capital as a province of the countryside that had taken the wrong turn. She was thinking about the levy. Her party wanted a charge on the new cable's revenues, to be spent on rural broadband, and she had made it known, discreetly, that she would trade it.

At a quarter to midnight, in a corridor with a broken water cooler, Vance said to her in a low voice that the Union could support the western route. It would require the levy. It would require, also, a guarantee that the money went to the villages and not to the ministry's consultants.

``You'll have the levy,'' said Ines.

``In the statute. Not in a promise.''

``In the statute.''

``I'd like to believe you.'' Vance looked at her with frank, tired appraisal. ``You know that if you lose the committee, you won't be in the next government. Tavoli will see to it. And Casimir Hale will say that you protected your constituency and put the country at risk.''

``I know.''

``Is it worth it?''

Ines considered. She had been asked the question by her husband, by her chief of staff, and by her mother, in those words or others, and had given each a different answer. She thought of the landing station at night, of its green lamps, of the men who walked its corridors with their hands in their pockets listening for faults. ``I think the country should have a cable that was laid by people who live beside it,'' she said. ``And I think that three years is a long time to wait for it. Both of these are true. I have decided which matters more.''

They returned to the committee room. The vote was taken at twelve minutes past midnight, by hand, and was carried by fourteen to twelve. Tavoli shook her hand with perfect grace. Dorrit Vance did not look at her at all. In the corridor, afterward, her chief of staff, a young man who had been up for twenty hours, said that she had won, and she said that she had won something, and went to find a taxi.

She sat in the back with her handbag on her knee. The city was empty. At the lights by the river she opened the bag and took out the brass tool, a small hooked thing worn smooth at the handle, and held it in her palm in the passing light. She did not think of anything in particular. She put it back, and the lights changed.

\numbered{28}{The Council}

\textsc{the oversight council met} four times a year in a hotel conference room, and had done so since Harbor's board established it, six years before, on the principle that a company which had done as much as Harbor should submit its hardest decisions to people who were not paid by it.

The members were paid by it. This was well known and openly stated, and the Council's charter explained, with great frankness, that independence was a matter of structure and not of income: their stipends were fixed for the term, their names could not be removed except by two thirds of their own number, and they were forbidden from seeking employment with the company for five years after leaving. Nobody had found a better arrangement. There were fourteen of them. They included a former appeals judge, Anselm Okoye, who chaired; a professor of constitutional law; a retired editor; two engineers; a bishop; and a number of others whom Warren, to his private shame, had difficulty remembering in the room.

He came before them on a cold morning in January with Nora and Priya Raman and a forty-page paper, to say that Harbor intended to give notice to the Directorate and to withdraw the custodian clause from the licence at its next renewal, eighteen months hence, and to tell the country's users why.

The Council was, he saw at once, divided, and the division was not along any line he had expected. The professor of constitutional law thought Harbor had no right to disclose; a licence was a contract, and a party that published its terms unilaterally weakened every contract it would ever sign. The editor thought it had a duty to disclose; a company that knew a thing and said nothing was part of the thing. One of the engineers observed that the Directorate would respond to notice by blocking Harbor in the country within the week, in which case thirty million people would lose a service that, whatever else it was, they used to speak to their families.

Judge Okoye listened to all of it with his hands folded. Then he asked Nora a question.

``Ms.~Adeyemi. I have read your memo. I note that it ends with a recommendation, and that it is signed. I have read forty of your memos, in my time here. I do not recall one that was signed in this way.''

``No,'' said Nora.

``May I ask why?''

She was pale, and she did not look at Warren. ``Because I realised that I had stopped being able to say what I thought. I could say what the options were. I could say which one I'd choose, in a sentence that began with `one might.' I think I'd forgotten the other kind of sentence.''

``I see.'' The judge appeared to reflect. ``I will say that it does not make the question easier. A man who says what he thinks may be wrong, and he may be wrong with great authority. There is something to be said for the older style.'' He looked down the table. ``However. The Council has a decision to make, and I propose that we make it.''

The Council did not endorse the withdrawal. It did not oppose it. It voted, nine to five, for a third course, which the retired editor, with a certain dry pride, had devised in the corridor: that Harbor give notice to the Directorate as proposed, and publish nothing to users until the Directorate had replied, or ninety days had passed, whichever came first. The users would be told; but not yet. The Council's recommendation was, by its charter, advisory.

``It is a compromise,'' said Priya, afterwards, on the pavement, with her scarf over her mouth.

``It's a good compromise,'' said Nora. ``That's what's awful about it.''

Warren said nothing. He was looking at the hotel across the street, whose neon sign had been installed, he was sure, by someone who had been asked only to make it visible from the road. It occurred to him that he could do what the Council advised, or what it had not, because the decision was his; and that he had come to the room to be told what to do, and been told, and that the telling had left him more alone than he had been.

\scenebreak
\noindent In January, because the Council had met in the same hemisphere as his daughter, Warren stayed on for a weekend, and on the Saturday morning he was called away from his coffee to deal with a bicycle.

Hanne's house stood at the end of a quiet street lined with bare lindens. He had not lived in it for five years, and he still entered it as a visitor who had been given an exceptionally detailed set of instructions. Hanne, who had been his wife and now was a friend of a precise and wary sort, was a structural engineer; they had met over a failed footbridge. She opened the door in a jumper and said that Pia had been waiting since seven.

The bicycle lay on its side in the hall. Its chain had come off, hung in a loop around the pedal, and Pia, eight years old and indignant, stood over it with her arms folded.

``It keeps doing it,'' she said.

He crouched. It took less than a minute to see what was wrong, and he saw it with a curious pleasure: the rear wheel had been refitted at a slight angle, so that the chain ran across a sprocket tooth that was worn on one side. He explained this to her, with his finger on the place. He showed her the tooth. He said that it was the wheel, not the chain, and that it could be fixed by loosening two nuts and pulling the wheel straight. She nodded, with great gravity. He began to fetch a spanner from the toolbox under the stairs.

``Papa,'' said Pia.

``Yes.''

``Do you want to know what happened?''

He paused. ``I think I know what happened. The wheel's crooked.''

``But do you want to know what happened?'' She sat down on the floor beside the bike, and began to tell him. There had been a hill. It was the one by the school, with the bump halfway. Her friend Soren had gone first and had not fallen off, and she had gone next, and gone faster, because she'd wanted to beat him, and at the bump the bike had made a noise and the chain had come off and she had come off, a bit. She showed him the scrape on her knee. Soren had laughed. She had walked the bike home, and she had not told Mama about the scrape because it was only a small one, and she had been cross for two days.

Warren listened to all of it, sitting on the floor of the hall, with the spanner in his hand. It occurred to him that he had come to the bicycle from the chain, which was a fact, and that she had come to it from Soren, which was another, and that her version had a sequence he had missed: the hill and the wish to win and the walking home. It was true. It did not change the fix. And yet he felt, listening, that he had been given the rest of the thing, and that the thing without it was a nut and a tooth.

``That sounds like a bad day,'' he said.

``It was.'' She leaned against the wall. ``You always go to the broken bit first.''

``Is that a problem?''

``No.~It's good. It's just that sometimes I don't want it fixed yet.'' She thought. ``I want somebody to know first.''

He put the spanner down. He did not fix the bicycle for another quarter of an hour. Hanne, passing the open door with a basket of washing, stopped for a moment and watched them, her eldest quarrel with him gone quiet in her face. Later, over a lunch of soup and bread, she said, in the idle voice she used for the most serious remarks, that Pia had been waiting four years for him to ask her what had happened first. He said that he had not known there was a first. Hanne smiled slightly and said, ``That's what I told the court.''

He wrote nothing on a card that day. He fixed the wheel, and it held. On the plane home he found himself returning, not to the tooth, but to the knee, and the hill, and the small, solemn sentence she had said, which seemed to him, in the dark above the ocean, to contain a whole theory of how organisations should answer the people who write to them.

\numbered{29}{Eleven Days}

\textsc{the notice was served} on a Monday. The Directorate replied on the Thursday, in a letter of four lines, which said that it noted the communication and would respond in due course. It was a better answer than anyone had expected, and for nine days nothing happened.

On the tenth day, a Wednesday, a man in the delta country was shot by police outside a polling station, and a video of it appeared on Harbor.

In the ordinary course it would have been removed within the hour. The footage was graphic, and the policy on graphic violence was clear, and the reviewers in the relevant language queue had been trained on it. The removal would have been correct, and in the normal sequence of events the Directorate would then, as it had done on three previous occasions, have designated the account that posted the video as ``of interest to the public order'', and Veridian Local would have taken custody of its content and its contacts.

The account that posted it belonged to a nineteen-year-old called Teodor Lazar. It had been followed, in the preceding week, by six hundred and eleven accounts, almost all new, and the Directorate's monthly designation list for February was due to be applied on Friday.

Priya found it. She was now, in an informal way that nobody had authorised, the person who looked at the grey cells: she had spent six weeks tracing the hill-district gaps to their sources, and she had a standing query that she ran every morning over the delta, looking for any pattern of followers that resembled the one in the annex's example list. On Wednesday evening it returned a hit. She read the account, then the followers, then the dates. She called Nora, and Nora called Warren.

``We could hold the custody transfer,'' Nora said. ``The licence requires us to comply with designations. It doesn't say when. If we delay the Friday transfer until the review of the account is complete, which is an ordinary internal process, we've done nothing improper.''

``How long does the review take?''

``As long as we want it to.''

``Is that a good answer?''

``It's an answer. It's not an honest one.'' Nora's voice was very even. ``It's what I would have written, a year ago, as a clause. It's a way of making a decision look like a delay.'' A pause. ``Warren, I think you have to decide whether to do it as a decision.''

He did it as a decision. He signed an order, on his own authority, that all custody transfers to Veridian Local from the delta country were suspended pending review, and he put his name at the bottom of it, and sent a copy to the Directorate with a covering note that said, in four sentences, what he had done and why. He did not write that it was a matter of judgment. He wrote that it was his.

The pause lasted eleven days. In that time the Directorate sent two letters and a delegation. Samir Qadri flew in from the capital with a lawyer, and sat across the table in a pale suit and said with great warmth that they had all been through a great deal together, and that he was sure there had been a misunderstanding. The delegation said that the suspension was a breach of the licence. Harbor's general counsel, a thin precise man named Corwin Haye, said that it was probably a breach of the licence, and that the exposure was real, and that he advised the order be lifted. The share price fell by three percent.

\scenebreak
\noindent Samir Qadri did not sleep on the flight home, and it was not because of the conversation with the lawyers.

He sat in the back of the car from the airport with his tie loosened, watching the lights of the ring road go past, and he held, in his lap, a small foil parcel that his wife had pressed on him at the door three days before. It contained, he now remembered, onions. They were the first of the season. He had meant to give them to Warren Brandt, who had asked about onions once, at the first lunch, so long ago, in a way Samir had found oddly moving. He had forgotten them in the hotel. He considered, for a moment, how absurd it was to have flown across half a continent to say that there had been a misunderstanding, with a parcel of onions in his bag.

He had meant it, though. That was what he could not explain to his cousin. There had been no misunderstanding, as such. There had been a clause, and a list, and a man in a grey jumper who had been, to Samir's surprise, correct about all of it.

He had been proud of the licence. He was a man who built infrastructure, whose family had laid the first fibre in the delta, and he had believed, honestly, that giving the Directorate a custodian was a fair trade for being allowed to operate at all. Countries do this. A state asks for a hand on the valve. He had negotiated it carefully, with attention to schedule four, and he had thought that the language was a safeguard. It was only in the last year, when the monthly lists began to include names that he recognised, a student he had once taught to swim, a cousin of the woman who ran his office, that he had understood how much it had come to mean. He had rung Yusuf. Yusuf had said, in the pleasant, final voice of a younger cousin who had become important, that it was a matter of public order and that Samir should not trouble himself.

The car turned off the ring road. He took out his telephone and wrote a message to the Directorate's counsel, and read it, and deleted it, and wrote another. The second said that the Directorate's application for an injunction against the suspension would be best filed with the court on the following Wednesday, at the next sitting. He knew perfectly well that there was a sitting on Friday. There was always a sitting on Friday. A man who knew the court as he did could have had the injunction in four days. He sent the message, and put the telephone in his pocket.

It would give them a week. It would give Teodor Lazar's followers, whoever they were, a week to close their accounts or to leave the country or to become a little less visible, which was all he had to give; and it would look, to everyone who looked, like the ordinary slowness of an ordinary court. He did not think of it as courage. He thought of it as the kind of thing one does with a door one has helped to build, and does quietly, so as not to be asked afterwards whether one had meant it.

At home he put the onions in the cold cupboard. His wife, still awake, asked how it had gone, and he said that it had gone as such things go, and that Mr.~Brandt was a serious man. She asked if he had given him the onions. He said that he had not. She said it was a pity, and he agreed, and went up to bed, where he lay for a long time listening to the sound of the city, and did not tell anyone, then or ever, about the Wednesday.

On the eleventh day the Directorate obtained an injunction from a court that sat in a building without windows, and the transfers resumed. By then Teodor Lazar's account had been closed by its owner, who had moved, with his family, to a cousin's house in another country. The six hundred and eleven followers had had eleven days to do the same, or to delete what they had written, or to do whatever it is that people do when they are warned. Harbor never learned what they did. Priya's query, run again on the day the pause ended, showed that four hundred and six of the accounts had been deleted. She did not know whether they had been deleted by their owners or by someone else.

``It was eleven days,'' Warren said to Diane, in her office, the day after. He looked, she thought, extremely tired. ``I paid a three percent drop in the stock, and the Directorate is going to punish us at renewal, and Corwin says we've damaged the licence. And I don't know if it mattered.''

``It mattered to some of them.''

``Some of them. Four hundred and six, maybe. Maybe none.''

``Warren.'' She came around the desk. ``You wanted to know if it works. It works exactly this much. It works in the number of days that you are willing to pay for. I have never seen a company do better than that, and I have seen most of them.''

He nodded slowly. It was, he reflected, almost the same thing she had said before about limits, and it was no more comforting for being true.

\scenebreak
\noindent He asked her to stay after the meeting. The room was empty, and the light in the long windows had gone the colour of tea.

``I think I should step back,'' Warren said. ``Not out. But back. I have been trying, for two years, to be someone who can run this thing in good conscience. I think I've improved. I don't think it's been enough. There are things in it that I would shut down if they were mine, and they aren't. I'm starting to think that the honest thing would be to hand it to someone who can hold it better, or to break it up, or to close the services that are worst.''

Diane sat down. She had known he would say it. She had prepared for it, in the way one prepares for a diagnosis, and she found that the preparation had been useless.

``May I make the case?'' she said. ``Against. I'd like you to hear all of it. Then do what you like.''

He nodded.

``Forty thousand people work here. That is not the argument. That is the least of it. The argument is that there are roughly nine hundred million people who use what you built to do things that were previously impossible or ruinous. They speak to their families across oceans. They check whether their water is safe. They find work. They organise their neighbourhoods. A woman in a market in the hill country, who never had a bank account, sells her goods to a city she has never seen. I could give you the numbers, and I would, but you know them. This is not a rounding error on a good intention. It is a change in what people can do.

``Consider what happens if you refuse to enter a country because its government is unpleasant. You do not, by refusing, improve the government. You leave thirty million people to be served by something worse, run by people who have no Council and no Roster and no clause about schedule four. Every time we have withdrawn, the replacement has been less scrupulous than we were. I'm not saying that to flatter us. I'm saying that it is the observed result.

``Consider what happens if you resign. I will tell you precisely what happens, because I have been through the contracts. Your successor is chosen by a board, three of whom have not read the annex, and five of whom wish to see the share price recover. The Roster lasts a quarter. The fund for Ilse Varga is cut in the second budget. And the people who depend on this, the schools that run their classes on it and the clinics that book their patients through it, find out that the man who minded has gone.

``Which is not to say you should stay for the sake of the share price. I'm saying you cannot withdraw. It's not a private act any more. It stopped being private about the time the first hospital rebuilt its appointments around us. You are in the position of a man who has built a bridge and discovered that the town now lives on the other side.''

She had stopped. He was looking at his hands.

``I think you're right,'' he said.

``I know.''

``And I think it doesn't answer it.'' He spoke slowly, as he did when he was finding the thing. ``Everything you've said is true, and it's all about what we're useful for. But I don't think being necessary is the same as being entitled. If a thing is indispensable and nobody can make it stop, that's a condition, not a justification. A tide is indispensable.'' He looked up. ``I didn't ask to be the person who decides what nine hundred million people see. I don't think I'd be the right person if I had. I wouldn't want a stranger to hold it. But I'm not sure I should either.''

``Then who?''

``I don't know.'' He smiled faintly. ``That's the part I keep writing down.''

Neither of them spoke. It struck her that they were both correct, and that the two positions did not touch at any point, like two lines drawn on different pieces of paper. A man could not be asked to run an institution against his conscience. An institution could not be allowed to stop because its maker had tired of it. She had thought, for most of her career, that one of those sentences must give way. She now suspected that neither would, and that the work was to live in the gap.

``I'll stay,'' he said, after some time. ``For now. I'm not saying it's right.''

``No,'' said Diane. ``It's only what there is.''

\numbered{30}{The Fellowship}

\textsc{the harbor fellowship for} Civic Accountability was announced in March, and it was Warren's idea, which was one of the reasons that nobody liked to say what was wrong with it.

Ten fellows, chosen from some nine hundred applicants, would spend a year embedded in Harbor's trust and policy functions, with full access to its internal data, and would produce at the end of that year a public report on any aspect of the company's civic effects that they chose. They would be paid by Harbor, at a rate the board thought handsome. Their findings would not be edited. The charter said so, in bold, on the first page, and Nora had written the clause herself.

Warren was moved by it. He had come to believe, since the hearing, that the problem was a shortage of people who could see inside the machine and say what they saw, and the fellowship was his answer: ten of them, with keys. He attended the first week's orientation, which was held in a bright room on the eleventh floor with fruit, and gave a short welcome. It was a good welcome. He said that he did not know whether they would find anything, and that he hoped they would. He said that he had been wrong about most of what he had believed at the start, and expected to be wrong about more. Several of the fellows laughed, and one wrote it down.

The report arrived the following March. It was a hundred and six pages long, and it was the work of ten extremely able people, and it was carefully, fairly and thoroughly argued, and Warren read it on a Sunday with a pencil.

It identified the hill-district gap, and the custodian clause, and the pay-per-check problem. It said, correctly, that the Roster was a real achievement and that the services which remained unowned were unowned for structural reasons. It was balanced, informed, and humane. And as he read, with growing unease, he began to see that the report's sentences had a shape. The shape was not in the facts. It was in how they ended. \emph{These findings, while concerning, should be weighed against the realities of operating at scale.} \emph{The appropriate remedy is a matter on which reasonable observers will differ.} \emph{Ultimately the determination rests with Harbor's leadership.}

He counted. It took him twenty minutes. Of the report's eleven recommendations, nine ended with a clause that placed the decision in his hands. The fellows had arrived with no idea what the clause was. Nobody had taught it to them. They had spent a year in the building, among reasonable people, producing documents that reasonable people would read, and the form had come to them as the weather comes: because it was the way that sentences, in this place, could be heard.

He turned to the last chapter, which was one of the two without it. It was headed ``What We Could Not See'', and it was signed by one author alone. It was eleven pages on the metrics. Priya Raman had written it. It said, plainly, that the dashboard did not contain the thing it claimed to measure, that four hundred and six metrics had become eleven screens, and that the screens were reassuring because the places where the signal was weakest had been quietly left out. She did not say it was a matter of judgment. She said what was missing, and how much, and from where, and she gave the numbers, and at the end she wrote: \emph{I do not know what the missing data would show. That is the finding.}

He put the pencil down. He sat for some time in the Sunday quiet of his flat, with the light going, and then he did what he did when he did not know what to do with a thing, which was to write it on a card. \emph{Ten fellows. Nine clauses. One not.} He pinned the card to the wall above his desk, next to the one about the exception and the mother, and the one about the hundred and ninety, and stood back to look at them: three small pieces of paper, in his small handwriting, on which he had written down, in the course of two years, what he could see and could not yet understand.

\partendnote{End of Part Three}

\notesheading{Notes to Part Three}

{\small

\emph{These notes are more personal than the earlier ones, and less reliable. In several places the narrator has two accounts of an event and has chosen one. Where the choice was close, the other is given here.}

\textbf{Chapter 21, ``four days.''} Warren says it was four days. His calendar says it was three, and that on the fourth he was in a dentist's chair. He has never been willing to concede the dentist. The narrator has gone with four, because Warren's memory of the whiteboard is exact, and because in most matters that can be checked his memory is better than anyone else's, and in most matters that cannot, it is worse.

\textbf{Chapter 21, ``Tobias Wren.''} Tobias's wife Imke says the first night's pages were not eleven but fourteen, and that he has been telling the story with the smaller number for so long that it has gone into his bones. He maintains eleven. She maintains fourteen and has the telephone bill to prove it, though the telephone bill proves only that the alerts went to a number.

\textbf{Interchapter, the Sleepover.} Anders Pell remembers the refrigerator as being in the garage. Warren remembers it as being in a cupboard under the stairs, behind the coats. Both agree that the salmon were lost. Warren says eleven pounds. Anders says it was four fish, which, he points out, is an entirely different quantity, and at his father's usual standards would have come to nearer twenty. The cake is not disputed. The cake was a bad loss and nobody has ever forgiven Warren, in the sense that nobody has thought to.

\textbf{Chapter 23, ``Rukhsana Hadid.''} She would like it recorded that she did not ask Warren for anything. She asked whether there was somebody she could telephone. Warren's recollection is that she asked for him by name. The email survives. She did not.

\textbf{Chapter 24, ``a fish stew.''} The stew was not Diane's. It was Lorne's, from a recipe he had learned from a boat builder in his youth, and Diane has been taking credit for it for thirty years on the principle that a person who is fed should know whom to thank. Lorne has said that he does not mind. Lorne has said this in the tone of a man who is keeping a ledger.

\textbf{Chapter 25, ``page twenty-six.''} It is the twenty-sixth page in Harbor's copy of the annex. In the Directorate's copy it is page thirty-one, because the Directorate's version includes an introduction, in the local language, of five pages. The two copies were never compared, until a lawyer did so on a quiet afternoon and discovered, to nobody's great surprise, that they agreed.

\textbf{The Samir scene, ``a Wednesday.''} The narrator has no source for this scene. No participant in the story knew of the delay, Samir has told no one, and the court's records show a filing on a Wednesday that no one has thought to question. The scene is offered as the likeliest reading of the evidence by a narrator who wishes, on this occasion, to be understood as guessing.

\textbf{Interchapter, the Lost Dog.} The grid was sixty squares, according to Warren. Mrs.~Aldous of square thirty-one remembers sixty-four, and says that Warren added four at the end for the vacant lots, and then did not mention it. Biscuit's whereabouts are agreed. The widower with the sausages was a Mr.~Prosser, a man of few words, who had written, in a note that reached the family on the fourth day, the single line: \emph{He was no trouble.}

\textbf{Chapter 28, ``nine to five.''} The vote was recorded. The abstentions were not. Judge Okoye has said that two members left the room before the count and returned afterwards, one to take a telephone call and one for reasons the judge declined to explain. The minutes record nine to five, a total of fourteen. A count taken with two absent would have been nine to three, or seven to five. The judge does not recall which.

\textbf{Chapter 28, ``the bicycle.''} Warren says the repair took under a minute. Hanne says that it took twenty, and that for fifteen of them he was sitting on the floor of the hall listening to Pia. She says that is the point. Pia, who was asked, says that she does not remember how long it took, only that it was the first time he had asked.

\textbf{Chapter 29, ``the order.''} Warren recalls signing at eleven in the morning. The system log of the signature gives sixteen forty. He has never been able to account for the five hours and forty minutes, though he has been heard to say that he spent them, in the interval, deciding. The narrator offers no view. The narrator has noticed, however, that Warren's recollection of the hour has moved earlier by a few minutes every time he has told the story.

\textbf{Chapter 30, ``nine clauses.''} Priya, asked by the narrator, counted ten. Warren, asked by the narrator, counted nine. They looked at the same report. The disagreement lies in the fourth recommendation, whose final sentence can be read either as a clause or as a conclusion, and which Warren did not recognise as either until it was pointed out, at which point he said that he thought it was a clause. That is also what Priya said.

}

\part{Priya Raman}

\numbered{31}{The Quiet Office}

\textsc{they gave priya raman} an office on the fourteenth floor, and it was the quietest room she had ever worked in.

The promotion came in April, four weeks after the fellows' report. She was made Director of Measurement Integrity, a title that had not existed the previous month and had been devised, with some ingenuity, by Gareth Linde, who explained that it reflected the importance of her chapter. The office had a window, a plant, a door that closed, and a desk the size of a small boat. It did not have a team. This was not stated, and when she asked about it Gareth said that headcount would follow the roadmap, which was in preparation.

She was, for some weeks, genuinely happy. She bought a plant for the windowsill, a small spiky thing that the shop assistant promised was impossible to kill, and watered it on Mondays with the solemnity of a person keeping a promise. She had been listened to, and she had been rewarded, and the two facts seemed to her to be connected in the way that effects follow causes. Her chapter was accepted. It was accepted in principle, a phrase she would later learn to treat as one treats the word \emph{probably} in a weather report, and it was assigned to a working group, which met on alternate Thursdays and was chaired by Gareth. The first recommendation, to collect the missing human metrics in the hill districts, was minuted as \emph{agreed}. The second, which required that Harbor stop excluding uncollected metrics from regional composites, was minuted as \emph{agreed, subject to dependencies}. The dependency was the custodian. The metrics were missing because the data belonged to Veridian Local, and Veridian was not within the working group's remit.

By June she understood the arrangement of the room. Nobody had said no to her. Everybody who could have helped her had said that it was a very good point, and put it on a list. Her office was quiet not because she was being ignored but because there was nobody who needed to disturb her; the stream of requests, questions and small emergencies which had filled her days as an analyst had stopped, as a river stops when a dam has been built upstream, and what remained was a stillness that felt, in the first hours, like respect.

She was asked to design a proxy. This was the working group's solution to the dependency: since the real data could not be had, she was to build an indicator that would approximate it from the data that could. She spent a fortnight on it, and it was a good proxy. It correlated at 0.7 with the metrics for the markets where both were available. She presented it, and the group thanked her, and the hill-district composite, with the proxy included, rose from 0.82 to 0.84.

That afternoon she sat at the enormous desk and looked at the plant. She had come to Harbor from a public health institute, where she had learned a lesson at twenty-three that she had not yet been able to put into words: that a measurement which cannot fall is not a measurement. A proxy built from the data you have, to stand in for the data you lack, will always agree with the data you have. It had been designed, in the kindest possible way, never to alarm anyone.

Warren came to see her in July. He knocked, which no senior person had done in months, and stood in the doorway holding two paper cups.

``I wanted to ask how it was going,'' he said.

``The composite is up two points.''

``Is that good?''

``It's meaningless.'' She heard herself say it, and then, because he was still waiting, she said the rest. She said that the proxy could not detect the thing it replaced. She said that every correct decision the working group had made had made the picture less accurate. She did not soften it with a clause. He listened with his head slightly to one side, as he did, and when she finished he put one of the cups on the desk and said, ``What would you need?''

``A count,'' said Priya. ``Of what we measure and what we don't. Against what we decided. All of it. For fourteen years.''

``Is that possible?''

``I don't know.''

He nodded, as though this were a satisfactory answer, which for him it was. ``Then do that,'' he said, and went out, leaving the second cup on the windowsill, where it remained for a week.

\numbered{32}{The Module}

\textsc{in august, every employee} at Harbor received an email assigning them a course called \emph{Do a Hotel: Asking the Right Question}, estimated duration twelve minutes, due in thirty days.

Priya did hers on a Tuesday, because it was the day she had told herself she would stop avoiding things. It was a good-looking module. The colours were the soft greens and warm oranges of a brand that had been told by a committee that it should feel like a friend. The graphics were warm. A cartoon man in a grey jumper, drawn without a face, stood in a corridor under a cartoon router, and a friendly voice said, ``At Harbor, we believe that curiosity is a responsibility.'' The man looked up. A thought bubble appeared. It said: \emph{I wonder how this works.}

The module had four stages, which it called Notice, Document, Escalate, and Reflect. In the first, the learner was invited to notice a small anomaly. In the second, they were taught to record it in the Observation Register, a tool built the previous year, in which each observation received a number and a status. In the third, they escalated it through the Roster to the owner of the relevant service. In the fourth, they reflected, with the aid of a short prompt, on what they had learned.

At no point did the cartoon man ask anyone anything. Priya went back through it to be sure. The man in the corridor noticed and wrote; he escalated through the correct channel; he reflected. The original act, which had been to walk down a corridor at two in the morning and find the person who knew, and ask, did not appear in the module at all. It had been replaced, one step at a time and each step reasonably, by a procedure for recording that one might have asked.

There was a quiz at the end, as there always was. The fifth question read:

\emph{An employee notices that a service is behaving unexpectedly. What is the BEST first step?}

\begin{enumerate}
\def\labelenumi{(\alph{enumi})}
\tightlist
\item
  Ask a colleague who may know.
\item
  Log the observation in the Observation Register.
\item
  Attempt to fix the issue directly.
\item
  Wait to see whether it recurs.
\end{enumerate}

She selected (a), because it was what she would have done, and the module told her, in a cheerful green box, that this was incorrect. The best first step was (b). \emph{Remember: documenting first ensures that insights are never lost!} She scored eighty percent. The one she missed was the only one that had been about the thing itself.

At the foot of the final screen there was a footnote, in a small grey typeface. \emph{The phrase `doing a hotel' is believed to originate in an informal remark by Harbor's founder, whose habit of investigating problems directly is a source of inspiration to the company. For the avoidance of doubt, employees should follow the procedure above and should not attempt to investigate systems for which they are not the registered owner.}

She read it twice. It occurred to her, with an odd detachment, that it was a nearly perfect specimen: a doctrine that had survived by having its original content replaced, a piece at a time, with its opposite, so that no stage of the process could be said to have betrayed it. Each person who had edited the module had improved it. She would have liked to find a villain. She looked, for a few minutes, in the credits, where the module's authors were listed, and found eleven names, of whom she knew four, and all of whom had been, when she thought about it, friendly.

\scenebreak
\noindent On Sunday Priya took the train to her sister's, as she did every second week, and the only part of her life that was not measured was waiting for her at the end of it.

Meena was thirty-one and a ward nurse, and her house was always slightly too warm and slightly too full: shoes in the hall, a drying rack in the kitchen, a child's drawing of a horse taped to the oven door. Her husband Callum, who taught mathematics at a secondary school, was cooking, in the manner of a man who had discovered that a recipe was a proof and could be followed. Their daughter Tara, who was six, met Priya at the door with a question she had clearly been saving.

``What do you do at work?''

``I count things,'' said Priya.

``What things?''

She thought about it. ``Things people forgot to count.''

Tara considered this seriously, and said that she had once forgotten to count the stairs and fallen down the last one. Priya said that this was exactly the kind of thing she meant. Tara, satisfied, took her by the hand to see the horse.

At dinner Meena talked about her week, as she always did, in the compressed and unsentimental way of someone who has seen a great deal at four in the morning. A man in bed nine had pulled out his line. A drug cart had been left unlocked for twenty minutes. A junior doctor had prescribed something to a patient whom the chart said was allergic to it, and Meena had caught it, and had filed an incident report.

``What happens to it?'' said Priya.

``It goes in the system.''

``And then?''

``And then it's in the system.'' Meena shrugged, helping Tara with her peas. ``There's a committee. It reviews them quarterly. It's a good system, honestly. I've never known anyone to be blamed for filing one.''

``Has anything ever changed because of one?''

Meena stopped, with the spoon in her hand, and thought about it with the attention she gave to dosages. ``No,'' she said. ``Not that I know. I think the point is that it's been written down. If something happens to the patient later, there's a record that someone knew.'' She looked at her sister. ``That's what it's for. It isn't for fixing. It's for afterwards.''

It was said without bitterness, and Priya felt it go in like a splinter. She said that she'd been working on something that sounded very similar. Callum, ladling rice, asked what, and she found herself explaining, in simple terms, the table with the two columns: how much of what a company measures is its own activity, and how little is what happens to the people at the end. He listened with the focus he gave to a student's wrong answer.

``It's the drunk and the lamppost,'' he said.

``The what?''

``A man's lost his keys, and he's looking for them under the lamppost. Someone asks if he dropped them there, and he says no, he dropped them in the park, but the light is better here.'' Callum smiled. ``It's a joke about economists. But it's a joke because everyone does it. You count where you can see.''

``It's more than that,'' Priya said. ``He isn't just looking in the wrong place. Somebody's told him he's found them.''

Their mother rang during the pudding, as she did every Sunday, on the speaker, from the other side of the country, and asked, as she did every Sunday, whether Priya was eating properly and whether the job was stable. Priya said that it was stable. She said that she'd been promoted, and heard her mother's pleasure travel down the line like warmth. Then her mother asked whether she was happy, and Priya, who had not been asked it in months, said ``I'm busy,'' which was not an answer, and her mother, who had spent forty years listening for the difference, let it go.

On the train home she took out her notebook and wrote, under the heading \emph{Method}, a line she had not planned. \emph{The people who are closest to the outcome file the reports. The people who can change the process read the summaries.} She looked at it for a time. It was the beginning of a hypothesis.

\numbered{33}{The Tally}

\textsc{it took her five} months.

She did it in the evenings, mostly, with the door closed. The fourteenth floor emptied at six, in stages, like a tide going out, and by seven the only sounds were the hum of the screens on standby and the occasional cough of the ventilation, which seemed to her, at that hour, to be the building turning over in its sleep. The method was simple. She would take every one of the four hundred and six metrics on the dashboard and ask two questions of it. What does this measure? And what, in the history of the company, has been decided on the basis of it? Then she would take the decisions, which were recorded in a different system, in memos and in the minutes of committees, and ask the converse: what was measured when this was decided?

The first result came in September, and she checked it three times. Of the four hundred and six metrics, three hundred and thirty-one measured something that Harbor did: how fast it responded, how many items it removed, how many appeals it resolved, how much it spent. The remaining seventy-five measured something that happened to a person afterwards: whether the correction reached the people who had seen the claim, whether the appeal resulted in the account being restored in time to matter, whether the reporter who was protected stayed protected. Of those seventy-five, sixty-one were among the metrics that were not being collected.

She sat for some time with the number. It was not a scandal. Nobody had decided it. The metrics for what Harbor did were easy, because Harbor could see them. The metrics for what happened afterward were hard, because afterward was a place Harbor could not see, and a dashboard is a mirror that has been placed where the light is. She made a table of it. It was a plain table, two columns and a total.

The second result was worse, and it was a result about absence.

She had set out to reconstruct the history of the regional safety notices, those requests from the Directorate that the hotline staff in the delta country had been asked to act on since the licence began. She knew that the quarterly reports gave volumes. She wanted the content: which categories, which accounts, how many had been declined. The logs that would have shown it were the hotline's operational logs. She went to the retention schedule to find out where they were kept. They were kept for ninety days.

She read the policy. It was a good policy. It had been written to protect the people who used the hotline, many of whom were frightened, and it ensured that the record of what they had said was destroyed before it could be demanded by a court or a government. It had been approved, according to the schedule, by the Global Policy function, in the autumn the licence was signed, on the signature of Nora Adeyemi.

Priya held the page for a while. The same ninety days that protected the callers from the Directorate also erased any evidence of what the Directorate had asked. The notices that had been declined, and the notices that had been obeyed, and the number of times that a hotline worker had been unable to say no, were gone, each quarter, as the quarter ended. No one had meant it. A privacy rule and a surveillance clause had been written by the same careful person in the same year, and they had met, years later, in a table of Priya's, to make a hole exactly the shape of the thing she was trying to count.

She made a note of the gap. It was a particular kind of finding: a quantity she could prove was unknowable. For the first time since April she was frightened, not of anyone, but of what the work was going to be. To say that something is missing is easy. To say that the person who is responsible for its absence did it while doing her job well, and to say it in a document that will be read by that person's friends, is a different matter.

\interchapter{The Night Clerk}

\textsc{the summer warren was} nine, his family stayed for a week in a seaside hotel with a lift.

It was a narrow old lift with a folding brass gate, and at night it hummed. The hum came up the shaft and through the wall of the room where Warren lay awake beside his sleeping brother, a low continuous note with a tremor in it, and at two in the morning he gave up trying to ignore it. He put on his sandals and went down the stairs in his pyjamas.

The night clerk sat behind the desk in the lobby, under a lamp, with a crossword. He was a man of about sixty with a drooping moustache, and he looked up with no expression at all when Warren appeared at the counter, which was at the height of Warren's chin.

``Excuse me,'' said Warren. ``Why does the lift hum?''

The clerk considered him. He did not ask whether Warren's parents knew where he was. He did not say that it was late, or that guests should not use the stairs at night, or that he was very glad Warren had asked and would pass the question on.

``Counterweight,'' he said. ``It's old and it needs oiling, and the man comes Thursday.''

Warren thought about this. It was a complete answer. It named the cause and gave the date of the remedy, and there was no part of it he could think to ask about. It had the satisfying finality of a key that fits. He said, ``Thank you,'' and the clerk said, ``Go to bed,'' not unkindly, and returned to his crossword.

He climbed the stairs. The lift was still humming when he got to the room, but it was a different noise now, a noise with a reason, and he fell asleep within minutes. In the morning he forgot about it. He would not think of the clerk again for thirty years, and he never learned the man's name. It seemed to him just what the world was like.

\numbered{35}{Offboarding}

\textsc{nora resigned in november,} and the form for it had twenty-three steps.

She had told Warren first, in his office, on a Thursday morning, in about four sentences. She had told Diane second. Neither of them had tried to talk her out of it, which she was grateful for and which also, a little, hurt. Then the machinery took over, as it does, and a ticket appeared in her queue with her name at the top and the words \emph{Offboarding: Senior Vice President, Global Policy and Trust}. One of the steps was \emph{Certify data handover}. The certifier, by a quirk of the workflow, was the nearest available person with the right access, and that was Priya.

They met in Nora's office on the seventh floor, now half packed. The glass room with the view of the car park had been stripped of its pictures. The pale rectangles where they had hung were still visible on the wall, clean and slightly brighter than their surroundings, like the places on a map where nobody has yet written anything. A cardboard box stood on the sofa. Priya brought her laptop and a printout of the checklist, and they went through it item by item, as one reads out the stations on a long line: the badge, the keys, the access to the archive, the archive itself.

``Thirty-one,'' said Nora, at the archive step. ``You'll want the memos.''

``I've read the memos.''

``Of course.'' Nora smiled slightly. ``You'd have found them. You find things.''

Priya put the laptop down. She had thought, on the way up, about how to say what she had to say, and she had planned a sentence that began with the ninety days, and she found she could not begin with it. ``Can I ask why?'' she said instead. ``I'm not certifying that. It's not on the form.''

Nora looked out of the window for a time. The car park was nearly empty. A gardener was moving along the verge with a strimmer.

``When I was a child,'' she said, ``my mother used to say I was too much, and my father used to say I was too sensitive. They said it affectionately. They said it at dinner, to guests, as a way of explaining me. Too much feeling, too much opinion, too much noise; and then, in the next breath, too tender, too easily hurt.'' She turned back. ``I don't think they knew what they were doing. They only wanted me to fit in the world, which is a thing you want for a child. And I learned, as a very good student does, the exact size that I could be. A bit less than too much. A bit more than too sensitive. A size at which people would listen and not flinch.''

``The clause,'' said Priya.

``I never thought of it like that. But yes.'' Nora sat on the arm of the sofa, next to the box. ``It is the correct size. I was proud of it. It's why I was good at this. You say the true thing, and you say it at precisely the volume the room can bear, and the room is grateful and does what it was going to do. I'd done it for so long that I couldn't hear it. I thought it was judgment.'' She laughed softly. ``When I signed the memo that winter, it was the first time in twenty years I'd said something at its actual size. I could feel how big it was. I think it frightened me more than the Directorate.''

Priya wanted to say that the ninety days were not her fault. She wanted to say it badly. She said instead, ``I've found something in the retention schedule. You approved it.''

``The hotline logs.''

``Yes.''

``I know what you've found,'' Nora said. ``I worked it out in the summer, after the annex. I've been waiting to see who would tell me.'' She looked, for a moment, very tired, and entirely steady. ``I was right to approve it. It protected people. And it did what you think it did. Both of those are true. I've decided I don't have to choose which one is the real me.''

It was, thought Priya afterwards, the most generous and the most frightening thing anyone had said to her at Harbor. She ticked the box marked \emph{Data handover certified}. The system sent a confirmation. It was, she noted, timestamped at 4:12 in the afternoon, and would be retained, under the schedule, for ninety days.

\scenebreak
\noindent On her last evening, Nora had Odile, Kobi and Ruth to dinner in the flat, among the boxes.

There was almost nothing left to cook on. She had packed the good pans and was working with one saucepan and a frying pan with a loose handle, and Odile, taking over, did the rice in the saucepan and a chicken in the frying pan with the grim competence of a woman who ran a ward. Kobi had his tooth fixed. He arrived with a small bag, and in it, wrapped in a tea towel, was a second walnut box, the twin of the first, with a lid that closed so softly that Nora lifted it twice to hear it.

``For the other house,'' he said. ``Wherever.''

She did not yet know where. She had a position in the spring, teaching a seminar at a university, and a loose offer from a foundation, and a very large amount of time. She told them this, plainly, with no clause, and noticed that Ruth was watching her do it.

``You've stopped decorating your sentences,'' Ruth said.

``Have I?''

``You used to put a little ribbon on the end. Not a bad one. A courteous one. I always thought it was lovely.'' She raised her glass. ``I find that I like it better with the ribbon off. It's rather startling.''

Odile said that she had a friend in the paediatric unit who had left a consultancy, and who had said that for a year afterward she kept saying \emph{it might be worth considering} to her own children. Everyone laughed, and Nora laughed hardest, and then said, to her own surprise, that she was frightened. Not of the leaving. Of what she would find out she thought, now that there was no room for it to be calibrated to.

``What do you think?'' said Kobi, slowly.

``About what?''

``About anything. Right now. Say it at the size it is.''

She took a breath. She looked round the table: the cheap plates, the half-packed shelves, three people who had never once asked her to explain herself. ``I think I did some good,'' she said. ``I think I did some harm. And I think I'm going to be more use to the people who come after me if I can say which was which without worrying about how it sounds.''

``That was big,'' said Odile.

``It was fine,'' said Kobi. ``It fit the room.''

They stayed until midnight. When she had locked the door behind them and sat down on a box, with the two walnut boxes on the floor beside her, the telephone rang, and it was her father. He asked if she had eaten. She said that she had. He asked if she was sure that she was doing the right thing, and she said, with a calm she had never before used with him, that she was not sure, and that she was going to do it anyway. There was a silence on the line. Then he said, in a voice she hardly knew, ``Your mother says to tell you that you were never too much.'' He cleared his throat. ``I would like to say that I agree.''

She sat for a long time after he had rung off. She did not interpret it.

\numbered{36}{The Draft}

\textsc{the draft of the} tally went to Legal on the second of February, and came back on the ninth with fifty-one comments.

Priya read them on the train. It was the 7:42, a slow service smelling of wet coats and coffee, and she read with her knees against the seat in front and her pencil behind her ear, as she had at university. They were extremely good. Corwin Haye's team had done a careful, courteous, line-by-line pass, and almost all of the comments were what they appeared to be: corrections of a date, a query about a number, a suggestion that a term be defined. They had caught two real errors of her own. She was grateful, and she said so, and she meant it.

The trouble was a smaller class of comment, about nine in all, which came from the same instinct and could not be refused on any ground she could find. They were all very mild. Where she had written \emph{the dashboard does not measure}, the comment suggested \emph{the dashboard does not currently capture}. Where she had written \emph{the custodian arrangement prevents}, it proposed \emph{the custodian arrangement may limit}. Where she had written, in the last paragraph, \emph{I do not know what the missing data would show. That is the finding}, a marginal note observed that the second sentence might be read as a conclusion in the absence of evidence, and offered \emph{That is a limitation of this analysis, and a matter for further review.}

She stared at that one for some time. It was true that it was a limitation. It was true that further review was a thing that could occur. The edit did not make the sentence false. It made it, she saw, the right size.

There was a more serious comment on page thirty-one. She had described the annex, and the list of designated accounts, and the fact that the custodian was owned by the cousin of the man who headed the Directorate. Legal had struck the paragraph. The note said that the licence's confidentiality provisions barred any disclosure of the annex's terms, and that publication would breach the licence and expose Harbor to termination. It was a correct statement of the contract. Corwin added, below it, in his own hand, \emph{I'm sorry. I would not have advised this, and I wish it were otherwise. Please call me.}

She did call him. She spent an hour on the telephone with a thin, exact man who spoke as though each sentence were being entered in a record, and who agreed with her on everything. He said the paragraph was true. He said it was important. He said that publishing it would end the licence, and the licence served thirty million people, and he did not feel that he was in a position to decide that those people ought to lose it because a draft had been too accurate.

``I'm not asking you to decide,'' Priya said. ``I'm asking you not to pretend it's missing.''

``I wouldn't,'' said Corwin. ``I'd say it was withheld.''

In the end she accepted six of the nine softening edits, because they were true, and refused three, including the last sentence, which stayed as she had written it. She kept a version of the struck paragraph with a short note saying that a passage had been removed at the request of counsel, for reasons of contract, and what the passage had concerned in general terms. It was not what she wanted. It was more than she had expected. When the document was released it contained six clauses she would not have written and one blank space she would not have left, and she understood that this was exactly what the institution had been designed to produce: a thing that was mostly true, in a form it could survive.

\scenebreak
\noindent Her error, when she made it, was a small one, and she made it out of the best of her habits.

The clinic was in market seven, on a road that the map called Hollow Road and the people who used it did not call anything. In the second draft of the tally she had cited it as the clearest case of what the grey cells meant. Its reporting to the national health programme was entirely on paper, delivered monthly by a courier, and so it appeared in no dashboard, in no composite and in no proxy. It was, she wrote, a place where a hundred and forty patients a week received care of which no one had any measure. She had thought the sentence was careful. It was true.

The email came from the clinic's coordinator, a woman named Bettina Nkemelu, in a tone of weary courtesy. She had read the draft, which had been shared with a handful of outside reviewers, and she wished to point out that the national programme paid its district clinics by a measure called data completeness, and that Hollow Road's was already the lowest in the region. A public document describing it as unmeasured would be read in the ministry as confirmation. The funding for the pilot, which paid for two of her four nurses, came up for review in September. She did not think Ms.~Raman meant any harm. She thought Ms.~Raman might like to see the place.

Priya went. It took a train, a bus that left twice a day, and forty minutes on the back of a motorbike. The clinic was a long low building with a tin roof and a queue under a tree. Bettina met her at the gate, a broad, quick woman of fifty with a stethoscope round her neck, and walked her through it without explanation: the waiting benches, the ledger with its columns ruled by hand, the cold box for the vaccines, which a nurse named Tamar opened to show that its thermometer had been checked at six that morning and initialled. The ledgers were meticulous. They were arranged so that a nurse could, in under a minute, find any patient who had come in the last four years. There were no gaps in them. Everything that Harbor's dashboard could not see, this room knew.

``You see,'' said Bettina, not unkindly, ``the numbers are not missing. They're only not where you were looking.''

Priya stood in the doorway holding her notebook. It was, she understood, the same mistake she had spent a year describing. She had taken the absence of a measurement to be the absence of the thing measured, and she had done it while holding a document whose last line said that she did not know what the missing data would show. She had been so occupied by the structure of what Harbor could not see that she had not gone to look at what was there. A hundred and forty patients a week were being well cared for by people with a pen. The grey cell had never been a finding about them. It had been a finding about Harbor.

She rewrote the passage on the train home. The clinic came out. In its place she put a general sentence about paper-based reporting, with no district named, and a note that absence of data should not be read as absence of care. It was a better sentence, and it was less vivid, and she knew that the first version would have been quoted. She did not know whether the second would be read at all. She thought of Corwin Haye's soft edits, and for the first time felt some sympathy for the instinct behind them: the wish to say a true thing in a form that would not hurt a person who had not been consulted.

\numbered{37}{The Ordinary Moment}

\textsc{priya's washing machine broke} on a Sunday in March, and she took her laundry to a laundromat on the corner of a street she hardly knew, at ten o'clock at night, because it was the only one open.

It was a bright, bare room with a row of machines along each wall, a rack of plastic chairs, and a vending machine that sold small boxes of detergent and, for reasons nobody could explain, hair ties. It smelled of warm cotton and hot dust and the sweet chemical flowers of a dozen competing detergents, and under the fluorescent tubes everyone looked slightly ill and slightly kind. A television on a bracket was showing a quiz programme with the sound off. In the far corner, in a chair that was too small for him, with a bag of towels at his feet, sat Warren Brandt.

She almost went out again. She stood at the door for a few seconds with her arms full, and he looked up, and neither of them said anything, and then he lifted a hand slightly, the way one acknowledges a neighbour.

``My machine's being replaced,'' he said. ``They said it would be a week. I've done two weeks.'' He seemed neither embarrassed nor pleased. ``I don't think I've been in one of these since university.''

She loaded hers into number six and put in the coins. He watched her do it with interest, as he watched most procedures, and she realised that he was looking at the mechanism of the coin slot. It was a small crowded room, and for the first time in a year there was no agenda, no one to present to, nobody whose approval the next sentence depended on. She sat down two chairs away.

``Pia called,'' he said, after a while. ``She told me about a bird she'd seen. I had to ask her to describe it three times, because the connection was bad.''

``Is it still failing?''

``At two. It's only ever at two.'' He folded a towel on his knee, with care, in thirds. ``I think I'm going to go back. To the hotel. I never asked the question. I wrote it down, and then there was the minister, and then everything, and I never went back to the desk.''

``Why not?''

``I think I was afraid of what it had turned into.'' He smiled slightly. ``I thought, if I ask it now, it will be a Question. It will have a meeting.''

Priya considered this. ``You could just ask the man at the desk.''

``That's what I keep telling myself.''

They sat. The machines turned. In the quiz programme a woman won something and wept silently. Priya was aware of an urge to ask him whether any of it would change, the clauses and the dashboard and the quiet office, and of a plain, accurate understanding that he did not know. He folded another towel. She did not ask. After forty minutes her cycle ended and she took her laundry out, warm and damp and slightly grey, and put it in the bag.

``Good night,'' she said.

``Good night, Priya.''

It was perhaps the only conversation they ever had in which neither of them tried to leave anything behind. Nothing was resolved by it, and nothing was meant to be, and when she got home she found that she had left the fabric softener in the machine, and did not go back for it.

\interchapter{The Right Answer, Late}

\textsc{in the spring of} his twelfth year, Warren asked his teacher why the school bell was always four minutes late.

He had been timing it for a month. It was not a complaint. It was a measurement. The clock in the hall said 9:00, the bell rang at 9:04, and the interval did not vary by more than a few seconds, and he wished to know the reason. Mr.~Aldren, who taught geography and had been teaching it for twenty years and had never been asked this, said that he did not know, but that it was a good question, and that he would find out.

He did not forget. He asked the secretary, who said it was the caretaker's department. He asked the caretaker, who said it was the board's. The board's clerk wrote down the question, and it went, in the manner of such things, into a folder. Warren waited. He did not mind waiting. At first he asked every Monday, and Mr.~Aldren said \emph{still looking}, and by June he had stopped asking, because it was clear that it was neither refused nor going to be answered, and that these were two different conditions, and that the second was less bearable.

In September he went to a different school. In the following March, at the age of thirteen, he received an envelope at home, forwarded from the old address, with a typed note on the board's letterhead. It said: \emph{Regarding your enquiry of last spring. The school bell is synchronised to a master clock which runs four minutes behind the hall clock owing to an uncorrected offset introduced when the system was installed in 1974. The offset is documented and was considered immaterial. We thank you for your interest.}

It was correct. It was precisely and entirely correct, and it said what it should have said, in one paragraph, with the date. He read it at the kitchen table and nodded. Then he put it back in its envelope and went out to play, because he was thirteen and the bell was somebody else's, and the question, which had once been the most interesting thing in the world, had become simply a fact he had been told. Nobody in his life had been waiting for the answer, including himself.

\numbered{39}{The Board}

\textsc{in the spring diane} did the thing she was best at, which was to turn an instinct into a procedure, and this time she was careful to do it well.

She had been thinking about the election pause for a year. Everyone had: it was the one memo that had worked, and the working had come, she now believed, from its form, not its author. It had given a single person the authority to hold an enforcement action for seventy-two hours, on one signature, with a name attached and a reviewer in the relevant language obliged to look within a day. It was a very small power. It had no appeal and no committee. It simply stopped a clock.

She wrote the Standing Hold in four paragraphs. Any director or above could invoke it on any enforcement action in any market, once per quarter, without asking permission, provided they recorded their name and a reason in one line. The hold lapsed at seventy-two hours unless a named reviewer extended it. Corwin Haye read the draft in an afternoon and said, to her surprise, that it was sound. Warren read it and wrote on the margin, in his draughtsman's hand, \emph{this is the best thing in the building}. The Roster had asked people to be answerable. The Standing Hold gave them a way to act.

It worked. In six months there were twenty-six holds, and the review of them, which she conducted herself on a Sunday, was gratifying: a fraudulent takedown reversed in a day, a journalist's account preserved, an opposition party's page kept up through a weekend when a campaign of reports had tried to silence it. No hold had been abused. Not one had been extended for a bad reason.

She was halfway down the second page of the review when she noticed the column of titles. Of the twenty-six, twenty-one had been invoked by vice presidents or above. Three by directors. Two by a single reviewer in a language queue who had been told by a friend, over lunch, that the power existed.

Diane sat with it for a long time. She had built something with no gatekeeper and no form, and it had gone exactly to the people who had always been told things first. The directors who knew about it were the directors who sat in the room where it was discussed. The reviewers in the hotline queues, the people who actually saw the harm at three in the morning in a language none of her executives spoke, did not know, and had not been told, and when she looked at the training materials she found that the module had been scheduled for the autumn, estimated duration twelve minutes.

She did not cancel it. She rewrote it, in four paragraphs, with no quiz, and sent it to every reviewer in the company on a Monday, and got back, over the following weeks, nine holds, of which seven were right. She counted that as a victory and was aware that it was not a solution. The thing that she had built was good. It was also true that a power available to everyone is used first by those already holding the telephone, and she had no procedure for that, and was beginning to suspect that nobody did.

\scenebreak
\noindent The Board of Directors of Harbor met on the first Thursday in May, in a room on the top floor, around a table made of a single piece of timber that someone had once told Warren was four hundred years old. It had the dull, deep sheen of wood polished for generations by sleeves, and when a director set down a cup the sound did not ring but sank.

There were eleven of them. Seven had been there from the beginning, or nearly, and four had been added in the last year, as part of what the chair described as a refresh. One of the four was Ines Lindqvist. She had left her ministry in the autumn after an election, and Harbor's nominating committee, to Warren's surprise and the surprise of several of the seven, had put her name forward. She had accepted. She sat at the far end of the table, in a plain dark jacket, and for the first hour she said nothing at all.

Priya was the fourth item on the agenda. She had been asked to present the tally in fifteen minutes, and she took twelve. She did it standing, without notes, with a single slide, the two-column table. She said what it said. She said that Harbor measured, in the main, what it did, and not what happened to people afterwards, and that the part of the picture that was weakest was exactly the part that was most consequential. She said that the logs that would have allowed an answer were not retained. She said that this was nobody's decision. Then she sat down, and the room was silent, and a director named Leopold Frisch, who had been on the board since the second year, said that it was an extraordinarily thorough piece of work and he hoped the company would reflect on it.

It was Lindqvist who spoke next. She spoke slowly, in good English, with the care of a person who is choosing every word because she intends to be held to it.

``Mr.~Brandt,'' she said. ``You may not remember. I put a question to your colleagues, in a conference room, on the ninth floor of a hotel, nearly four years ago. I asked whose name was on the decision.''

``I remember,'' said Warren.

``I was told a function.'' She smiled very slightly. ``It was an excellent answer. I have been, for the last year, in a position to see what a function does. I would like to ask it again, to you, in this room. Whose name is on the decision?''

There was a pause in which nobody moved. Corwin Haye, who sat at the end of the table as secretary, put down his pen.

Warren did not look at Nora's old memo, or at his notes. He had none. He looked at Lindqvist, and he answered the question she had asked. He said that for the decision in the delta country to pause the custody transfers, the name was his, and he had signed it. He said that for the licence, the name was Nora Adeyemi's and Samir Qadri's and, because he had approved it without reading the annex, his own. He said that for the retention schedule the name was Nora's, but that the policy review that had let it stand for four years was his responsibility, and the failure to read it was also his. He said that for the dashboard it was Gareth Linde's, and Priya's recommendation, and the working group's, but that the group had been convened under his instruction and that the instruction had been to surface what mattered. He said that he did not think this was a satisfying list, and that he did not think a list could be. He said that the honest answer was that there were a great many names, and that he had them here, in a document he had written the night before, and that they were all his.

He said it plainly, without hurrying, and when he stopped, the room was quiet in a different way. It was the answer Lindqvist had asked for. It was complete, and it was correct, and it was what a chief executive would say who had really understood. Priya, in the corner, felt the hair rise on her arms.

``Thank you,'' said Lindqvist.

Corwin Haye cleared his throat. He said, in his careful voice, that he thought the board would want to record its appreciation of the Chief Executive's candour. He said that the statement was of course a significant one, with implications for liability and for the company's positions in several jurisdictions, and that it would be appropriate for it to be reviewed by counsel before it was minuted. He proposed that the minutes record that the Chief Executive had made a statement regarding accountability, and that the full text be held pending that review.

``How long?'' said Lindqvist.

``It is difficult to say. Eight to eleven weeks, I should think.''

The chair put it to the vote. It passed ten to one. Lindqvist's was the one. She did not make a speech. She put her pen down, folded her hands on the old timber, and looked at Warren, who looked back at her, and for a moment they both seemed to be waiting for the other to be surprised.

\scenebreak
\noindent On the night of the Board, Priya telephoned her sister from the station platform, though she had no idea what to say.

Meena answered on the fifth ring, in the middle of something, with a child shouting in the background about shoes. She said she was on her way to bed and asked if everything was all right. Priya said that it was, and that it wasn't, and that she had just stood in a room of eleven people and told them what they didn't know, and that the answer they had given her was to take eight to eleven weeks to think about it.

``Did they say it was a good piece of work?''

``They said it was extraordinarily thorough.''

``Then you're fine.'' Meena laughed, kindly. ``That's what they say.''

``I know. That's the thing. I was prepared for them to be angry, or defensive, or to try to explain it away. I wasn't prepared for them to be nice.'' The train arrived on the opposite platform and left again, and she stood in the quiet after it. ``A man told the truth in that room. The chief executive. He said it all, out loud, every name. And it was held for review.''

``By the lawyers.''

``Yes.''

``For eleven weeks.''

``Yes.''

There was a pause on the line, and in it Priya heard her sister's voice change, from the voice that dealt with children to the voice that dealt with the four in the morning. ``Do you remember what I said about the incident reports? That they're for afterwards?''

``I do.''

``I was wrong, a bit,'' said Meena. ``I've been thinking about it since you came. I told you nothing ever changes because of one. But there was a sister on the ward when I started, Dorothea, who kept every report she ever filed in a ring binder. Thirty years. When they changed the drug chart, she went in with the binder, and she put it on the table. It took her four years. I never saw anyone argue with a binder.'' Meena yawned. ``I'm not saying it's the same. I'm saying that someone has to keep the thing, in case there's an afterwards.''

``The logs are deleted at ninety days.''

``Then keep your own.''

Priya stood on the platform with the phone against her ear, thinking of the table with two columns, and of the printed copy of it in her bag, and of the fact that she had never in her life made a second copy of anything. A train came in. She said she would call on Sunday. She got on, and sat, and took out the notebook, and under the heading she had written in the spring she added a single line: \emph{Keep the binder.}

\scenebreak
\noindent The call in May worked perfectly.

It was a Sunday morning in his flat, and a Sunday evening in hers, and for the first time in his memory, there was nothing wrong. The picture was sharp. The sound was clean. Pia sat at the kitchen table in her pyjamas with a glass of milk, and Warren, who had cleared his diary and turned off his telephone and had thought, the night before, about three questions he wished to ask, asked them. He asked about the heron drawing, and about Soren and the hill, and whether the knee had healed. She told him. He listened with his whole face. He remembered that she disliked mushrooms and had a friend called Clementine who was afraid of dogs, and he asked after Clementine, and Pia stared at him through the screen with a look of faint astonishment.

They talked for forty minutes. It was the best conversation he had ever had with her, and when it ended, he sat looking at the dark screen with the pleasure of a man who has finally got a machine to do what it was built to do.

She rang back at once. She had not got down from the chair. Her face was serious, and a little pink.

``Papa. I don't want to do any more calls.''

He felt something drop in his chest, quite quietly, like a stone into a well. ``Of course,'' he said. ``If you'd rather. We can do fewer.''

``No.~I don't mean fewer.'' She frowned, hunting for it. ``They're good. You're good at them now. You didn't used to be. But it's still a screen.'' She looked at the table. ``I want you to be here. For the thing with the bike. When I'm cross. I want you to be here when it happens.''

He said, carefully, that he could come in August, for three weeks. She said that three weeks was visiting. He said that he would come more often. She said, with the unanswerable patience of eleven, that she knew he would try, and that she wanted him to say whether he would come and stay.

There was a movement in the background, and Hanne appeared at the edge of the screen and put a hand on Pia's shoulder. ``Pia,'' she said gently. ``That's enough for tonight. Go and brush your teeth.'' Pia went. Hanne remained, standing, her face a little to the side of the camera, in the way of a person who intends to say something once.

``She's right, you know,'' she said. ``And she's not entitled to ask it.''

``No.''

``I mean that kindly. You have become a better father. I have watched it, I have not enjoyed watching it, because it is a rather large thing to watch someone become what you needed for years. But I'm not going to move. She has a school here and a life, and I have built mine around it, and what you've made of yourself is real and does not change that.'' She paused. ``I'm telling you because you will do the arithmetic. You'll think, I've improved, therefore. It doesn't follow. Whether she should have you is a matter for her. Not for me, and not for you.''

``I know,'' said Warren, and found that he did.

``Good.'' She looked tired, and honest. ``August. Come for the whole of it. Don't fix anything.''

He sat in the flat with the dark screen for a long time. The answer was perfectly clear. It had no mechanism. There was no cause to find, and no unit to replace, and nothing to count. He took out a card and for some time did not write on it. Then he wrote \emph{Be there}, and looked at it, and did not pin it up, because it did not seem to him that it belonged on the wall with the others.

\numbered{40}{The Hotel Question}

\textsc{he went back in} June.

It was not an official trip. There was a renewal meeting in the capital, the licence's third postponement, and he arrived the night before, and he asked the travel office, with an unusual specificity that caused a ripple of speculation, whether he could have the same hotel. He asked for the ninth floor. The room, when he reached it, was not the same room, but the bed had been made with the same excessive tightness, and the stationery was on the desk.

He lay on the bed for an hour in his clothes. The curtains were the same heavy gold, and the air conditioning, he noticed, made the same low roar with the same ticking inside it, as though it had been waiting. At two minutes to two he rang Pia, who was eleven now and had been told that he would call, and she answered at once, in her kitchen, at one in the afternoon, holding a drawing.

She had drawn a bird. It was a heron, she said, and she held it up to the screen. He was in the middle of asking her what its legs were for when the picture stuttered, and held, and her mouth stopped on a vowel. It was seven minutes past two. The grey wheel turned. The word \emph{reconnecting} appeared.

He sat for forty seconds, as he knew he must. Then he got up and put on his jacket, and checked that the key card was in the pocket, and went out into the corridor, past the interlocking leaves and the closed brass tulips, and stood under the white disc with the amber light. He counted its blinks. The flutter came on the ninth. He took the stairs, because he wanted to, and arrived at the lobby at a quarter past two, where a young man in a waistcoat sat behind the desk with a book.

His name was Cassian, according to the badge. He was perhaps twenty-five, and he looked up with an expression that was perfectly open, neither wary nor eager. Warren said, ``Excuse me. The wireless on the ninth floor drops at seven minutes past two every night. Is it one unit for the floor or one for the whole building?''

``For the building,'' said Cassian. ``It's a scheduled restart. The controller reboots every night from five past two to clear its memory, and it takes the floors down one at a time, so the ninth goes at about seven past. It's all back in ten minutes. It's a fault in the firmware. The supplier says there's an update, but our contract only allows us to change units at the renewal. That's in September.'' He considered. ``I think it would be about eleven weeks.''

It was a sentence and a half, and then a little more. It named the cause. It named the mechanism. It gave a date. Warren listened to all of it, and found that there was nothing in it that he could think to ask, and that it was exactly right. He thanked him. He asked, because it occurred to him, whether a guest could pay for the update himself. Cassian said that the unit was owned by the supplier, and the supplier's terms did not allow it, and Warren said that he understood. It was an entirely reasonable arrangement. It would take eleven weeks.

He was half turned toward the stairs when he saw her. Nora Adeyemi was coming down from the mezzanine, in a long coat, with a conference lanyard round her neck, one hand on the banister, and she stopped on the fourth step from the bottom when she saw him. The hotel hosted a governance conference every June, and she had spoken at it three years running, and she had a panel in the morning. She had not known he was in the city.

She looked from him to the young man at the desk, and she saw at once, as she saw everything, what it must be. She saw the amber light and the stationery and the hour. Her mind began, as it always began, to assemble a meaning: that he had come back to the beginning; that he was settling something; that he was asking, now, in a different room, what he had once been prevented from asking. She felt the first words of the interpretation form, a graceful and perceptive sentence, with a gentle clause at its end.

She did not say it. She stood on the fourth stair with her mouth very slightly open, and let the sentence go, and she looked at the two of them, a man at a counter asking about a router and a young man telling him the truth, and she watched what actually happened. It was nothing. A question was asked. It was answered. It would take eleven weeks.

``Hello, Nora,'' said Warren.

``Hello.'' Her voice was quite steady. ``Is it the router?''

``It restarts at two. They can't change it until September.''

``I see,'' said Nora. And then, because she could think of nothing to add that was not an interpretation, she said it again, more quietly. ``I see.''

They went up together in the lift. Neither spoke. On the ninth floor she got out at the door to the stairs, and said good night, and he went along the corridor to his room, and sat down on the edge of the bed, and picked up his phone. The call was still on the screen. He pressed the button. The grey wheel turned for a moment, and then it said \emph{reconnecting}, and he watched it, and waited, with the patience that was not patience, for a small face to return.

\partendnote{End}

\notesheading{Notes to Part Four}

{\small

\emph{These notes are different in kind. They concern what can and cannot be established. Each records a record, or the absence of one. The narrator regrets the pedantry, which is the point.}

\textbf{Chapter 31, ``four hundred and six.''} The dashboard export lists four hundred and three metrics. Priya's table says four hundred and six. The difference is three metrics that appeared in the interface but not in the export, and which the export tool, according to a footnote in its documentation, ``may omit where status is under review.'' They were under review for the whole of the year. There is no record of what they measured.

\textbf{Chapter 33, ``the retention schedule.''} The schedule survives in three versions. The first, from the period of the licence, states ninety days. The second, a revision, states ninety days ``unless held.'' The third, which no one can date, states ninety days ``unless held or otherwise required by applicable law.'' Priya used the second. The third was found later, in a different system, by a person who had been looking for something else.

\textbf{Chapter 33, ``sixty-one of those seventy-five.''} In a later reconciliation, a data analyst on the measurement team reached sixty-three. She had counted two metrics that Priya classified as process measures. Priya, shown the analyst's list, agreed that the two were ambiguous, and that she had chosen to exclude them because their definitions mentioned, in passing, an outcome. Both counts are defensible. Only one appears in the report.

\textbf{Interchapter, the Night Clerk.} There is no hotel of that description in the town. There is a hotel with a similar name, which was demolished. The lift was removed in the nineteen-seventies. The clerk's name is recorded in none of the guest books that survive. Warren's recollection is nevertheless clear, and the narrator has found no reason to doubt it, except that it is the only thing in the book for which there is no corroboration.

\textbf{Chapter 35, ``twenty-three steps.''} The checklist has twenty-two steps in the version stored in the HR system. A twenty-third, \emph{Certify data handover}, was added by a workflow edit a month before Nora resigned. The edit was made by a person whose name has since been removed from the system in the ordinary course of an unrelated migration. Priya's confirmation, timestamped 4:12 in the afternoon, was retained for ninety days and is therefore not available. Its existence rests on her memory, Nora's, and the entry in a ticket that was closed.

\textbf{Chapter 36, ``fifty-one comments.''} Legal's file contains forty-eight. Three of the comments were made in a chat window and not in the document. Priya remembers all fifty-one, and has said that she can recall the exact wording of the three that were not recorded. The narrator has relied on her. A reader may prefer to rely on the forty-eight.

\textbf{Chapter 37, ``the laundromat.''} Warren and Priya both place the encounter in March. There is a receipt, found in Priya's coat, for a cycle at a laundromat on a street whose name she cannot recall, dated the eleventh of February. The coat was not hers. She had borrowed it. Warren has no receipt, since he does not keep them.

\textbf{Chapter 39, ``whose name is on the decision.''} No transcript exists of the answer. The minutes record that ``the Chief Executive made a statement regarding accountability.'' Three attendees have given accounts. In the first, Warren said ``my name.'' In the second, he said ``mine, and then I have a list.'' In the third, he said nothing at first, and then asked whether he could answer the question in a different order. All three agree that he named the delta licence, the retention schedule, the dashboard, and the pause. They disagree about the order, and about whether he named the working group. The narrator has combined them, and in so doing has produced an account that none of the three would fully recognise.

\textbf{Chapter 39, ``ten to one.''} The vote was ten in favour, one against. The one against is on the record as Lindqvist. There is also a note, in a different hand, that one director who voted in favour did so after a pause of eleven seconds, and that this was timed by a person in the room who was not asked to time it and who has declined to say why.

\textbf{Chapter 40, ``eleven weeks.''} The night manager said eleven. His written note to the front office, filed the following morning, says ten to twelve. The contract renewal date, September, is in the hotel's records. The firmware update was delivered in the last week of August. The restart at five past two was discontinued, according to the hotel, on the second of September, and, according to Warren, on the fourth. The narrator has no means of resolving the discrepancy, and notes only that the unit, when last seen, was in a storeroom in a plastic bag.

\textbf{Chapter 40, the final call.} There is no record of it. Warren's phone logs a connection beginning at 2:31 a.m. that lasted four minutes and thirty-one seconds. The video file was not retained. The narrator, who has described the grey wheel and the word \emph{reconnecting}, did so on the basis of what Warren said and what the logs permit. Whether a small face returned, and what she said, is outside the evidence. The narrator would like to say that it did.
}
